\section{Propriétés constructibles}
\label{section:IV.9-fr}

Soient $S$ un préschéma, $f:X\to S$ un morphisme de
\emph{présentation finie} (1.6.1), $\mathcal F$ un $\mathcal O_X$-Module
quasi-cohérent de \emph{présentation finie}. On se propose dans ce paragraphe
de donner des critères assurant, par exemple, que l'ensemble des $s\in S$ tels
que le $k(s)$-préschéma
$X_s=f^{-1}(s)=X\times_S\operatorname{Spec}(k(s))$ a une certaine propriété,
ou tels que le $\mathcal O_{X_s}$-Module
$\mathcal F_s=\mathcal F\otimes_{\mathcal O_X}k(s)$ a une certaine propriété,
est \emph{localement constructible}, ou tout au moins
\emph{ind-constructible} (1.9.4)~; on verra que c'est le cas pour la plupart
des propriétés qui s'introduisent en Géométrie algébrique. En supposant
seulement $f$ \emph{localement de présentation finie}, on donnera aussi (9.9)
des critères pour que l'ensemble des points $x\in X$ où la fibre $X_{f(x)}$
(ou le $\mathcal O_{X_{f(x)}}$-Module $\mathcal F_{f(x)}$) a une certaine
propriété, soit \emph{localement constructible}. On verra au \S~12 que ces
résultats, joints à l'hypothèse supplémentaire que $f$ est \emph{plat}
(resp. \emph{propre et plat}), permettent de prouver que les ensembles
considérés dans $X$ (resp. dans $S$) sont même \emph{ouverts} dans de nombreux
cas.

\subsection{Le principe de l'extension finie}
\label{subsection:IV.9.1-fr}

\begin{proposition}[9.1.1]
\label{IV.9.1.1-fr}
\textnormal{(Principe de l'extension finie).} Soient $k$ un corps,
$\mathcal C$ un ensemble d'extensions de $k$. On suppose vérifiées les
conditions suivantes~:
\begin{enumerate}[label=(\roman*)]
  \item Si $K\in\mathcal C$ et s'il existe un $k$-homomorphisme $K\to K'$
  (où $K'$ appartient à l'univers où l'on se place), alors
  $K'\in\mathcal C$.
  \item Si $K\in\mathcal C$, il existe une sous-extension $K'$ de $K$, de
  type fini sur $k$, telle que $K'\in\mathcal C$.
  \item Si $K\in\mathcal C$ est le corps de fractions d'une $k$-algèbre
  intègre de type fini $A$, il existe $f\in A-\{0\}$ tel que pour tout idéal
  maximal $\mathfrak m$ de $A$ on ait $A_f/\mathfrak m\in\mathcal C$.
\end{enumerate}
\oldpage[IV-3]{55}
Soit alors $\Omega$ une extension algébriquement close de $k$ (dans l'univers
considéré). Les conditions suivantes sont équivalentes~:
\begin{enumerate}[label=\emph{\alph*)}]
  \item $\mathcal C$ est non vide.
  \item Il existe $K'\in\mathcal C$ qui est une extension finie de $k$.
  \item On a $\Omega\in\mathcal C$.
\end{enumerate}
\end{proposition}

La condition (i) implique évidemment que \emph{b)} entraîne \emph{c)}, et
\emph{c)} entraîne trivialement \emph{a)}~; prouvons que \emph{a)} entraîne
\emph{b)}. En vertu de (ii) et (iii) il existe une extension
$K'\in\mathcal C$ qui est de la forme $A/\mathfrak m$, où $A$ est une
$k$-algèbre de type fini sur $k$ et $\mathfrak m$ un idéal maximal de $A$.
On sait, par le théorème des zéros de Hilbert (Bourbaki, \emph{Alg. comm.},
chap.~V, \S~3, n\textsuperscript{o}~1, cor.~2 du th.~1) que $K'$ est une
extension finie de $k$.

\begin{corollary}[9.1.2]
\label{IV.9.1.2-fr}
Sous les hypothèses (i), (ii), (iii) de (9.1.1), si $k$ est algébriquement
clos et si $\mathcal C$ n'est pas vide, $\mathcal C$ contient toutes les
extensions de $k$ dans l'univers envisagé.
\end{corollary}

En effet, comme \emph{a)} entraîne \emph{c)}, on a $k\in\mathcal C$, et la
conclusion résulte de l'hypothèse (i).

\begin{remark}[9.1.3]
\label{IV.9.1.3-fr}
En pratique, on vérifiera la condition (ii) de (9.1.1) en notant que $K$ est
limite inductive de ses sous-extensions de type fini $K_\alpha$, et en
appliquant les résultats du \S~8, compte tenu éventuellement du fait que pour
$K_\alpha\subset K_\beta$, $K_\beta$ est fidèlement plat sur $K_\alpha$.
Fréquemment l'ensemble $\mathcal C$ est formé des \emph{corps} appartenant à
un ensemble $\mathcal C'$ de $k$-algèbres qui vérifie la condition suivante~:
\begin{enumerate}[label=(i bis)]
  \item Si $A\in\mathcal C'$ et s'il existe un homomorphisme de $k$-algèbres
  $A\to A'$ (où $A'$ appartient à l'univers où l'on se place), alors
  $A'\in\mathcal C'$.
\end{enumerate}
Lorsqu'il en est ainsi, la condition (i) est trivialement satisfaite, et on
vérifiera généralement la condition (iii) de (9.1.1) en notant que le corps
des fractions $K$ de $A$ est la limite inductive des algèbres $A_f$ (pour la
relation $D(f)\supset D(g)$), et en appliquant les résultats du \S~8, compte
tenu éventuellement du fait que les morphismes $D(g)\to D(f)$ sont des
immersions ouvertes.

Par contre, lorsque (i bis) n'est pas vérifiée, la démonstration de (iii) est
souvent plus délicate, et est liée à des critères de constructibilité qui
seront développés plus loin.
\end{remark}

Voici des exemples typiques d'application du principe de l'extension finie~:

\begin{proposition}[9.1.4]
\label{IV.9.1.4-fr}
Soient $k$ un corps, $\Omega$ une extension algébriquement close de $k$,
$X$ et $Y$ deux préschémas de type fini sur $k$. Les conditions suivantes
sont équivalentes~:
\begin{enumerate}[label=\emph{\alph*)}]
  \item Il existe un $\Omega$-morphisme $X_{(\Omega)}\to Y_{(\Omega)}$
  (resp. un $\Omega$-morphisme possédant l'une (déterminée) des propriétés
  (i) à (xiv) de (8.10.5)).
  \item Il existe une extension finie $k'$ de $k$ et un $k'$-morphisme
  $X_{(k')}\to Y_{(k')}$ (resp. un $k'$-morphisme possédant la propriété
  considérée).
  \item Il existe une extension $K$ de $k$ et un $K$-morphisme
  $X_{(K)}\to Y_{(K)}$ (resp. un $K$-morphisme possédant la propriété
  considérée).
\end{enumerate}
\end{proposition}

On applique la remarque (9.1.3) en prenant pour $\mathcal C'$ l'ensemble de
toutes les $k$-algèbres $A$ (de l'univers où on est placé) telles qu'il existe
un $A$-morphisme $X\otimes_k A\to Y\otimes_k A$ (resp. un morphisme ayant
celle des propriétés de (8.10.5) que l'on considère). La condition (i bis) de
(9.1.3) est alors vérifiée grâce au fait que les

\oldpage[IV-3]{56}
propriétés envisagées dans (8.10.5) sont toutes stables par changement de base.
La condition (iii) de (9.1.1) est donc satisfaite en vertu de (8.8.2, (i))
(resp. (8.10.5)), puisque $\operatorname{Spec}(k)$ est quasi-compact et
quasi-séparé. Reste à vérifier la condition (ii) de (9.1.1) qui découle encore
de (8.8.2, (i)) et de (8.10.5), en considérant $K$ comme limite inductive de
ses sous-extensions de type fini. On conclut donc par (9.1.1).

En particulier, s'il existe une extension $K$ de $k$ et un $K$-isomorphisme
$X_{(K)}\xrightarrow{\sim}Y_{(K)}$, on dit que $X$ et $Y$ sont
\emph{géométriquement isomorphes}.

Le corollaire suivant généralise (II, 6.6.5)~:

\begin{corollary}[9.1.5]
\label{IV.9.1.5-fr}
Soient $k$ un corps, $X$ un $k$-préschéma. S'il existe une extension $K$ de
$k$ telle que $X_{(K)}$ soit projectif (resp. quasi-projectif) sur $K$, alors
$X$ est projectif (resp. quasi-projectif) sur $k$.
\end{corollary}

Le morphisme $\operatorname{Spec}(K)\to\operatorname{Spec}(k)$ étant
fidèlement plat et quasi-compact, il résulte déjà de (2.7.1, (v)) que $X$ est
de type fini sur $k$. L'hypothèse signifie qu'il existe une immersion fermée
(resp. une immersion) $X_{(K)}\to\mathbf P_K^r=\mathbf P_k^r\otimes_k K$
(II, 5.5.4, (ii) et 5.5.3)~; appliquant (9.1.4) pour la propriété (iv)
(resp. (ii)) de (8.10.5), on en déduit qu'il y a une extension finie $k'$ de
$k$ et une immersion fermée (resp. une immersion)
$X_{(k')}\to\mathbf P_{k'}^r$, autrement dit $X_{(k')}$ est projectif
(resp. quasi-projectif) sur $k'$. On conclut alors par (II, 6.6.5).

\begin{proposition}[9.1.6]
\label{IV.9.1.6-fr}
Soient $k$ un corps, $\Omega$ une extension algébriquement close de $k$,
$X$ un préschéma de type fini sur $k$, $\mathcal F$, $\mathcal G$ deux
$\mathcal O_X$-Modules cohérents. Supposons qu'il existe un isomorphisme
$\mathcal F\otimes_k\Omega\xrightarrow{\sim}\mathcal G\otimes_k\Omega$.
Alors il existe une extension finie $k'$ de $k$ et un isomorphisme
$\mathcal F\otimes_k k'\xrightarrow{\sim}\mathcal G\otimes_k k'$.
\end{proposition}

Le raisonnement est le même que dans (9.1.4), en appliquant (8.5.2, (i))
(on utilise ici, dans la démonstration de la propriété (iii) de (9.1.1), le
fait que les morphismes $D(g)\to D(f)$ (avec les notations de (9.1.3)) sont
des immersions ouvertes, et \emph{a fortiori} des morphismes plats).

\subsection{Propriétés constructibles et ind-constructibles}
\label{subsection:IV.9.2-fr}

\begin{definition}[9.2.1]
\label{IV.9.2.1-fr}
Soit $\mathbf P(X,\mathcal F,k)$ une relation. On dit que $\mathbf P$ est une
propriété \emph{constructible} (resp. \emph{ind-constructible}) de préschémas
algébriques si les deux conditions suivantes sont vérifiées~:
\begin{enumerate}[label=(\roman*)]
  \item Si $k$ est un corps, $X$ un préschéma algébrique sur $k$,
  $\mathcal F$ un $\mathcal O_X$-Module cohérent, $k'$ une extension de $k$,
  alors, pour que $\mathbf P(X,\mathcal F,k)$ soit vraie, il faut et il suffit
  que $\mathbf P(X_{(k')},\mathcal F\otimes_k k',k')$ soit vraie.
  \item Soient $S$ un préschéma intègre noethérien, de point générique $\eta$,
  $u:X\to S$ un morphisme de type fini, $\mathcal F$ un $\mathcal O_X$-Module
  cohérent. Pour tout $s\in S$, posons
  $X_s=u^{-1}(s)=X\times_S\operatorname{Spec}(k(s))$,
  $\mathcal F_s=\mathcal F\otimes_{\mathcal O_S}k(s)$. Soit $E$ l'ensemble
  des $s\in S$ tels que $\mathbf P(X_s,\mathcal F_s,k(s))$ soit vraie.
  Alors l'un des ensembles $E$, $S-E$ (resp. l'ensemble $E$) contient un
  ouvert non vide (et par suite est un voisinage de $\eta$) (resp. contient
  un ouvert non vide s'il contient $\eta$).
\end{enumerate}
\end{definition}

\oldpage[IV-3]{57}
\begin{remarks}[9.2.2]
\label{IV.9.2.2-fr}
\begin{enumerate}[label=(\roman*)]
  \item Il s'agit là bien entendu d'une convention de langage de nature
  métamathématique et non d'une définition mathématique proprement dite. On a
  des «~définitions~» analogues pour des relations entre $k$, un ou plusieurs
  $k$-préschémas algébriques, des $k$-morphismes entre ces préschémas, des
  Modules cohérents sur ces préschémas ou des homomorphismes entre ces Modules.

  \item Nous aurons aussi à considérer des relations où figurent des
  \emph{parties constructibles} de préschémas. Par exemple, soit
  $\mathbf P(X,Z,k)$ une relation~; nous dirons (par abus de langage) que
  $\mathbf P$ est une propriété constructible (resp. ind-constructible) de la
  partie constructible $Z$ de $X$ si les deux conditions suivantes sont
  vérifiées~:
  \begin{enumerate}[label=$\arabic*^{\mathrm o}$]
    \item Si $k$ est un corps, $X$ un préschéma algébrique sur $k$, $Z$ une
    \emph{partie constructible de $X$}, $k'$ une extension de $k$, alors, pour
    que $\mathbf P(X,Z,k)$ soit vraie, il faut et il suffit que
    $\mathbf P(X_{(k')},p^{-1}(Z),k')$ soit vraie
    ($p:X_{(k')}\to X$ étant la projection canonique).
    \item Soient $S$ un préschéma intègre noethérien, de point générique
    $\eta$, $u:X\to S$ un morphisme de type fini, $Z$ une
    \emph{partie constructible de $X$}. Pour tout $s\in S$, posons
    $X_s=u^{-1}(s)$, $Z_s=Z\cap X_s$. Soit $E$ l'ensemble des $s\in S$ tels
    que $\mathbf P(X_s,Z_s,k(s))$ soit vraie. Alors l'un des ensembles $E$,
    $S-E$ (resp. l'ensemble $E$) contient un ouvert non vide (resp. contient
    un ouvert non vide s'il contient $\eta$).
  \end{enumerate}

  On notera que dans la condition $2^{\mathrm o}$ il faut supposer que $Z$ est
  une partie constructible de $X$, et non seulement que $Z_s$ est une partie
  constructible de $X_s$ pour tout $s$~; la première de ces deux propriétés
  entraîne la seconde (1.8.2), mais non réciproquement.

  \item Si $\mathbf P$ est une propriété constructible, il est clair qu'il en
  est de même de «~non $\mathbf P$~». Si $\mathbf P$, $\mathbf Q$ sont deux
  propriétés constructibles (resp. ind-constructibles), il en est de même des
  propriétés «~$\mathbf P$ ou $\mathbf Q$~» et «~$\mathbf P$ et
  $\mathbf Q$~»~; en effet, si, sous les hypothèses de (9.2.1, (ii)), $E$,
  $E'$ sont deux parties de $S$ et si $E$ contient un ouvert non vide, il en
  est de même de $E\cup E'$, et si $S-E$ et $S-E'$ contiennent chacun un
  ouvert non vide, il en est de même de
  $S-(E\cup E')=(S-E)\cap(S-E')$.

  \item Soit $\mathbf P(X,\mathcal F,k)$ une relation vérifiant la condition
  (9.2.1, (i))~; soient $S$ un préschéma, $u:X\to S$ un morphisme de type
  fini~; avec les notations de (9.2.1, (ii)), soit $E$ l'ensemble des $s\in S$
  tels que $\mathbf P(X_s,\mathcal F_s,k(s))$ soit vraie. Soit d'autre part
  $g:S'\to S$ un morphisme quelconque, et posons
  $X'=X_{(S')}$, $\mathcal F'=\mathcal F\otimes_{\mathcal O_S}\mathcal O_{S'}$~;
  alors il résulte de la transitivité des fibres (I, 3.6.4) et de la condition
  (9.2.1, (i)) que \emph{l'ensemble des $s'\in S'$ tels que
  $\mathbf P(X'_{s'},\mathcal F'_{s'},k(s'))$ soit vraie est égal à
  $g^{-1}(E)$}. Cela s'étend aussitôt au cas où il figure dans $\mathbf P$
  plusieurs préschémas, Modules, morphismes de préschémas ou homomorphismes de
  Modules, ainsi qu'aux propriétés du type considéré dans (ii).

  \item Comme nous le verrons dans la suite de ce paragraphe et dans le courant
  du reste du chap.~IV, la plupart des propriétés $\mathbf P$ vérifiant la
  condition (9.2.1, (i)) vérifient aussi (9.2.1, (ii)). Comme exceptions
  possibles, notons la propriété d'être \emph{projectif}, ou
  \emph{quasi-projectif}, ou \emph{affine}, ou \emph{quasi-affine} sur le
  corps de base (pour un préschéma algébrique)~; nous verrons (9.6.2) que ces
  propriétés sont \emph{ind-constructibles}, mais nous donnerons plus tard un
  exemple où $S$ est une partie ouverte non vide de $\operatorname{Spec}(\mathbf Z)$
  (ou une partie ouverte d'une courbe elliptique sur un corps fini) et où
  toutes les fibres $X_s$

  \oldpage[IV-3]{58}
  sauf la fibre générique $X_\eta$ sont projectives sur $k(s)$ (tous les
  préschémas $X_s$ étant de dimension 2).
\end{enumerate}
\end{remarks}

\begin{proposition}[9.2.3]
\label{IV.9.2.3-fr}
Soient $\mathbf P$ une propriété constructible (resp. ind-constructible) de
préschémas algébriques, $S$ un préschéma, $X$ un préschéma de présentation
finie sur $S$, $\mathcal F$ un $\mathcal O_X$-Module quasi-cohérent de
présentation finie. Alors l'ensemble $E$ des $s\in S$ tels que
$\mathbf P(X_s,\mathcal F_s,k(s))$ soit vraie est localement constructible
(resp. ind-constructible). En outre, si $S$ est irréductible de point générique
$\eta$, un des deux ensembles $E$, $S-E$ est un voisinage de $\eta$ dans $S$
(resp. $E$ est un voisinage de $\eta$ s'il contient ce point).
\end{proposition}

Pour démontrer ces assertions, on peut se borner au cas où
$S=\operatorname{Spec}(A)$ est affine. On sait alors qu'il existe un
sous-anneau $A_0$ de $A$ qui est une $\mathbf Z$-algèbre de type fini, un
$A_0$-préschéma de type fini $X_0$ et un $\mathcal O_{X_0}$-Module cohérent
$\mathcal F_0$ tels que $X$ soit isomorphe à $X_0\otimes_{A_0}A$ et
$\mathcal F$ à $\mathcal F_0\otimes_{A_0}A$ (8.9.1). Soit
$p:S\to S_0=\operatorname{Spec}(A_0)$ le morphisme correspondant à
l'injection $A_0\to A$, et soit $E_0$ l'ensemble des $s_0\in S_0$ tels que
\[
\mathbf P((X_0)_{s_0},(\mathcal F_0)_{s_0},k(s_0))
\]
soit vraie~; alors, en vertu de (9.2.2, (iv)), on a $E=p^{-1}(E_0)$~; on peut
par suite (1.8.2) se borner au cas où $S$ est le spectre d'une
$\mathbf Z$-algèbre de type fini, donc un schéma noethérien. Utilisons le
critère de constructibilité $(\mathbf 0_{\mathrm{III}},9.2.3)$ (resp. le
critère d'ind-constructibilité (1.9.10))~; on est alors ramené, en utilisant
comme ci-dessus (9.2.2, (iv)) et en remplaçant $S$ par un sous-schéma fermé
intègre de $S$, au cas où $S$ est noethérien et intègre, et où il faut prouver
que $E$ est rare dans $S$ ou contient un ouvert non vide de $S$ (resp. que $E$
contient un ouvert non vide de $S$ s'il contient le point générique)~; mais
cela est garanti en vertu de la condition (9.2.1, (ii)).

On notera que l'on n'a utilisé (9.2.1, (ii)) que lorsque $S$ est le spectre
d'un anneau intègre de type fini sur $\mathbf Z$. Il est clair d'autre part
que l'énoncé de (9.2.3) s'applique aussi lorsqu'il figure dans $\mathbf P$
plusieurs préschémas, Modules sur ces préschémas, morphismes de préschémas ou
homomorphismes de Modules. Il s'applique encore lorsqu'il y figure des parties
(en nombre fini) des préschémas considérés, pourvu qu'on impose à ces parties
la condition d'être \emph{localement constructibles}. En effet, la restriction
au cas où $S$ est affine montre qu'on peut se borner au cas où ces parties sont
\emph{constructibles}~: on applique alors (8.3.11) qui montre (avec les
notations précédentes) qu'une partie constructible de $X$ est l'image
réciproque d'une partie constructible de $X_0$ pour un choix convenable de
$A_0$.

\begin{corollary}[9.2.4]
\label{IV.9.2.4-fr}
Soient $\mathbf P$ une propriété constructible (resp. ind-constructible) de
préschémas algébriques, $X$, $Y$ deux $S$-préschémas de présentation finie,
$f:X\to Y$ un $S$-morphisme. Pour tout $s\in S$, posons
$X_s=X\times_S\operatorname{Spec}(k(s))$,
$Y_s=Y\times_S\operatorname{Spec}(k(s))$,
$f_s=f\times1:X_s\to Y_s$. Alors l'ensemble $E$ des $s\in S$ tels que pour
tout $y\in Y_s$, la propriété $\mathbf P(f_s^{-1}(y),k(y))$ soit vraie, est
localement constructible (resp. ind-constructible).
\end{corollary}

En effet, soit $Z$ l'ensemble des $y\in Y$ tels que
$\mathbf P(f^{-1}(y),k(y))$ soit vraie. Comme les fibres $f^{-1}(y)$ et
$f_s^{-1}(y)$ sont isomorphes, on voit que $E$ est l'ensemble des $s\in S$
tels que $Y_s\subset Z$~; si $g:Y\to S$ est le morphisme structural, on a
donc $E=S-g(Y-Z)$.

\oldpage[IV-3]{59}
Or, $f$ est de présentation finie (1.6.2, (v)), donc il résulte de (9.2.3) que
$Z$ est localement constructible (resp. ind-constructible) dans $Y$, donc
$Y-Z$ est localement constructible (resp. pro-constructible) dans $Y$. Comme
$g$ est de présentation finie, $g(Y-Z)$ est localement constructible
(resp. pro-constructible) dans $S$, en vertu du théorème de Chevalley (1.8.4)
(resp. de (1.9.5, (vii)))~; donc $E$ est localement constructible
(resp. ind-constructible) dans $S$.

\begin{remark}[9.2.5]
\label{IV.9.2.5-fr}
On notera que si $\mathbf P$ est une propriété de préschémas algébriques pour
laquelle la prop. (9.2.3) est vraie, alors $\mathbf P$ satisfait aussi à la
condition (9.2.1, (ii))~: cela résulte en effet du fait que dans un espace
irréductible noethérien, un ensemble constructible est rare ou contient un
ouvert non vide $(\mathbf 0_{\mathrm{III}},9.2.3)$.
\end{remark}

\begin{proposition}[9.2.6]
\label{IV.9.2.6-fr}
Désignons par $\mathbf P$ une des propriétés suivantes d'un $k$-préschéma
algébrique $X$~:
\begin{enumerate}[label=(\roman*)]
  \item $X$ est vide.
  \item $X$ est fini sur $k$.
  \item $X$ est radiciel sur $k$.
  \item $\dim(X)$ appartient à une partie donnée $\Phi$ de l'ensemble
  $\overline{\mathbf Z}=\mathbf Z\cup\{-\infty\}$.
\end{enumerate}
Alors $\mathbf P$ est constructible.
\end{proposition}

Il est clair que (i) et (ii) sont des cas particuliers de (iv), en prenant
respectivement pour $\Phi$ l'ensemble $\{-\infty\}$ et l'ensemble
$\{-\infty,0\}$. On n'a donc qu'à démontrer (iii) et (iv). Dans chacun de ces
deux cas la condition (i) de (9.2.1) est remplie en vertu de (2.7.1, (xv)) et
(4.1.4). D'autre part, dans le cas (iii), la propriété $\mathbf P$ vérifie la
conclusion de (9.2.3) en vertu de (1.8.7)~; il reste donc à voir qu'il en est de
même dans le cas (iv). Cela va résulter de la proposition plus précise
suivante~:

\begin{proposition}[9.2.6.1]
\label{IV.9.2.6.1-fr}
Si $f:X\to S$ est un morphisme de présentation finie, la fonction
$s\mapsto\dim(f^{-1}(s))$ est localement constructible.
\end{proposition}

La question est locale sur $S$, donc on peut supposer que
$S=\operatorname{Spec}(A)$ est affine et prouver que pour tout $n$, l'ensemble
$E$ des $s\in S$ tels que $\dim(X_s)=n$ est constructible. Le même
raisonnement que dans (9.2.3) ramène au cas où $A$ est noethérien et intègre,
et il suffit alors de prouver le

\begin{corollary}[9.2.6.2]
\label{IV.9.2.6.2-fr}
Soient $S$ un préschéma intègre noethérien de point générique $\eta$,
$f:X\to S$ un morphisme de type fini. Alors il existe un voisinage de $\eta$
dans $S$ tel que la fonction $s\mapsto\dim(X_s)$ soit constante dans ce
voisinage.
\end{corollary}

Les images par $f$ des composantes irréductibles (en nombre fini) de $X$ qui
ne rencontrent pas $X_\eta$, sont contenues dans des parties fermées de $S$ ne
contenant pas $\eta$ (puisque $S$ est intègre
$(\mathbf 0_{\mathrm I},2.1.5)$), donc (en remplaçant $S$ par un voisinage
ouvert de $\eta$) on peut se borner au cas où toutes les composantes
irréductibles $X_i$ de $X$ rencontrent $X_\eta$~; désignons encore par $X_i$
le sous-préschéma fermé réduit de $X$ ayant $X_i$ pour espace sous-jacent~;
comme $\dim(X_s)=\sup_i(\dim((X_i)_s))$ (4.1.1), on peut se borner au cas où
$X$ est irréductible. Il existe alors un recouvrement fini $(U_j)$ de $X$ par
des

\oldpage[IV-3]{60}
ouverts affines partout denses, et les nombres $\dim((U_j)_\eta)$ sont tous
égaux à $n=\dim(X_\eta)$ (4.1.1.3)~; on peut par suite se borner au cas où $X$
est affine, donc aussi $X_\eta$. Il existe alors, en vertu de (4.1.2), un
ouvert non vide $W$ de $X$ tel que $W\supset X_\eta$, et un
$k(\eta)$-morphisme fini surjectif
$h:W_\eta\to\mathbf V_{k(\eta)}^n
(=\operatorname{Spec}(k(\eta)[T_1,\ldots,T_n]))$~; en appliquant
(8.8.2, (i)) et la méthode de (8.1.2, \emph{a)}), on en déduit (en remplaçant
au besoin $S$ par un voisinage de $\eta$) que l'on a $h=g_\eta$, où
$g:W\to\mathbf V_A^n(=\operatorname{Spec}(A[T_1,\ldots,T_n]))$ est un
$S$-morphisme, et on peut supposer ce morphisme fini et surjectif en vertu de
(8.10.5, (vi) et (x)). On en conclut que pour tout $s\in S$, le morphisme
$g_s:W_s\to\mathbf V_{k(s)}^n$ est fini et surjectif, donc que
$\dim(W_s)=n$ (4.1.2).

\subsection{Propriétés constructibles de morphismes de préschémas algébriques}

\begin{proposition}[9.3.1]
\label{IV.9.3.1-fr}
Soit $\mathbf P(X,k)$ une propriété constructible (resp. ind-constructible)
de préschémas algébriques. Désignons par $\mathbf P'(f,X,Y,k)$ la relation
suivante~: $f:X\to Y$ est un $k$-morphisme de $k$-préschémas algébriques tel
que pour tout $y\in Y$, on ait la propriété
$\mathbf P(f^{-1}(y),k(y))$. Alors $\mathbf P'$ est une propriété
constructible (resp. ind-constructible).
\end{proposition}

En effet, comme $\mathbf P$ vérifie la condition (9.2.1, (i)), il en est de
même de $\mathbf P'$ en vertu de la transitivité des fibres (I, 3.6.4)~; par
ailleurs le fait que $\mathbf P'$ vérifie la condition (9.2.1, (ii)) résulte
de (9.2.4), en raison de la remarque (9.2.5).

\begin{proposition}[9.3.2]
\label{IV.9.3.2-fr}
Désignons par $\mathbf P$ une des propriétés suivantes d'un $k$-morphisme
$f:X\to Y$ de $k$-préschémas algébriques~:
\begin{enumerate}[label=(\roman*)]
  \item $f$ est surjectif.
  \item $f$ est quasi-fini.
  \item $f$ est radiciel.
  \item Pour tout $y\in Y$, $\dim(f^{-1}(y))$ appartient à $\Phi$ (notation
  de (9.2.6)).
\end{enumerate}
Alors $\mathbf P$ est constructible.
\end{proposition}

Cela résulte aussitôt de (9.3.1) et (9.2.6) si l'on tient compte de ce que $f$
est de type fini (1.5.4, (v)), de la caractérisation des morphismes radiciels
(I, 3.5.8), et de celle des morphismes quasi-finis (II, 6.2.2).

\begin{proposition}[9.3.3]
\label{IV.9.3.3-fr}
Supposons vérifiées les hypothèses de (8.8.1) dont nous gardons les
notations~; supposons en outre que $S_\alpha$ soit quasi-compact, $X_\alpha$
et $Y_\alpha$ de présentation finie sur $S_\alpha$, et soit
$f_\alpha:X_\alpha\to Y_\alpha$ un $S_\alpha$-morphisme. Soit $\mathbf P$ une
propriété ind-constructible de morphismes de préschémas algébriques. Pour tout
$s\in S$ (resp. $s_\lambda\in S_\lambda$) posons
$X_s=X\times_S\operatorname{Spec}(k(s))$,
$Y_s=Y\times_S\operatorname{Spec}(k(s))$,
$f_s=f\times1:X_s\to Y_s$ (resp.
$X_{\lambda,s_\lambda}=X_\lambda\times_{S_\lambda}
\operatorname{Spec}(k(s_\lambda))$,
$Y_{\lambda,s_\lambda}=Y_\lambda\times_{S_\lambda}
\operatorname{Spec}(k(s_\lambda))$,
$f_{\lambda,s_\lambda}=f_\lambda\times1:
X_{\lambda,s_\lambda}\to Y_{\lambda,s_\lambda}$). Alors, afin que pour
tout $s\in S$ on ait la propriété $\mathbf P(f_s)$, il faut et il suffit
qu'il existe $\lambda\geqslant\alpha$ tel que pour tout
$s_\lambda\in S_\lambda$, on ait $\mathbf P(f_{\lambda,s_\lambda})$.
\end{proposition}

En effet, soit $E$ (resp. $E_\lambda$) l'ensemble des $s\in S$ (resp.
$s_\lambda\in S_\lambda$) tels que la propriété $\mathbf P(f_s)$ (resp.
$\mathbf P(f_{\lambda,s_\lambda})$) soit vraie~; il résulte de
(9.2.2, (iv)) que l'on a
$E_\mu=u_{\lambda\mu}^{-1}(E_\lambda)$ pour $\lambda\leqslant\mu$, et
$E=u_\lambda^{-1}(E_\lambda)$~; en outre, en vertu de (9.2.3), $E$ (resp.
$E_\lambda$) est ind-constructible dans $S$ (resp. $S_\lambda$)~; la
proposition résulte donc de (8.3.4) appliqué à la partie ind-constructible
$E$ de $S$.

\oldpage[IV-3]{61}
On généralise sans peine ce résultat à des propriétés $\mathbf P$ du type
considéré dans (9.2.3, (i) et (ii)).

\begin{remark}[9.3.4]
\label{IV.9.3.4-fr}
La conjonction de (9.3.3) et de (9.3.2, (ii)) démontre l'assertion
(8.10.5, (xi)).
\end{remark}

\begin{proposition}[9.3.5]
\label{IV.9.3.5-fr}
Soit $\mathbf P(X,Y,k)$ la propriété~: «~$X$ et $Y$ sont deux préschémas de
type fini sur le corps $k$, et il existe une extension $k'$ de $k$ et un
$k'$-morphisme $g:X_{(k')}\to Y_{(k')}$ vérifiant $\mathbf Q(g)$~», où
$\mathbf Q$ est une des propriétés (i) à (xiv) de (8.10.5). Alors $\mathbf P$
est une propriété ind-constructible.
\end{proposition}

La définition de $\mathbf P$ montre en effet que la condition (9.2.1, (i))
est satisfaite, compte tenu de ce que la propriété $\mathbf Q(g)$ est stable
par changement du corps de base, et de ce que deux extensions de $k$ peuvent
toujours être considérées comme des sous-extensions d'une troisième extension
de $k$. Pour vérifier (9.2.1, (ii)), on peut se borner au cas où
$S=\operatorname{Spec}(A)$ est affine~; si $K=k(\eta)$, corps des fractions de
$A$, il existe par hypothèse et en vertu de (9.1.4) une extension finie $K'$ de
$K$ et un $K'$-morphisme
$g':(X_\eta)_{(K')}\to(Y_\eta)_{(K')}$, vérifiant $\mathbf Q(g')$, et $K'$ est
évidemment le corps des fractions d'une $A$-algèbre intègre finie $A'$. Si
l'on pose $S'=\operatorname{Spec}(A')$, il résulte alors de (8.10.5) qu'il y
a un voisinage $U'$ du point générique $\eta'$ de $S'$ tel que, si l'on pose
$X'=X\otimes_A A'$, $Y'=Y\otimes_A A'$, il existe, pour tout $s'\in U'$, un
morphisme $X'_{s'}\to Y'_{s'}$ ayant la propriété $\mathbf Q$. Mais le
morphisme $h:S'\to S$ est fini, donc fermé, et comme
$h^{-1}(\eta)=\{\eta'\}$, $h(U')$ contient un voisinage ouvert $U$ de $\eta$
dans $S$~; pour tout $s\in U$, il y a donc un $s'\in U'$ tel que $h(s')=s$,
et comme
$X'_{s'}=X_s\otimes_{k(s)}k(s')$,
$Y'_{s'}=Y_s\otimes_{k(s)}k(s')$, la propriété
$\mathbf P(X_s,Y_s,k(s))$ est vraie pour tout $s\in U$.

\begin{env}[9.3.6]
\label{IV.9.3.6-fr}
\textbf{Exemple.} Prenons par exemple pour $\mathbf Q$ la propriété d'être un
isomorphisme. Alors, en conjuguant (9.3.5) et (9.3.3), on a la propriété
suivante~: les notations et hypothèses étant celles de (8.8.1), $S_\alpha$
étant quasi-compact, $X_\alpha$ et $Y_\alpha$ de présentation finie sur
$S_\alpha$, afin que, pour tout $s\in S$, $X_s$ et $Y_s$ soient
géométriquement isomorphes (9.1.4), il faut et il suffit qu'il existe
$\lambda\geqslant\alpha$ tel que, pour tout $s_\lambda\in S_\lambda$,
$X_{\lambda,s_\lambda}$ et $Y_{\lambda,s_\lambda}$ soient géométriquement
isomorphes.
\end{env}

On a un résultat analogue lorsque les préschémas que l'on considère sont munis
de «~lois de composition~» d'une certaine espèce
$(\mathbf 0_{\mathrm{III}},8.2.1)$, par exemple des «~préschémas en groupes~»,
«~préschémas en anneaux~», etc. $(\mathbf 0_{\mathrm{III}},8.2.3)$. Alors
l'énoncé précédent est encore valable lorsque par «~isomorphisme~» on entend
des isomorphismes de préschémas qui sont des homomorphismes pour les lois de
composition considérées $(\mathbf 0_{\mathrm{III}},8.2.2)$~; il suffit ici
d'utiliser non seulement (8.10.5) mais aussi (8.8.2, (i)) en remarquant que la
notion d'homomorphisme pour une loi de composition s'exprime en écrivant que
des diagrammes de morphismes de préschémas sont commutatifs (il faut bien
entendu que les morphismes de transition $X_\mu\to X_\lambda$ et
$Y_\mu\to Y_\lambda$ pour $\lambda\leqslant\mu$ soient des homomorphismes
pour les lois de composition envisagées).

On peut aussi, au lieu de considérer des morphismes de préschémas comme dans
(9.3.5), considérer des homomorphismes de Modules, en utilisant (9.1.6) au
lieu de (9.1.5).

\oldpage[IV-3]{62}

\subsection{Constructibilité de certaines propriétés des modules}
\label{subsection:IV.9.4-fr}

\begin{env}[9.4.1]
\label{IV.9.4.1-fr}
\textbf{Notations.} Dans ce numéro et les suivants jusqu'à la fin du \S~9,
nous utiliserons systématiquement les notations suivantes~: étant donné un
morphisme $f:X\to S$, nous poserons, pour tout $s\in S$,
$X_s=f^{-1}(s)=X\times_S\operatorname{Spec}(k(s))$~; pour tout
$\mathcal O_X$-Module quasi-cohérent $\mathcal F$, $\mathcal F_s$ désignera
le $\mathcal O_{X_s}$-Module
$\mathcal F\otimes_{\mathcal O_S}k(s)$, et pour tout homomorphisme
$u:\mathcal F\to\mathcal G$ de $\mathcal F$ dans un $\mathcal O_X$-Module
quasi-cohérent $\mathcal G$, $u_s:\mathcal F_s\to\mathcal G_s$ sera le
morphisme $p^*(u)$, en désignant par $p$ la projection canonique $X_s\to X$.
Pour toute section $g$ de $\mathcal F$ au-dessus de $X$, on désignera par
$g_s$ l'image de $g$ par l'homomorphisme canonique
$\Gamma(X,\mathcal F)\to\Gamma(X_s,\mathcal F_s)$. Pour toute partie $Z$ de
$X$, on notera $Z_s$ l'image réciproque
$p^{-1}(Z)=Z\cap X_s$ (\textbf{I}, 3.6.1). Enfin, si $Y$ est un second
$S$-préschéma et $h:X\to Y$ un $S$-morphisme, on notera $h_s$ le morphisme
$h\times1:X_s\to Y_s$.
\end{env}

\begin{proposition}[9.4.2]
\label{IV.9.4.2-fr}
Soient $S$ un préschéma intègre de point générique $\eta$, $f:X\to S$ un
morphisme de présentation finie, $\mathcal F$, $\mathcal G$, $\mathcal H$
trois $\mathcal O_X$-Modules quasi-cohérents de présentation finie. Soient
$u:\mathcal F\to\mathcal G$, $v:\mathcal G\to\mathcal H$ deux
homomorphismes de $\mathcal O_X$-Modules, et supposons que la suite
$\mathcal F_\eta\xrightarrow{u_\eta}\mathcal G_\eta
\xrightarrow{v_\eta}\mathcal H_\eta$ soit exacte. Alors il existe un
voisinage ouvert $U$ de $\eta$ dans $S$ tel que, pour tout $s\in U$, la suite
$\mathcal F_s\xrightarrow{u_s}\mathcal G_s
\xrightarrow{v_s}\mathcal H_s$ soit exacte.
\end{proposition}

Avec les notations générales de (9.2.1), il s'agit ici de la relation
$\mathbf P(X,\mathcal F,\mathcal G,\mathcal H,u,v,k)$~: «~$X$ est un
préschéma algébrique sur le corps $k$,
$\mathcal F\xrightarrow{u}\mathcal G\xrightarrow{v}\mathcal H$ une suite
exacte de $\mathcal O_X$-Modules quasi-cohérents~». Comme, pour toute extension
$k'$ de $k$, la projection canonique $X_{(k')}\to X$ est un morphisme
fidèlement plat (2.2.13), la condition (9.2.1, (i)) est satisfaite (2.2.7). En
vertu de (9.2.3), on peut se borner au cas où $S$ est intègre et
\emph{noethérien}, auquel cas $X$ est noethérien, et $\mathcal F$,
$\mathcal G$, $\mathcal H$ sont des $\mathcal O_X$-Modules cohérents.
L'hypothèse implique qu'il existe un voisinage ouvert $U$ de $\eta$ dans $S$
tel que la suite
$\mathcal F|f^{-1}(U)\to\mathcal G|f^{-1}(U)\to\mathcal H|f^{-1}(U)$ soit
\emph{exacte}, en vertu de (8.5.8, (ii)) appliqué suivant la méthode générale
de (8.1.2, \emph{a)}), et l'on peut donc déjà supposer que la suite
$\mathcal F\xrightarrow{u}\mathcal G\xrightarrow{v}\mathcal H$ est exacte~;
il suffit évidemment de prouver que l'on a
$\operatorname{Ker}(v_s)=(\operatorname{Ker}v)_s$ et
$\operatorname{Im}(u_s)=(\operatorname{Im}u)_s$ pour tout $s$ voisin de
$\eta$ dans $S$~; par suite (tenant compte de ce que les
$\mathcal O_X$-Modules $\mathcal F$, $\mathcal G$, $\mathcal H$ sont
\emph{cohérents} et de $(\mathbf 0_{\mathrm I},5.3.4)$) on est ramené à
prouver la proposition dans le cas particulier où la suite
\[
0\to\mathcal F\xrightarrow{u}\mathcal G\xrightarrow{v}\mathcal H\to0
\]
est exacte. Or, il y a alors un ouvert $U$ dans $S$ contenant $\eta$ et tel
que $\mathcal H|f^{-1}(U)$ soit \emph{plat} sur $U$ (6.9.1)~; il résulte donc
de (2.1.8) que pour tout $s\in U$, la suite
$0\to\mathcal F_s\to\mathcal G_s\to\mathcal H_s\to0$ est exacte, ce qui
achève la démonstration.

\begin{corollary}[9.4.3]
\label{IV.9.4.3-fr}
Soient $S$ un préschéma intègre, de point générique $\eta$, $f:X\to S$ un
morphisme de présentation finie. Soit
$\mathcal L^\bullet=(\mathcal L^i)_{i\in\mathbf Z}$ un complexe de
$\mathcal O_X$-Modules quasi-cohérents de présentation finie. Pour tout $i$,
il existe un voisinage ouvert $U$ de $\eta$ dans $S$ tel que les homomorphismes
canoniques
\begin{equation}
\label{IV.9.4.3.1-fr}
(\mathcal H^i(\mathcal L^\bullet))_s\to
\mathcal H^i(\mathcal L_s^\bullet) \tag{9.4.3.1}
\end{equation}
soient bijectifs pour tout $s\in U$.
\end{corollary}

\oldpage[IV-3]{63}

On peut évidemment se borner à un complexe à trois termes de degrés
$-1$, $0$, $+1$~:
\[
\mathcal M\xrightarrow{u}\mathcal N\xrightarrow{v}\mathcal P
\]
et à $i=0$~; l'homomorphisme à considérer est alors l'homomorphisme canonique
$(\operatorname{Ker}v/\operatorname{Im}u)_s\to
\operatorname{Ker}(v_s)/\operatorname{Im}(u_s)$. Utilisant (8.9.1) et
(8.5.2, (i)), on voit qu'on peut se ramener au cas où $S$ (donc aussi $X$)
est \emph{noethérien}, et par suite $\mathcal M$, $\mathcal N$, $\mathcal P$
des $\mathcal O_X$-Modules \emph{cohérents}~; alors
$\operatorname{Im}(u)$ et $\operatorname{Ker}(v)$ sont aussi cohérents
$(\mathbf 0_{\mathrm I},5.3.4)$ et en outre il existe un voisinage $U$ de
$\eta$ tel que pour $s\in U$, on ait
$\operatorname{Ker}(v_s)=(\operatorname{Ker}(v))_s$ et
$\operatorname{Im}(u_s)=(\operatorname{Im}(u))_s$ (9.4.2)~; la conclusion
résulte alors de (9.4.2) appliqué à la suite exacte
\[
0\to\operatorname{Im}(u)\to\operatorname{Ker}(v)\to
\operatorname{Ker}(v)/\operatorname{Im}(u)\to0,
\]
compte tenu de ce que $\mathcal O_\eta=k(\eta)$ (puisque $S$ est intègre) et
par conséquent la suite
\[
0\to(\operatorname{Im}u)_\eta\to(\operatorname{Ker}v)_\eta\to
(\operatorname{Ker}v/\operatorname{Im}u)_\eta\to0
\]
est exacte.

\begin{proposition}[9.4.4]
\label{IV.9.4.4-fr}
Soient $f:X\to S$ un morphisme de présentation finie, $\mathcal F$,
$\mathcal G$, $\mathcal H$ trois $\mathcal O_X$-Modules quasi-cohérents de
présentation finie. Soient $u:\mathcal F\to\mathcal G$,
$v:\mathcal G\to\mathcal H$ deux homomorphismes de $\mathcal O_X$-Modules.
Alors l'ensemble $E$ des $s\in S$ tels que la suite
$\mathcal F_s\xrightarrow{u_s}\mathcal G_s
\xrightarrow{v_s}\mathcal H_s$ soit exacte est localement constructible.
\end{proposition}

Compte tenu de (9.2.3), il faut établir que la propriété $\mathbf P$ considérée
dans (9.4.2) est constructible. On a déjà remarqué dans la démonstration de
(9.4.2) que la condition (9.2.1, (i)) est satisfaite pour cette propriété, et
il reste à vérifier la condition (9.2.1, (ii)). Supposons donc $S$ intègre
noethérien, de point générique $\eta$, et prouvons que $E$ ou $S-E$ est un
voisinage de $\eta$. Si $\eta\in E$, notre assertion résulte de (9.4.2), et
l'on peut donc se borner au cas où $\eta\notin E$, autrement dit, la suite
$\mathcal F_\eta\xrightarrow{u_\eta}\mathcal G_\eta
\xrightarrow{v_\eta}\mathcal H_\eta$ n'est pas exacte. Distinguons alors
deux cas~:

$1^\circ$ Posons $w=v\circ u$, et supposons d'abord que
$w_\eta=v_\eta\circ u_\eta\neq0$. Comme $\mathcal F$, $\mathcal G$,
$\mathcal H$ sont cohérents, il en est de même de
$\mathcal N=\operatorname{Ker}(w)$ $(\mathbf 0_{\mathrm I},5.3.4)$~; il
résulte donc de (9.4.2) appliqué à la suite exacte
$0\to\mathcal N\to\mathcal F\xrightarrow{w}\mathcal H$ qu'il y a un
voisinage $U$ de $\eta$ dans $S$ tel que, pour $s\in S$,
$\operatorname{Ker}(w_s)=\mathcal N_s$~; en restreignant $S$, on peut donc
supposer cette relation vérifiée pour tout $s\in S$. Soit $j$ l'injection
canonique $\mathcal N\to\mathcal F$, et posons
$\mathcal M=\operatorname{Coker}(j)$~; l'exactitude à droite du foncteur
$\mathcal F\mapsto\mathcal F_s$ entraîne que
$\mathcal M_s=\operatorname{Coker}(j_s)$ pour tout $s\in S$. L'hypothèse
$w_\eta\neq0$ signifie que $\mathcal M_\eta\neq0$~; comme $\mathcal M$ est
cohérent $(\mathbf 0_{\mathrm I},5.3.4)$, il résulte de (1.8.6) qu'il y a un
voisinage ouvert $U$ de $\eta$ dans $S$ tel que $\mathcal M_s\neq0$ pour
$s\in U$, donc $w_s\neq0$ pour $s\in U$, et \emph{a fortiori} $S-E$ est un
voisinage de $\eta$.

$2^\circ$ Supposons que $w_\eta=0$~; en vertu de (8.5.2, (i)), appliqué
suivant la méthode générale de (8.1.2, \emph{a)}), il existe un voisinage
ouvert $U$ de $\eta$ tel que $w|f^{-1}(U)=0$~; remplaçant $S$ par $U$, on peut
déjà supposer $w=0$ dans $X$. Alors
$\mathcal F\xrightarrow{u}\mathcal G\xrightarrow{v}\mathcal H$ est un
complexe à trois termes $\mathcal L^\bullet$, auquel on peut appliquer
(9.4.3)~; par hypothèse on a
$(\mathcal H^0(\mathcal L^\bullet))_\eta=
\mathcal H^0(\mathcal L^\bullet_\eta)\neq0$, et
$\mathcal H^0(\mathcal L)$ est cohérent
$(\mathbf 0_{\mathrm I},5.3.4\text{ et }5.3.3)$, donc il résulte de (1.8.6)
qu'il y a un voisinage ouvert $U$ de $\eta$ tel que
$(\mathcal H^0(\mathcal L^\bullet))_s\neq0$ pour tout $s\in U$~; mais comme
on peut supposer que
$\mathcal H^0(\mathcal L_s^\bullet)=(\mathcal H^0(\mathcal L^\bullet))_s$
pour $s\in U$ par (9.4.3), on voit de nouveau que $S-E$ est un voisinage de
$\eta$ dans $S$.

\oldpage[IV-3]{64}

\begin{corollary}[9.4.5]
\label{IV.9.4.5-fr}
Soient $f:X\to S$ un morphisme de présentation finie, $\mathcal F$,
$\mathcal G$ deux $\mathcal O_X$-Modules quasi-cohérents de présentation finie,
$u:\mathcal F\to\mathcal G$ un homomorphisme de $\mathcal O_X$-Modules.
Alors l'ensemble des $s\in S$ tels que $u_s$ soit injectif
\textnormal{(}resp. surjectif, bijectif, nul\textnormal{)} est localement
constructible.
\end{corollary}

Il suffit d'appliquer (9.4.4) aux suites
$0\to\mathcal F\xrightarrow{u}\mathcal G$,
$\mathcal F\xrightarrow{u}\mathcal G\to0$,
$0\to\mathcal F\xrightarrow{u}\mathcal G\to0$,
$\mathcal F\xrightarrow{u}\mathcal G\xrightarrow{1_\mathcal G}\mathcal G$.

\begin{corollary}[9.4.6]
\label{IV.9.4.6-fr}
Soient $f:X\to S$ un morphisme de présentation finie, $\mathcal F$ un
$\mathcal O_X$-Module quasi-cohérent de présentation finie. Soit $h$ une
section de $\mathcal F$ au-dessus de $X$~; pour tout $s\in S$, soit $h_s$ la
section correspondante de $\mathcal F_s$ au-dessus de $X_s$ \textnormal{(}pour
le morphisme projection $X_s\to X$\textnormal{)}. Alors, l'ensemble des
$s\in S$ tels que $h_s=0$ est localement constructible.
\end{corollary}

Il suffit de remarquer que $h$ correspond à un homomorphisme
$u:\mathcal O_X\to\mathcal F$ $(\mathbf 0_{\mathrm I},5.1.1)$ et $h_s$ à
l'homomorphisme $u_s$.

\begin{proposition}[9.4.7]
\label{IV.9.4.7-fr}
Soient $f:X\to S$ un morphisme de présentation finie, $\mathcal F$ un
$\mathcal O_X$-Module quasi-cohérent de présentation finie. L'ensemble $E$
\textnormal{(}resp. $E'$\textnormal{)} des $s\in S$ tels que $\mathcal F_s$
soit un $\mathcal O_{X_s}$-Module localement libre \textnormal{(}resp.
localement libre de rang $n$\textnormal{)} est localement constructible.
\end{proposition}

Si $X$ est un préschéma algébrique sur un corps $k$, $\mathcal F$ un
$\mathcal O_X$-Module cohérent, $k'$ une extension de $k$, alors, pour que
$\mathcal F$ soit localement libre (resp. localement libre de rang $n$), il
faut et il suffit qu'il en soit de même de $\mathcal F\otimes_k k'$, puisque la
projection $X_{(k')}\to X$ est un morphisme fidèlement plat (2.2.7). Autrement
dit, la condition (9.2.1, (i)) est vérifiée pour les propriétés dont on veut
démontrer la constructibilité, et il reste à vérifier (9.2.1, (ii))~; on peut
donc encore supposer que $S$ est affine, noethérien et intègre. Il y a de
nouveau quatre cas à envisager~:

$1^\circ$ $\eta\in E$. Il résulte de (8.5.5), appliqué suivant la méthode
générale de (8.1.2, \emph{a)}), qu'il existe un voisinage ouvert $U$ de $\eta$
dans $S$ tel que $\mathcal F|f^{-1}(U)$ soit localement libre~;
\emph{a fortiori} $\mathcal F_s$ est localement libre pour tout $s\in U$.

$2^\circ$ $\eta\in E'$. Même raisonnement que dans le $1^\circ$.

$3^\circ$ $\eta\in S-E$. Comme $\mathcal F_\eta$ est un
$\mathcal O_{X_\eta}$-Module cohérent, dire qu'il n'est pas localement libre
équivaut à dire qu'il \emph{n'est pas plat} sur $\mathcal O_{X_\eta}$
(Bourbaki, \emph{Alg. comm.}, chap.~II, \S~5, n\textsuperscript{o}~2,
cor.~2 du th.~1). Le fait que $S-E$ est un voisinage de $\eta$ résultera donc
du lemme plus général suivant (appliqué au cas où $g$ est l'identité)~:

\begin{lemma}[9.4.7.1]
\label{IV.9.4.7.1-fr}
Soient $S$ un préschéma intègre noethérien, de point générique $\eta$, $X$,
$Y$ deux $S$-préschémas de type fini sur $S$, $g:X\to Y$ un $S$-morphisme,
$\mathcal F$ un $\mathcal O_X$-Module cohérent. Si $\mathcal F_\eta$ n'est
pas $g_\eta$-plat, alors il existe un voisinage ouvert $U$ de $\eta$ dans $S$
tel que pour tout $s\in U$, $\mathcal F_s$ ne soit pas $g_s$-plat.
\end{lemma}

Compte tenu de (2.1.2) et de Bourbaki, \emph{Alg. comm.}, chap.~I\textsuperscript{er},
\S~2, n\textsuperscript{o}~3, Remarque~1, l'hypothèse signifie qu'il existe un
ouvert non vide $V$ de $Y_\eta$, et un homomorphisme injectif
$v:\mathcal M\to\mathcal N$ de $\mathcal O_V$-Modules \emph{cohérents}, tel
que l'homomorphisme
$1\otimes v:\mathcal F_\eta\otimes_{\mathcal O_V}\mathcal M\to
\mathcal F_\eta\otimes_{\mathcal O_V}\mathcal N$ ne soit pas injectif. On a
$V=Y_\eta\cap W$, où $W$ est ouvert dans $Y$ (\textbf{I}, 3.6.1), et il
résulte de (8.5.2, (i) et (ii)), appliqué suivant la méthode de
(8.1.2, \emph{a)}), qu'il existe un voisinage ouvert $U_0$ de $\eta$ dans $S$,
deux $\mathcal O_Z$-Modules

\oldpage[IV-3]{65}

cohérents $\mathcal G$, $\mathcal H$ (où $Z=W\cap h^{-1}(U_0)$,
$h:Y\to S$ étant le morphisme structural) et un $\mathcal O_Z$-homomorphisme
$u:\mathcal G\to\mathcal H$ tel que $\mathcal M=\mathcal G_\eta$,
$\mathcal N=\mathcal H_\eta$ et $v=u_\eta$~; on peut donc supposer $U_0$ pris
tel que, pour $s\in U_0$, $u_s:\mathcal G_s\to\mathcal H_s$ soit injectif
(9.4.5). Mais pour tout $s\in U_0$, l'homomorphisme
$1\otimes u_s:\mathcal F_s\otimes_{\mathcal O_{Z_s}}\mathcal G_s\to
\mathcal F_s\otimes_{\mathcal O_{Z_s}}\mathcal H_s$ n'est autre que
$(1\otimes u)_s$~; l'hypothèse que $(1\otimes u)_\eta$ est non injectif
entraîne donc (9.4.5) l'existence d'un ouvert non vide $U\subset U_0$ tel que
pour tout $s\in U$, $(1\otimes u)_s$ soit non injectif, et par suite
$\mathcal F_s$ n'est pas $g_s$-plat pour tout $s\in U$.

$4^\circ$ $\eta\in S-E'$. Il est clair que $S-E\subset S-E'$, et si
$\eta\in S-E$, $S-E'$ est \emph{a fortiori} un voisinage de $\eta$ d'après le
$3^\circ$. Supposons donc $\eta\in E$, donc $\mathcal F_\eta$ localement
libre~; ces hypothèses entraînent que $X_\eta$ est non connexe, et que les
rangs du $\mathcal O_{X_\eta}$-Module localement libre $\mathcal F_\eta$ ne
sont pas les mêmes sur les diverses composantes connexes de $X_\eta$. Or, il
résulte de (8.4.2), appliqué suivant la méthode de (8.1.2, \emph{a)}), que l'on
peut supposer (en remplaçant $S$ par un voisinage ouvert de $\eta$) que $X$ et
$X_\eta$ ont le même nombre de composantes connexes, les composantes connexes
de $X_\eta$ étant les intersections de $X_\eta$ et des composantes connexes de
$X$. La conclusion résulte alors du raisonnement fait dans le $2^\circ$,
appliqué à chacune des composantes connexes de $X$ (qui sont en nombre fini).

\begin{remark}[9.4.7.2]
\label{IV.9.4.7.2-fr}
Le lemme (9.4.7.1) sera plus tard généralisé et débarrassé d'hypothèses
noethériennes (11.2.8).
\end{remark}

\begin{proposition}[9.4.8]
\label{IV.9.4.8-fr}
Soient $S$ un préschéma localement noethérien, $f:X\to S$ un morphisme de type
fini, $\mathcal F$ un $\mathcal O_X$-Module cohérent. On suppose que pour tout
$s\in S$, $X_s$ soit un préschéma localement intègre. Alors l'ensemble $E$
\textnormal{(}resp. $E'$\textnormal{)} des $s\in S$ tels que $\mathcal F_s$
soit un $\mathcal O_{X_s}$-Module de torsion \textnormal{(}resp. un
$\mathcal O_{X_s}$-Module sans torsion\textnormal{)} est localement
constructible.
\end{proposition}

On peut évidemment supposer $S=\operatorname{Spec}(A)$ affine et noethérien et
prouver que $E$ (resp. $E'$) est alors constructible en utilisant le critère
$(\mathbf 0_{\mathrm{III}},9.2.3)$~; remplaçant $S$ par le sous-préschéma fermé
réduit de $S$ ayant pour espace sous-jacent une partie fermée irréductible $Y$
de $S$, et notant (\textbf{I}, 3.6.4) que pour $s\in Y$, la fibre
$(X_{(Y)})_s$ s'identifie canoniquement à $X_s$ et le faisceau
$(\mathcal F\otimes_{\mathcal O_S}\mathcal O_Y)_s$ à $\mathcal F_s$, on voit
qu'on est ramené au cas où $S$ est intègre de point générique $\eta$, et à
prouver que $E$ ou $S-E$ (resp. $E'$ ou $S-E'$) est un voisinage de $\eta$
dans $S$. Notons en outre que $X$ est réunion finie d'ouverts affines $U_i$,
de type fini sur $S$, et chacun des $(U_i)_s$ est induit sur un ouvert de
$X_s$, donc localement intègre~; en outre, si $(U_i)_\eta$ est vide (resp. non
vide), on sait que $(U_i)_s$ est aussi vide (resp. non vide) dans un voisinage
de $\eta$ (9.2.6). On peut donc supposer tous les $(U_i)_\eta$ non vides et
intègres, et dire que $\mathcal F_s$ est de torsion (resp. sans torsion)
équivaut à dire que chacun des $\mathcal F_s|(U_i)_s$ est de torsion (resp.
sans torsion). On est donc ramené au cas où $X=\operatorname{Spec}(B)$ est
affine, et $\mathcal F=\widetilde M$, où $M$ est un $B$-module de type fini~;
on posera $B_s=B\otimes_A k(s)$, $M_s=M\otimes_A k(s)$ et l'on peut supposer
$B_\eta$ intègre. Nous avons quatre cas à envisager~:

$1^\circ$ $\eta\in E$~; $M_\eta$ est donc un $B_\eta$-module de torsion de
type fini, et il y a par suite un $h\neq0$ dans $B_\eta$ tel que
$hM_\eta=0$~; en vertu de (8.5.2, (i)), appliqué suivant la méthode

\oldpage[IV-3]{66}

de (8.1.2, \emph{a)}), on peut (en remplaçant au besoin $S$ par un voisinage
de $\eta$) supposer que $h\in\Gamma(X_\eta,\mathcal O_{X_\eta})$ est de la
forme $g_\eta$, où $g\in\Gamma(X,\mathcal O_X)$~; soit
$u:\mathcal F\to\mathcal F$ l'endomorphisme de $\mathcal F$ défini par la
multiplication par $g$~; par hypothèse, on a $u_\eta=0$, donc (9.4.5)
l'endomorphisme $u_s:\mathcal F_s\to\mathcal F_s$ défini par la multiplication
par $g_s$, est nul dans un voisinage de $\eta$. D'autre part, soit
$v:\mathcal O_X\to\mathcal O_X$ l'endomorphisme défini par la multiplication
par $g$~; comme $B_\eta$ est intègre et $h=g_\eta\neq0$, $v_\eta$ est injectif,
et il résulte donc de (9.4.5) que $v_s$ est un endomorphisme injectif de
$\mathcal O_{X_s}$ pour $s$ voisin de $\eta$, autrement dit, $g_s$ est un
élément $\mathcal O_{X_s}$-régulier pour ces valeurs de $s$~; donc
$\mathcal F_s$ est de torsion dans un voisinage de $\eta$.

$2^\circ$ $\eta\in S-E$. Dire qu'un $B_\eta$-module de type fini $M_\eta$
n'est pas un module de torsion signifie que son quotient $M_\eta/T$ par son
sous-module de torsion est $\neq0$, et comme c'est un $B_\eta$-module sans
torsion de type fini, il est isomorphe à un sous-module d'un $B_\eta$-module
$B_\eta^n$~; il existe par suite un homomorphisme
$w:M_\eta\to B_\eta$ qui est $\neq0$. Appliquant (8.5.2, (i)) suivant la
méthode de (8.1.2, \emph{a)}), on en déduit (en remplaçant au besoin $S$ par un
voisinage de $\eta$) qu'il existe un homomorphisme
$v:\mathcal F\to\mathcal O_X$ tel que $v_\eta=\widetilde w$. L'hypothèse
$v_\eta\neq0$ entraîne donc (9.4.5) que $v_s\neq0$ dans un voisinage de
$\eta$, et comme $X_s$ est localement intègre, $\mathcal F_s$ n'est pas de
torsion pour ces valeurs de $s$.

$3^\circ$ $\eta\in E'$. Comme $M_\eta$ est un $B_\eta$-module de type fini
sans torsion, il existe un homomorphisme injectif
$w:M_\eta\to B_\eta^n$. Utilisant (8.5.2, (i)) et (9.4.5) comme dans le
$2^\circ$ (en restreignant au besoin $S$), on en déduit qu'il existe un
homomorphisme $v:\mathcal F\to\mathcal O_X^n$ tel que $v_\eta=\widetilde w$,
et que pour $s$ voisin de $\eta$,
$v_s:\mathcal F_s\to\mathcal O_{X_s}^n$ est injectif~; pour ces valeurs de
$s$, $\mathcal F_s$ est donc sans torsion.

$4^\circ$ $\eta\in S-E'$. Soit $T$ le sous-module de torsion de $M_\eta$~;
par hypothèse $T\neq0$, et $T$ est de type fini puisque $M_\eta$ est
noethérien. Utilisant cette fois (8.5.2, (i) et (ii)) on voit (en restreignant
au besoin $S$) qu'il existe un $\mathcal O_X$-Module cohérent $\mathcal G$ et
un homomorphisme injectif $u:\mathcal G\to\mathcal F$ tels que
$\mathcal G_\eta=\widetilde T$ et que $u_\eta$ soit l'injection canonique
$\widetilde T\to\mathcal F_\eta$. Il résulte alors du $1^\circ$ et de (1.8.6)
que dans un voisinage de $\eta$, $\mathcal G_s$ est un
$\mathcal O_{X_s}$-Module de torsion $\neq0$, et d'autre part il résulte de
(9.4.5) que dans un voisinage de $\eta$, $u_s$ est injectif. On en conclut que
dans un voisinage de $\eta$, le sous-Module de torsion de $\mathcal F_s$ n'est
pas nul. C.Q.F.D.

\begin{remark}[9.4.9]
\label{IV.9.4.9-fr}
La propriété «~$X$ est un $k$-préschéma algébrique localement intègre~» ne
vérifie pas la condition (9.2.1, (i)), et il n'est donc pas certain que
l'énoncé (9.4.8) soit encore valable quand on ne fait aucune hypothèse sur $S$
et qu'on suppose seulement que $f$ est un morphisme de présentation finie et
$\mathcal F$ un $\mathcal O_X$-Module de présentation finie. Considérons
toutefois le cas particulier suivant~: $S_0$ étant un préschéma localement
noethérien, soient $f_0:X_0\to S_0$ un morphisme de type fini, tel que les
fibres $(X_0)_{s_0}$ soient localement intègres (pour tout $s_0\in S_0$), et
$\mathcal F_0$ un $\mathcal O_{X_0}$-Module cohérent~; soit $g:S\to S_0$ un
morphisme quelconque, posons
$X=X_0\times_{S_0}S$, $\mathcal F=\mathcal F_0\otimes_{\mathcal O_{S_0}}
\mathcal O_S$, $f=f_0\times1_S:X\to S$, et \emph{supposons} que pour tout
$s\in S$, la fibre $X_s$ soit encore localement intègre. Alors l'ensemble $E$
(resp. $E'$) des $s\in S$ tels que $\mathcal F_s$ soit de torsion (resp. sans

\oldpage[IV-3]{67}

torsion) est encore \emph{localement constructible}. En effet, soit $s\in S$
et soit $s_0=g(s)$~; il suffira (compte tenu de (1.8.2)) de prouver que, pour
que $\mathcal F_s$ soit de torsion (resp. sans torsion), il faut et il suffit
que $(\mathcal F_0)_{s_0}$ le soit. Or,
$\operatorname{Supp}(\mathcal F_s)$ est l'image réciproque de
$\operatorname{Supp}((\mathcal F_0)_{s_0})$ par la projection
$p:X_s\to(X_0)_{s_0}$ (\textbf{I}, 9.1.13)~; comme $p$ est fidèlement plat et
quasi-compact, dire que $\operatorname{Supp}(\mathcal F_s)$ contient un point
maximal de $X_s$ équivaut à dire que
$\operatorname{Supp}((\mathcal F_0)_{s_0})$ contient un point maximal de
$(X_0)_{s_0}$ (1.1.5 et 2.3.4)~; d'où notre assertion en ce qui concerne
l'ensemble $E$ (\textbf{I}, 7.4.6). Si le sous-Module de torsion $\mathcal T$
de $(\mathcal F_0)_{s_0}$ n'est pas nul,
$\mathcal T\otimes_{k(s_0)}k(s)$ (qui est de torsion d'après ce qui précède)
n'est pas nul et s'identifie à un sous-Module de $\mathcal F_s$ (2.2.7), donc
le sous-Module de torsion de $\mathcal F_s$ n'est pas nul. Enfin, si
$(\mathcal F_0)_{s_0}$ est sans torsion, on peut supposer (en considérant un
ouvert affine de $(X_0)_{s_0}$) que $(\mathcal F_0)_{s_0}$ est isomorphe à un
sous-Module d'un $\mathcal O_{(X_0)_{s_0}}^n$, donc $\mathcal F_s$ est
isomorphe à un sous-Module d'un $\mathcal O_{X_s}^n$ (2.2.7), et ceci établit
notre assertion concernant $E'$.
\end{remark}

\subsection{Constructibilité de propriétés topologiques}
\label{subsection:IV.9.5-fr}

\begin{proposition}[9.5.1]
\label{IV.9.5.1-fr}
Soient $f:X\to S$ un morphisme de présentation finie, $Z$ une partie
localement constructible de $X$. Alors l'ensemble $E$ des $s\in S$ tels que
$Z_s\neq\varnothing$ est localement constructible.
\end{proposition}

En effet, on a $E=f(Z)$, et il suffit d'appliquer le th. de Chevalley (1.8.4).

\begin{corollary}[9.5.2]
\label{IV.9.5.2-fr}
Si $Z'$, $Z''$ sont deux parties localement constructibles de $X$, l'ensemble
des $s\in S$ tels que $Z'_s\subset Z''_s$ \textnormal{(}resp.
$Z'_s=Z''_s$\textnormal{)} est localement constructible.
\end{corollary}

En effet, la relation $Z'_s\subset Z''_s$ équivaut à
$(Z'\cap(\complement Z''))_s=\varnothing$ et
$Z'\cap\complement Z''$ est localement constructible.

\begin{proposition}[9.5.3]
\label{IV.9.5.3-fr}
Soient $f:X\to S$ un morphisme de présentation finie, $Z$, $Z'$ deux parties
localement constructibles de $X$ telles que $Z\subset Z'$. Alors l'ensemble
$E$ des $s\in S$ tels que $Z_s$ soit dense dans $Z'_s$ est localement
constructible dans $S$.
\end{proposition}

Il faut vérifier les deux conditions de (9.2.2, (iii)). En ce qui concerne la
première, considérons un préschéma $X$ algébrique sur un corps $k$, deux parties
constructibles $Z$, $Z'$ de $X$ telles que $Z\subset Z'$, et une extension
$k'$ de $k$. Alors la projection canonique $p:X_{(k')}\to X$ est un morphisme
fidèlement plat et quasi-compact, et l'on a donc
$p^{-1}(\overline Z)=\overline{p^{-1}(Z)}$ et
$p^{-1}(\overline{Z'})=\overline{p^{-1}(Z')}$ en vertu de (2.3.10)~; comme
$p$ est surjectif, la relation $\overline Z=\overline{Z'}$ est équivalente à
$\overline{p^{-1}(Z)}=\overline{p^{-1}(Z')}$.

Vérifions maintenant la seconde condition, et supposons donc $S$ affine,
noethérien et intègre, de point générique $\eta$. Distinguons deux cas~:

$1^\circ$ $\eta\in S-E$, autrement dit, $Z_\eta$ n'est pas dense dans
$Z'_\eta$~; il existe donc dans $X$ un ouvert $V$ tel que
$V\cap Z_\eta=\varnothing$ et $V\cap Z'_\eta\neq\varnothing$. Comme $X$ est
noethérien, $V$ est localement constructible, donc il en est de même de
$V\cap Z$, et en vertu de (9.5.1), il y a un voisinage $U$ de $\eta$ dans $S$
tel que pour tout $s\in U$, on ait $(V\cap Z)_s=\varnothing$ et
$(V\cap Z')_s\neq\varnothing$~; cela entraîne que $Z_s$ n'est pas dense dans
$Z'_s$ pour $s\in U$, autrement dit $U\subset S-E$.

\oldpage[IV-3]{68}

$2^\circ$ $\eta\in E$, donc $Z_\eta$ est dense dans $Z'_\eta$. Montrons
d'abord que l'on peut supposer $Z'$ \emph{fermé}. En effet, $Z_\eta$ est dense
dans $\overline{Z'_\eta}$ (adhérence prise dans $X_\eta$)~; posons
$V_\eta=X_\eta-\overline{Z_\eta}$, qui est ouvert dans $X_\eta$ et ne
rencontre pas $Z'_\eta$~; on peut supposer $V_\eta$ de la forme $V\cap X_\eta$,
où $V$ est ouvert (donc constructible) dans $X$, et l'hypothèse
$V_\eta\cap Z'_\eta=\varnothing$ entraîne alors
$V_s\cap Z'_s=\varnothing$ pour tout $s$ voisin de $\eta$ en vertu de (9.5.1).
Remplaçant $S$ par un voisinage ouvert de $\eta$, on peut donc supposer que
$V\cap Z'=\varnothing$, donc $V\cap\overline{Z'}=\varnothing$ (adhérence prise
dans $X$), et par suite $(\overline{Z'})_\eta=\overline{Z'_\eta}$, d'où notre
assertion. L'ensemble $Z'$ est alors réunion de ses composantes irréductibles
en nombre fini, et en restreignant encore $S$ à un voisinage de $\eta$, on
peut supposer que toutes les composantes irréductibles $Z'_i$ de $Z'$
rencontrent $X_\eta$, d'où résulte que $X_\eta$ contient les points génériques
des $Z'_i$ $(\mathbf 0_{\mathrm I},2.1.8)$. Dire que $Z_s$ est dense dans
$Z'_s$ équivaut alors à dire que chacun des $(Z\cap Z'_i)_s$ est dense dans
$(Z'_i)_s$, et l'on est ainsi ramené au cas où $Z'$ est irréductible.
Remplaçant alors au besoin $X$ par le sous-préschéma réduit ayant $Z'$ pour
espace sous-jacent, on voit qu'on peut supposer que $Z'=X$ et que $X$ est
\emph{intègre} et domine $S$. Enfin, en recouvrant $X$ par un nombre fini
d'ouverts affines $W_j$ et remplaçant $Z$ par $Z\cap W_j$, on peut supposer
que $X=\operatorname{Spec}(A)$, où $A$ est un anneau noethérien \emph{intègre}.
Comme $X_\eta$ est noethérien intègre et que $Z_\eta$ est constructible dans
$X_\eta$ et dense dans $X_\eta$, $Z_\eta$ contient un ouvert non vide de
$X_\eta$ $(\mathbf 0_{\mathrm{III}},9.2.2)$, que l'on peut supposer de la forme
$(D(t))_\eta$, où $t$ est un élément $\neq0$ de $A$. En remplaçant au besoin
$S$ par un voisinage de $\eta$, on peut en outre, en vertu de la relation
$(D(t))_\eta\subset Z_\eta$, supposer que $D(t)\subset Z$ (9.5.2). Enfin,
comme l'homothétie de rapport $t_\eta$ dans $\mathcal O_{X_\eta}$ est
injective, il résulte de (9.4.5) que pour $s$ voisin de $\eta$, $t_s$ est
$\mathcal O_{X_s}$-régulier, donc $(X_s)_{t_s}$ est dense dans $X_s$, et
\emph{a fortiori} il en est de même de $Z_s$ qui contient $(X_s)_{t_s}$.

\begin{corollary}[9.5.4]
\label{IV.9.5.4-fr}
Soient $f:X\to S$ un morphisme de présentation finie, $Z$ une partie
localement constructible de $X$. L'ensemble $E$ des $s\in S$ tels que $Z_s$
soit fermé \textnormal{(}resp. ouvert, resp. localement fermé\textnormal{)}
dans $X_s$ est localement constructible dans $S$.
\end{corollary}

Dire que $Z_s$ est ouvert dans $X_s$ signifie que
$(X-Z)_s=X_s-Z_s$ est fermé dans $X_s$, et comme $X-Z$ est localement
constructible, on peut se borner à considérer l'ensemble des $s\in S$ tels que
$Z_s$ soit fermé et l'ensemble des $s\in S$ tels que $Z_s$ soit localement
fermé.

Vérifions encore dans chaque cas les deux conditions de (9.2.2, (ii)). La
première résulte du fait que $p:X_{(k')}\to X$ est fidèlement plat et
quasi-compact, et de (2.3.12) et (2.3.14). Vérifions donc la seconde condition,
$S$ étant supposé affine, noethérien et intègre, de point générique $\eta$.
Posons $Z'=\overline Z$~; le même raisonnement que dans (9.5.3) montre que
$Z'_\eta$ est égal à l'adhérence de $Z_\eta$ dans $X_\eta$~; en vertu de
(9.5.3), il y a donc un voisinage $U$ de $\eta$ tel que pour $s\in U$, $Z_s$
soit dense dans $Z'_s$, ce dernier étant fermé dans $X_s$. Dire que $Z_s$ est
fermé dans $X_s$ signifie alors que $Z''_s=\varnothing$, où $Z''=Z'-Z$~; il
résulte donc de (9.5.1) que l'ensemble $E$ des $s\in S$ où
$Z''_s=\varnothing$ est tel que $E$ ou $S-E$ contienne un voisinage de
$\eta$. Dire que $Z_s$ est localement fermé dans $X_s$ signifie que $Z''_s$
est fermé dans $X_s$~; il suffit donc d'appliquer le résultat précédent en
remplaçant $Z$ par $Z''$, qui est localement constructible dans $X$.

\oldpage[IV-3]{69}

\begin{proposition}[9.5.5]
\label{IV.9.5.5-fr}
Soient $f:X\to S$ un morphisme de présentation finie, $Z$ une partie
localement constructible de $X$ telle que, pour tout $s\in S$, $Z_s$ soit
localement fermé dans $X_s$. Pour tout $s\in S$, soit
$D(s)\subset\mathbf Z\cup\{-\infty\}$ l'ensemble des dimensions des
composantes irréductibles de $Z_s$. Alors la fonction $s\mapsto D(s)$ est
localement constructible dans $S$.
\end{proposition}

Soit $\Phi$ une partie finie de $\mathbf Z\cup\{-\infty\}$~; il s'agit de
montrer que l'ensemble des $s\in S$ tels que $D(s)=\Phi$ est localement
constructible~; compte tenu de (9.2.3), nous avons encore à vérifier les deux
conditions de (9.2.2, (iii)).

En ce qui concerne la première, il s'agit de voir que si $X$ est un préschéma
algébrique sur un corps $k$, $Z$ une partie localement fermée de $X$, $k'$ une
extension de $k$, $p:X_{(k')}\to X$ la projection canonique, alors l'ensemble
des dimensions des composantes irréductibles de $Z$ est le même que celui des
dimensions des composantes irréductibles de $p^{-1}(Z)$~; compte tenu de
l'existence d'un sous-préschéma de $X$ ayant $Z$ pour espace sous-jacent
(\textbf{I}, 5.2.1), cela résulte de (4.2.8).

Pour la seconde condition de (9.2.2, (iii)), on est dans le cas où $S$ est
noethérien et intègre de point générique $\eta$, et $f:X\to S$ est un
morphisme de type fini, de sorte que $X$ est noethérien. Le sous-espace
$Z_\eta$ de l'espace noethérien $X_\eta$ est par hypothèse localement fermé,
donc a un nombre fini de composantes irréductibles $Z_{i\eta}$, qui sont
localement fermées dans $X_\eta$. Il existe par suite pour chaque indice $i$
une partie localement fermée $Z_i$ de $X$ telle que $(Z_i)_\eta=Z_{i\eta}$,
donc si $Z'=\bigcup_i Z_i$, on a $Z'_\eta=Z_\eta$. Mais comme $Z$ et $Z'$ sont
localement constructibles, on peut, en remplaçant $S$ par un voisinage de
$\eta$, supposer que $Z'=Z$ (9.5.1). En outre, pour $i\neq j$,
$Z_{i\eta}\cap Z_{j\eta}$ est rare et \emph{fermé} dans $Z_{i\eta}$~; donc
(9.5.3 et 9.5.4), on peut encore supposer, en restreignant $S$, que pour
$j\neq i$, $(Z_i)_s\cap(Z_j)_s$ est rare et fermé dans $(Z_i)_s$. Comme $Z_i$
est localement fermé dans $X$, il y a un sous-préschéma de $X$ ayant $Z_i$
pour espace sous-jacent (et encore noté $Z_i$), qui est de type fini sur
$S$ (\textbf{I}, 6.3.5). Posons
$U_i=Z_i-\bigcup_{j\neq i}(Z_i\cap Z_j)$, qui est ouvert dans $Z_i$ et tel que,
pour tout $s\in S$, $(U_i)_s$ soit ouvert partout dense dans $(Z_i)_s$~; en
outre, par construction, les $U_i$ sont deux à deux sans point commun. Comme
$(Z_i)_s$ est un $k(s)$-préschéma algébrique, l'ensemble des dimensions des
composantes irréductibles de $Z_s=\bigcup_i(Z_i)_s$ est le même que l'ensemble
des dimensions des composantes irréductibles de la réunion des $(U_i)_s$
(4.1.1.3), chacune de ces composantes étant déjà une composante irréductible
de l'un des $(U_i)_s$. On peut donc se borner au cas où $Z=U_i$~; en outre,
comme $Z_\eta$ est alors irréductible, il n'y a qu'une seule composante
irréductible de $Z$ qui rencontre $X_\eta$ $(\mathbf 0_{\mathrm I},2.1.8)$ et
l'on peut évidemment, en restreignant $Z$, supposer $Z$ irréductible. On est
finalement ramené à prouver le cas particulier suivant de (9.5.5)~:

\begin{corollary}[9.5.6]
\label{IV.9.5.6-fr}
Soient $S$ un préschéma noethérien et irréductible de point générique $\eta$,
$X$ un préschéma irréductible, $f:X\to S$ un morphisme dominant de type fini.
Alors il existe un voisinage $U$ de $\eta$ dans $S$ tel que, pour tout
$s\in U$, toutes les composantes irréductibles de $X_s$ soient de dimension
$n=\dim(X_\eta)$.
\end{corollary}

On peut évidemment se borner au cas où $S=\operatorname{Spec}(A)$ est affine,
$A$ étant donc noethérien~; remplaçant $f$ par $f_{\mathrm{red}}$, qui est de
type fini (1.5.4, (vii)), on peut

\oldpage[IV-3]{70}

supposer $A$ intègre et $X$ réduit, donc intègre puisqu'il est irréductible. On
sait (4.1.2) qu'il existe un ouvert $W$ dense dans $X_\eta$ et un
$k(\eta)$-morphisme \emph{fini et surjectif}
$h:W\to\mathbf V((k(\eta))^n)=\operatorname{Spec}(B')$, où
$B'=k(\eta)[T_1,\ldots,T_n]$ ($T_i$ indéterminées). Si $V$ est un ouvert de
$X$ tel que $V\cap X_\eta=W$, on sait (9.5.3) que pour $s$ voisin de $\eta$
dans $S$, $V_s$ est un ouvert dense dans $X_s$, et l'on peut par suite (4.1.1.3)
se borner au cas où $V=X$, $W=X_\eta$. Posons
$B=A[T_1,\ldots,T_n]$ et $Y=\operatorname{Spec}(B)=\mathbf V_A^n$, de sorte
que $\operatorname{Spec}(B')=Y_\eta$~; il résulte de (8.8.2, (i)) et
(8.10.5, (vi) et (x)), appliqués suivant la méthode de (8.1.2, \emph{a)}),
qu'en remplaçant au besoin $S$ par un voisinage de $\eta$, on peut supposer que
$h=g_\eta$, où $g:X\to Y$ est un morphisme fini surjectif~; autrement dit, on
a $X=\operatorname{Spec}(C)$, où $C$ est une $B$-algèbre finie et
l'homomorphisme $B\to C$ correspondant à $g$ est \emph{injectif}~; comme $C$
est un anneau intègre, $C$ est donc un $B$-module de type fini \emph{sans
torsion}, et $C_\eta=C\otimes_A k(\eta)$ est donc un module de type fini sans
torsion sur $B_\eta=B'=B\otimes_A k(\eta)=k(\eta)[T_1,\ldots,T_n]$ (étant un
module de fractions dont les dénominateurs sont dans
$A-\{0\}\subset B-\{0\}$). Il résulte donc de (9.4.8) qu'il y a un voisinage
$U$ de $\eta$ dans $S$ tel que pour $s\in S$,
$C_s=C\otimes_A k(s)$ soit un module de type fini \emph{sans torsion} sur
$B_s=B\otimes_A k(s)=k(s)[T_1,\ldots,T_n]$, et en particulier l'homomorphisme
$g_s:B_s\to C_s$ est injectif. Comme aucun élément de $B_s$ n'est diviseur de
0 dans $C_s$, pour tout idéal premier minimal $\mathfrak p_i$ de $C_s$ (dont
les éléments sont diviseurs de 0 dans $C_s$), on a nécessairement
$\mathfrak p_i\cap B_s=0$, donc l'homomorphisme canonique
$B_s\to C_s/\mathfrak p_i$ est injectif. On en déduit que pour chaque
composante irréductible $Z_i=\operatorname{Spec}(C_s/\mathfrak p_i)$ de $X_s$,
la restriction à $Z_i$ de $g$ est un morphisme fini et dominant $Z_i\to Y_s$,
donc \emph{surjectif} (\textbf{II}, 6.1.10). On conclut par suite de (4.1.2)
que $\dim Z_i=n$, ce qui achève la démonstration.

\begin{remark}[9.5.7]
\label{IV.9.5.7-fr}
On aura soin de noter que sous les hypothèses de (9.5.6) il peut se faire que
$X_s$ \emph{ne soit irréductible pour aucun} $s\neq\eta$ dans un voisinage de
$\eta$, autrement dit, la propriété «~$X$ est un $k$-préschéma algébrique
irréductible~» \emph{n'est pas constructible}. Prenons par exemple
$S=\operatorname{Spec}(k[T])$, où $k$ est un corps algébriquement clos, $T$ une
indéterminée~; on a donc $k(\eta)=K=k(T)$. Soit $L$ une extension finie
séparable de $K$ de degré $>1$, et soit $X$ la \emph{fermeture intégrale} de
$S$ dans $L$ (\textbf{II}, 6.3.4)~; on a donc $X=\operatorname{Spec}(B)$, où
$B$ est la fermeture intégrale de $k[T]$ dans $L$. On sait que $B$ est un
anneau de Dedekind, et que tous les idéaux maximaux de $k[T]$, sauf un nombre
fini, sont non ramifiés dans $B$~; comme en outre le corps résiduel de tout
idéal maximal de $B$ est nécessairement $k$ (puisque c'est une extension finie
de $k$), on voit que pour presque tous les idéaux maximaux $\mathfrak j_s$ de
$k[T]$, $B_s=B/\mathfrak j_sB$ est composé direct de $[L:K]$ corps isomorphes
à $k$, autrement dit $X_s$ n'est pas irréductible, bien que
$X_\eta=\operatorname{Spec}(L)$ le soit. Le même exemple montre que la
propriété «~$X$ est un $k$-préschéma algébrique intègre~» n'est pas
constructible. Enfin il en est de même de la propriété «~$X$ est un
$k$-préschéma algébrique réduit~». Il suffit pour le voir de prendre encore
$S=\operatorname{Spec}(k[T])$, où cette fois $k$ est un corps algébriquement
clos de caractéristique $p>0$, et pour $X$ la fermeture intégrale de $S$ dans
$L=K^{1/p}$ (où $K=k(T)$), de sorte que $X=\operatorname{Spec}(B)$ avec
$B=k[T^{1/p}]$ ($k$ étant parfait)~; tout idéal maximal de $k[T]$ est de la
forme $(T-\alpha)$ avec $\alpha\in k$~; l'unique idéal de $B$ au-dessus de
l'idéal $(T-\alpha)$ est l'idéal principal
$(T^{1/p}-\alpha^{1/p})$ et il est immédiat

\oldpage[IV-3]{71}

que l'anneau quotient $B/(T-\alpha)B$ admet par suite des éléments nilpotents~;
autrement dit, $X_s$ n'est réduit pour aucun $s\neq\eta$, tandis que
$X_\eta=\operatorname{Spec}(L)$ est intègre.

Nous verrons un peu plus loin (9.7) que l'on obtient par contre des propriétés
constructibles lorsqu'on considère les notions «~géométriques~»
correspondantes aux notions de préschéma irréductible, réduit ou intègre (4.5
et 4.6).
\end{remark}

\subsection{Constructibilité de certaines propriétés des morphismes}
\label{subsection:IV.9.6-fr}

\begin{proposition}[9.6.1]
\label{IV.9.6.1-fr}
Soient $X$, $Y$ deux $S$-préschémas de présentation finie, $f:X\to Y$ un
$S$-morphisme. Soit $E$ l'ensemble des $s\in S$ pour lesquels $f_s$ a l'une
des propriétés suivantes~: être~:
\begin{enumerate}[label=(\roman*),widest=xii]
  \item surjectif~;
  \item dominant~;
  \item séparé~;
  \item propre~;
  \item radiciel~;
  \item fini~;
  \item quasi-fini~;
  \item une immersion~;
  \item une immersion fermée~;
  \item une immersion ouverte~;
  \item un isomorphisme~;
  \item un monomorphisme.
\end{enumerate}
Alors $E$ est localement constructible dans $S$.
\end{proposition}

Les assertions de (i), (v) et (vii) ne sont insérées que pour mémoire, ayant
déjà été établies dans (9.3.2).

(ii)~: Notons d'abord que $f$ est de présentation finie (1.6.2, (v)), donc, en
vertu du théorème de Chevalley (1.8.4), $f(X)$ est localement constructible dans
$Y$. Par ailleurs, on a $f_s(X_s)=(f(X))_s$ (\textbf{I}, 3.6.1)~; l'ensemble
des $s$ tels que $f_s$ soit dominant est l'ensemble des $s$ tels que
$(f(X))_s$ soit dense dans $Y_s$~; la conclusion résulte donc dans ce cas de
(9.5.3).

(iii)~: Comme $f:X\to Y$ est de présentation finie, l'immersion diagonale
$\Delta_f:X\to X\times_YX$ est de présentation finie (1.6.2, (iv) et (v))~; il
résulte donc de (1.8.4.1) que $\Delta_f(X)=\Delta_X$ est une partie localement
constructible de $X\times_YX$. Dire que $f_s$ est séparé signifie (compte tenu
de (\textbf{I}, 5.3.4)) que $(\Delta_X)_s$ est fermé dans
$X_s\times_{Y_s}X_s=(X\times_YX)_s$~; la conclusion résulte donc ici de
(9.5.4).

(iv)~: Montrons que la propriété «~$X$ et $Y$ sont des préschémas algébriques
sur un corps $k$, et $f:X\to Y$ un $k$-morphisme propre~» est constructible. La
condition (9.2.1, (i)) est vérifiée en vertu de (2.7.1, (vii)). On peut donc se
borner au cas où $S$ est affine, noethérien et intègre, de point générique
$\eta$, et on doit prouver que $E$ ou $S-E$ est un voisinage de $\eta$.
Supposons d'abord que $\eta\in E$~; il résulte de (8.10.5, (xii)) appliqué
suivant la méthode de (8.1.2, \emph{a)}) qu'en remplaçant $S$ par un voisinage
de $\eta$, on peut

\oldpage[IV-3]{72}

supposer que $f$ est lui-même un morphisme propre~; on sait alors qu'il en est
de même de $f_s$ pour tout $s\in S$ (\textbf{II}, 5.4.2, (iii)). Supposons au
contraire que $\eta\in S-E$, et distinguons deux cas~:

\begin{enumerate}[label=\upshape{$\arabic{enumi}^\circ$}]
  \item Supposons que $f_\eta$ soit non séparé~; alors il résulte de (iii) que
  $f_s$ est non séparé (et \emph{a fortiori} non propre) dans un voisinage de
  $\eta$.

  \item Supposons $f_\eta$ séparé~; $Y$ est réunion finie d'ouverts affines
  $V_i$, et pour que $f_s$ soit propre, il faut et il suffit que chacune de ses
  restrictions $f_s^{-1}((V_i)_s)\to(V_i)_s$ le soit (\textbf{II}, 5.4.1)~;
  on peut donc se borner au cas où $Y$ est \emph{affine}, donc un schéma. Dire
  que $f_\eta$ n'est pas propre signifie (\textbf{II}, 5.6.3) qu'il existe un
  morphisme \emph{de type fini} $h:Z\to Y_\eta$ tel que le morphisme
  $(f_\eta)_{(Z)}:X_\eta\times_{Y_\eta}Z\to Z$ ne soit pas fermé. Comme $Y$ est
  un schéma, on déduit de (8.8.2, (i) et (ii)) (en restreignant au besoin $S$ à
  un voisinage de $\eta$) qu'il existe un morphisme de type fini $g:Y'\to Y$
  tel que $Z$ soit isomorphe à $Y'_\eta$, et que $g_\eta=h$~; si l'on pose
  $X'=X\times_YY'$, $f'=f_{(Y')}:X'\to Y'$, on a
  $(f_\eta)_{(Z)}=f'_\eta$, et par hypothèse il existe donc une partie fermée
  $M'$ de $X'_\eta$ telle que $f'_\eta(M')$ ne soit pas fermé dans $Y'_\eta$.
  Or $M'$ est la trace sur $X'_\eta$ d'une partie fermée $N'$ de $X'$~; comme
  $X'$ est noethérien et $f'$ de type fini, $f'(N')$ est constructible dans
  $Y'$ (1.8.4) et par hypothèse
  $(f'(N'))_\eta=f'_\eta(N'_\eta)=f'_\eta(M')$ n'est pas fermé dans
  $Y'_\eta$. On conclut alors de (9.5.4) qu'il existe un voisinage $U$ de
  $\eta$ dans $S$ tel que pour $s\in U$,
  $(f'(N'))_s=f'_s(N'_s)$ ne soit pas fermé dans $Y'_s$~; autrement dit, le
  morphisme $f'_s$ n'est pas fermé, et \emph{a fortiori} le morphisme $f_s$
  n'est pas propre.
\end{enumerate}

(vi) La propriété est conjonction des propriétés (iv) et (vii) (8.11.1), donc
la proposition résulte dans ce cas de ce qui a déjà été prouvé.

(ix) La vérification de la condition (9.2.1, (i)) résulte de (2.7.1, (xi)). On
peut donc se borner au cas où $S$ est affine, noethérien et intègre de point
générique $\eta$, et à prouver que $E$ ou $S-E$ est un voisinage de $\eta$. Si
$\eta\in E$, on peut (en remplaçant $S$ par un voisinage de $\eta$) supposer
que $f$ est une immersion fermée en vertu de (8.10.5, (iv)), et alors il est
clair que $f_s$ est une immersion fermée pour tout $s\in S$ (\textbf{I},
4.3.2). Supposons donc que $\eta\in S-E$ et distinguons deux cas~:

\begin{enumerate}[label=\upshape{$\arabic{enumi}^\circ$}]
  \item $f_\eta$ n'est pas un morphisme fini. Alors il résulte de (vi) que dans
  un voisinage de $\eta$, $f_s$ n'est pas fini, ni \emph{a fortiori} une
  immersion fermée.

  \item $f_\eta$ est fini~; alors, en vertu de (8.10.5, (x)), on peut supposer
  (en restreignant au besoin $S$) que $f$ lui-même est un morphisme \emph{fini}.
  Dans ce cas $\mathcal A=\mathcal A(X)=f_*(\mathcal O_X)$ est un
  $\mathcal O_Y$-Module cohérent (\textbf{II}, 6.1.3) et
  $f=\mathcal A(u)$, où $u:\mathcal O_Y\to\mathcal A$ est un homomorphisme de
  $\mathcal O_Y$-Algèbres (\textbf{II}, 1.1.2)~; comme $f_\eta$ n'est pas une
  immersion fermée par hypothèse, $u_\eta$ n'est pas surjectif (\textbf{II},
  1.4.10), donc (9.4.5) il existe un voisinage de $\eta$ dans lequel $u_s$
  n'est pas surjectif, et par suite (\textbf{I}, 4.2.3) $f_s$ n'est pas une
  immersion fermée.
\end{enumerate}

(viii) La vérification de la condition (9.2.1, (i)) se fait comme dans (ix), en
utilisant cette fois (2.7.1, (x)) et le fait que toute immersion d'un préschéma
noethérien dans un autre est quasi-compacte. On est donc ramené au cas où $S$
est affine, noethérien et intègre de point générique $\eta$, et à prouver que
$E$ ou $S-E$ est un voisinage de $\eta$. Si $\eta\in E$, on conclut comme dans
(ix) à l'aide de (8.10.5, (ii)). Si $\eta\in S-E$, on distingue de nouveau
deux cas~:

\begin{enumerate}[label=\upshape{$\arabic{enumi}^\circ$}]
  \item $f_\eta(X_\eta)$ n'est pas une partie localement fermée de $Y_\eta$.
  Comme $f(X)$ est constructible dans $Y$ (1.8.4) et que
  $f_s(X_s)=(f(X))_s$, on déduit de (9.5.4) que pour $s$ voisin de $\eta$,
  $f_s(X_s)$ n'est pas localement fermé dans $Y_s$, et \emph{a fortiori}
  $f_s$ n'est pas une immersion.

  \item\oldpage[IV-3]{73} $f_\eta(X_\eta)$ est localement fermé dans $Y_\eta$.
  Comme $f(X)$ est constructible dans $Y$ (1.8.4) et qu'il en est de même de
  $\overline{f(X)}$ puisque $Y$ est noethérien, il résulte de (8.3.11) qu'en
  restreignant au besoin $S$, on peut supposer que $f(X)$ est localement fermé
  dans $Y$. Il y a alors un ouvert $V$ de $Y$ contenant $f(X)$ et dans lequel
  $f(X)$ est fermé. Comme $Y$ est noethérien, $V$ est de type fini sur $S$, et
  en remplaçant $Y$ par $V$, on peut donc se ramener au cas où $f(X)$ est
  \emph{fermé} dans $Y$. Mais alors $f_s(X_s)=(f(X))_s$ est fermé dans $Y_s$
  pour tout $s\in S$, et dire que $f_s$ est une immersion équivaut à dire que
  $f_s$ est une immersion fermée~; on est donc ramené à ce qui a été démontré
  dans (ix).
\end{enumerate}

(x) Utilisant cette fois (2.7.1, (ix)) et (8.10.5, (iii)), on est ramené au cas
où $S$ est affine, noethérien, intègre de point générique $\eta$, et où
$\eta\in S-E$. Distinguons trois cas~:

\begin{enumerate}[label=\upshape{$\arabic{enumi}^\circ$}]
  \item $f_\eta(X_\eta)$ n'est pas ouvert dans $Y_\eta$. Comme $f(X)$ est
  constructible dans $Y$ (1.8.4), on déduit de (9.5.4) que $f_s(X_s)$ n'est
  pas ouvert dans $Y_s$ pour $s$ voisin de $\eta$, et \emph{a fortiori} $f_s$
  n'est pas une immersion ouverte.

  \item $f_\eta(X_\eta)$ est ouvert dans $Y_\eta$, mais $f_\eta$ n'est pas une
  immersion. Il résulte alors de (viii) que pour $s$ voisin de $\eta$ dans $S$,
  $f_s$ n'est pas une immersion, ni \emph{a fortiori} une immersion ouverte.

  \item $f_\eta(X_\eta)$ est ouvert dans $Y_\eta$ et $f_\eta$ est une
  immersion. Comme $f(X)$ est constructible dans $Y$, il résulte de (8.3.11)
  qu'en restreignant au besoin $S$, on peut déjà supposer que $f(X)$ est ouvert
  dans $Y$. Comme $Y$ est noethérien, le sous-préschéma induit sur $f(X)$ est
  de type fini sur $S$, donc on peut se ramener au cas où $f$ est \emph{surjectif}
  en remplaçant $Y$ par $f(X)$. Par hypothèse, $f_\eta$ est une immersion
  fermée, donc on peut, comme dans (ix), supposer que $f$ est une immersion
  \emph{fermée}, et par suite que $X$ est un sous-préschéma fermé de $Y$ défini
  par un Idéal cohérent $\mathcal J$ de $\mathcal O_Y$. Par hypothèse
  $f_\eta$ n'est pas un isomorphisme, donc $\mathcal J_\eta\neq0$~; on en
  conclut (9.4.5) que dans un voisinage de $\eta$, on a $\mathcal J_s\neq0$,
  et par suite l'immersion fermée surjective $f_s$ n'est pas ouverte.
\end{enumerate}

(xi) La propriété est conjonction des propriétés (i) et (x) et résulte donc de
ce qui a été démontré.

(xii) En vertu de (\textbf{I}, 5.3.8), dire que $f_s$ est un monomorphisme
signifie que
$\Delta_{f_s}=(\Delta_f)_s:X_s\to X_s\times_{Y_s}X_s=(X\times_YX)_s$ est un
isomorphisme, et comme $X\times_YX$ est un $S$-préschéma de présentation finie
(1.6.2, (iv)), la conclusion résulte de (xi).

\begin{proposition}[9.6.2]
\label{IV.9.6.2-fr}
Soient $X$, $Y$ deux $S$-préschémas de présentation finie, $f:X\to Y$ un
$S$-morphisme.
\begin{enumerate}[label=\upshape{I)}]
  \item Soit $E$ l'ensemble des $s\in S$ pour lesquels $f_s$ a l'une des
  propriétés suivantes~: être~:
  \begin{enumerate}[label=(\roman*)]
    \item affine~;
    \item quasi-affine~;
    \item projectif~;
    \item quasi-projectif.
  \end{enumerate}
  Alors $E$ est ind-constructible dans $S$.
\end{enumerate}

\oldpage[IV-3]{74}

\begin{enumerate}[label=\upshape{II)}]
  \item Soit $\mathcal L$ un $\mathcal O_X$-Module inversible. Alors l'ensemble
  $E'$ des $s\in S$ tels que $\mathcal L_s$ soit un $\mathcal O_{X_s}$-Module
  ample (resp. très ample) relativement à $f_s$, est ind-constructible dans
  $S$.
\end{enumerate}
\end{proposition}

I) Vérifions les conditions (i) et (ii) de (9.2.1). En ce qui concerne la
condition (9.2.1, (i)), elle résulte, pour les propriétés (i) et (ii), de
(2.7.1, (xiii) et (xiv))~; pour les propriétés (iii) et (iv), elle résulte de
(9.1.5). Vérifions ensuite la condition (9.2.1, (ii)), en supposant donc $S$
noethérien, intègre et de point générique $\eta$, et que
$f_\eta:X_\eta\to Y_\eta$ ait l'une des propriétés (i) à (iv) de l'énoncé.
Appliquant (8.10.5, (viii), (ix), (xiii) et (xiv)) suivant la méthode de
(8.1.2, \emph{a)}), on voit d'abord qu'il existe un voisinage ouvert $U$ de
$\eta$ tel que, si $V$ et $W$ sont les images réciproques de $U$ dans $X$ et
$Y$ respectivement par les morphismes structuraux, la restriction $V\to W$ de
$f$ a celle des propriétés (i) à (iv) que l'on considère. La conclusion
résulte alors de ce que ces propriétés sont toutes stables par changement de
base.

II) On procède de la même façon. La condition (9.2.1, (i)) résulte cette fois
de (2.7.2). Pour la condition (9.2.1, (ii)), avec les mêmes notations que dans
I), il résulte de (8.10.5.2) que pour un voisinage $U$ de $\eta$ dans $S$, la
restriction $\mathcal L|V$ est ample (resp. très ample) relativement à la
restriction $V\to W$ de $f$. La conclusion résulte encore de la stabilité des
propriétés considérées par changement de base (\textbf{II}, 4.6.13 et
4.4.10).

On peut améliorer (9.6.2, II) sous certaines conditions~:

\begin{proposition}[9.6.3]
\label{IV.9.6.3-fr}
Soient $X$, $Y$ deux $S$-préschémas de présentation finie, $f:X\to Y$ un
$S$-morphisme propre, $\mathcal L$ un $\mathcal O_X$-Module inversible. Alors
l'ensemble $E$ (resp. $E'$) des $s\in S$ tels que $\mathcal L_s$ soit un
$\mathcal O_{X_s}$-Module ample (resp. très ample) relativement à
$f_s:X_s\to Y_s$ est localement constructible dans $S$.
\end{proposition}

Soient $k$ un corps, $Z$, $T$ deux préschémas algébriques sur $k$, $\mathcal M$
un $\mathcal O_Z$-Module inversible, $g:Z\to T$ un $k$-morphisme de type fini~;
alors, si $\mathbf P(Z,T,\mathcal M,g,k)$ désigne la relation~: «~$\mathcal M$ est ample
(resp. très ample) relativement à $g$~», on a déjà remarqué dans (9.6.2) que
$\mathbf P(Z,T,\mathcal M,g,k)$ vérifie la condition (9.2.1, (i)) en vertu de (2.7.2).
On sait déjà d'autre part que $E$ et $E'$ sont ind-constructibles. Il reste donc
à voir que si $S$ est noethérien, intègre, de point générique $\eta$ et si
$\eta\in S-E$ (resp. $\eta\in S-E'$), alors $S-E$ (resp. $S-E'$) contient un
voisinage de $\eta$. Nous considérerons séparément le cas de $E'$ et celui de
$E$.

I) \emph{Cas de $E'$.} --- Notons que puisque $f$ est séparé et $Y$
quasi-compact, il existe un entier $h$ tel que, pour $q>h$, on ait
$R^qf_*(\mathcal L)=0$ (\textbf{III}, 1.4.12)~; d'autre part, puisque $f$ est
propre et $Y$ noethérien, les $R^qf_*(\mathcal L)$ sont tous des
$\mathcal O_Y$-Modules cohérents (\textbf{III}, 3.2.1)~; comme ils sont nuls
sauf pour un nombre fini de valeurs de $q$, le théorème de platitude générique
(6.9.1) montre qu'en restreignant $S$ à un voisinage de $\eta$, on peut
supposer que $\mathcal L$ et les $R^qf_*(\mathcal L)$ sont tous $S$-plats. On
conclut alors de (\textbf{III}, 6.9.9) que l'homomorphisme canonique
\begin{equation}
\label{IV.9.6.3.1-fr}
 f_*(\mathcal L)\otimes_{\mathcal O_S} k(s)
 \longrightarrow (f_s)_*(\mathcal L_s)
 \tag{9.6.3.1}
\end{equation}
est un \emph{isomorphisme}.

\oldpage[IV-3]{75}

Cela étant, il résulte de (\textbf{II}, 4.4.4) que dire que $\mathcal L_s$
n'est pas très ample relativement à $f_s$ signifie que~: \emph{ou bien}
l'homomorphisme canonique
$(f_s)^*((f_s)_*(\mathcal L_s))\to\mathcal L_s$ n'est pas surjectif~; \emph{ou
bien} l'homomorphisme précédent est surjectif et le morphisme canonique
$r:X_s\to\mathbf P((f_s)_*(\mathcal L_s))$ n'est pas une immersion. Compte
tenu de l'isomorphisme (9.6.3.1), ces conditions s'écrivent respectivement sous
la forme~: $1^\circ$ l'homomorphisme canonique
$(f^*(f_*(\mathcal L)))_s\to\mathcal L_s$ n'est pas surjectif~; $2^\circ$
l'homomorphisme précédent est surjectif et le morphisme canonique
$X_s\to\mathbf P((f_*(\mathcal L))_s)$ n'est pas une immersion.

Supposons d'abord que l'homomorphisme canonique
$(f^*(f_*(\mathcal L)))_\eta\to\mathcal L_\eta$ ne soit pas surjectif. Comme
$f_*(\mathcal L)$ est cohérent, il en est de même de $f^*(f_*(\mathcal L))$,
et alors il résulte de (9.4.5) que pour tout $s$ dans un voisinage de $\eta$,
l'homomorphisme $(f^*(f_*(\mathcal L)))_s\to\mathcal L_s$ n'est pas
surjectif, ce qui prouve dans ce cas que $S-E'$ est un voisinage de $\eta$.

Supposons en second lieu que l'homomorphisme canonique
$(f^*(f_*(\mathcal L)))_\eta\to\mathcal L_\eta$ soit surjectif mais que le
morphisme $X_\eta\to\mathbf P((f_*(\mathcal L))_\eta)$ ne soit pas une
immersion. Alors le même raisonnement que ci-dessus montre d'abord que pour
tout $s$ assez voisin de $\eta$, l'homomorphisme
$(f^*(f_*(\mathcal L)))_s\to\mathcal L_s$ est surjectif~; d'autre part, en
vertu de (9.6.1, (viii)), pour $s$ assez voisin de $\eta$, le morphisme
$X_s\to\mathbf P((f_*(\mathcal L))_s)$ n'est pas une immersion, ce qui achève
la démonstration dans le cas de $E'$.

II) \emph{Cas de $E$.} --- Considérons d'abord le cas particulier où $Y=S$~:

\begin{corollary}[9.6.4]
\label{IV.9.6.4-fr}
Soit $f:X\to S$ un morphisme propre de présentation finie, $\mathcal L$ un
$\mathcal O_X$-Module inversible. Alors, l'ensemble $E$ des $s\in S$ tels que
$\mathcal L_s$ soit ample (relativement à $f_s$) est ouvert dans $S$, et
$\mathcal L|f^{-1}(E)$ est ample relativement à la restriction
$f^{-1}(E)\to E$ de $f$.
\end{corollary}

Comme la condition (9.2.1, (i)) est vérifiée par la propriété $\mathbf P$ définie plus
haut, le résultat de (9.2.2, (iv)) et le raisonnement de (9.2.3) montrent qu'on
peut se borner au cas où $S$ est \emph{noethérien}~; mais alors le résultat
découle de (\textbf{III}, 4.7.1) et de la stabilité de l'amplitude par
changement de base (\textbf{II}, 4.6.13).

\begin{corollary}[9.6.5]
\label{IV.9.6.5-fr}
Sous les hypothèses de (9.6.4), pour que $\mathcal L$ soit ample relativement
à $f$, il faut et il suffit que, pour tout $s\in S$, $\mathcal L_s$ soit ample
relativement à $f_s$.
\end{corollary}

\phantomsection\label{IV.9.6.6-fr}%
\textbf{(9.6.6) Fin de la démonstration de (9.6.3).} --- Revenons au cas
général, $S$ étant noethérien, intègre et de point générique
$\eta\in S-E$. Comme $f$ est propre, il résulte de (9.6.4) que l'ensemble $V$
des $y\in Y$ tels que $\mathcal L_y$ soit un $\mathcal O_{X_y}$-Module ample,
relativement au morphisme $f_y:X_y\to\operatorname{Spec}(k(y))$ est ouvert,
donc $F=Y-V$ est fermé dans $Y$. Cela étant, comme $f_s$ est propre et que,
pour tout $y\in Y$ au-dessus de $s\in S$, on a
$\mathcal L_y=(\mathcal L_s)_y$, il résulte de (9.6.5) que pour que $s$
appartienne à $S-E$, il faut et il suffit que l'on ait $F_s\neq\varnothing$.
Mais comme $F$ est fermé dans $Y$ et que $F_\eta\neq\varnothing$, il résulte
de (9.5.1) que l'ensemble des $s\in S$ tels que $F_s\neq\varnothing$ est un
voisinage de $\eta$ dans $S$. C.Q.F.D.

\begin{remarks}[9.6.7]
\label{IV.9.6.7-fr}
(i) Pour chacune des propriétés $\mathbf P$ considérées dans (9.6.1) la proposition
(9.3.3) est applicable, et ces propriétés (pour les morphismes $f_s$) sont donc
«~stables~» par passage d'une limite projective d'un système projectif
essentiellement affine (8.13.4) de préschémas $X_\lambda$ à l'un convenable
d'entre eux.

(ii) Soit $Z$ une partie \emph{localement constructible} de $X$ telle que, pour
tout $s\in S$, $Z\cap X_s$ soit \emph{ouvert} dans $X_s$, et désignons par
$Z_s$ le sous-préschéma de $X_s$ induit sur l'ouvert $Z\cap X_s$.

\oldpage[IV-3]{76}

Alors, dans les propositions (9.6.1) et (9.6.2, (I)), on peut partout remplacer
$f_s$ par sa restriction $f_s|Z_s:Z_s\to Y_s$ sans changer les conclusions. En
effet, la vérification de (9.2.1, (i)) se fait comme dans (9.6.1) et (9.6.2).
D'autre part, dans la réduction au cas où $S$ est noethérien, faite dans
(9.2.3), si $Z=q^{-1}(Z_0)$, où $q:X\to X_0$ est la projection canonique et
$Z_0$ une partie constructible de $X_0$ (8.3.11), le fait que
$(Z_0)_{s_0}$ soit ouvert dans $(X_0)_{s_0}$, pour $s_0=p(s)$, découle de
(2.4.10) et de ce que la projection $X_s\to(X_0)_{s_0}$ est surjective. On est
donc ramené à vérifier (9.2.1, (ii)) sous les nouvelles hypothèses. Or, comme
$Z_\eta$ est ouvert dans $X_\eta$, il existe un ouvert $Z'\subset X$ tel que
$Z_\eta=Z'\cap X_\eta$~; comme $X$ est alors noethérien, $Z'$ est
constructible, et il en est de même de $Z$ par hypothèse~; on conclut donc de
(9.5.2) et de ($\mathbf 0_{\mathrm{III}}$, 9.2.2) qu'il y a un voisinage $U$
de $\eta$ dans $S$ tel que $Z_s=Z'_s$ pour $s\in U$. Remplaçant $S$ par $U$,
on peut donc se borner au cas où $Z$ est ouvert dans $X$, et alors on est
ramené à ce qui a été prouvé dans (9.6.1) et (9.6.2).
\end{remarks}

\subsection{Constructibilité des propriétés de séparabilité, d'irréductibilité
géométrique et de connexité géométrique}
\label{subsection:IV.9.7-fr}

\begin{lemma}[9.7.1]
\label{IV.9.7.1-fr}
Soient $S$ un préschéma irréductible de point générique $\eta$, $f:X\to S$
un morphisme de présentation finie. Si $X_\eta$ a $n$ composantes irréductibles
(resp. $n'$ composantes connexes), il existe un voisinage ouvert $U$ de $\eta$
dans $S$ tel que pour tout $s\in U$, $X_s$ ait au moins $n$ composantes
irréductibles (resp. $n'$ composantes connexes).
\end{lemma}

On peut se borner au cas où $S$ est affine, donc $X$ quasi-compact et
quasi-séparé. Soient $V_i$ ($1\leqslant i\leqslant n$) les intérieurs des
composantes irréductibles de $X_\eta$, qui sont deux à deux disjoints et
quasi-compacts, et dont la réunion est dense dans $X_\eta$~; en vertu de
(8.2.11), appliqué suivant la méthode de (8.1.2, \emph{a)}), il existe pour
chaque $i$ un ouvert quasi-compact $W_i$ dans $X$ tel que
$W_i\cap X_\eta=V_i$~; comme $X$ est quasi-séparé, les intersections
$W_i\cap W_j$ sont quasi-compactes (1.2.7), donc, en remplaçant $S$
par un voisinage de $\eta$, on peut supposer que $W_i\cap W_j=\varnothing$ pour
$i\neq j$ en vertu de (8.3.3). En outre, comme les $W_i$ sont constructibles,
il y a un voisinage $U$ de $\eta$ dans $S$ tel que pour tout $s\in U$ les
$(W_i)_s$ soient non vides (9.5.1 et 9.2.3) et que la réunion des $(W_i)_s$ soit
\emph{dense} dans $X_s$ (9.5.3 et 9.2.3). Cela étant, les composantes
irréductibles (en nombre fini) des $(W_i)_s$ sont aussi les composantes
irréductibles de la réunion des $(W_i)_s$ $(\mathbf 0_{\mathrm I},2.1.7)$, donc
les adhérences dans $X_s$ de ces composantes sont les composantes
irréductibles de $X_s$ $(\mathbf 0_{\mathrm I},2.1.6)$ et leur nombre est
évidemment $\geqslant n$.

Soient maintenant $C_j$ ($1\leqslant j\leqslant n'$) les composantes connexes
de $X_\eta$, qui sont des ouverts quasi-compacts de $X_\eta$, deux à deux
disjoints. Si on remplace les $V_i$ par les $C_j$ dans le raisonnement
précédent, on voit (utilisant deux fois (8.3.3)) qu'on peut supposer $X$ réunion
de $n'$ ouverts quasi-compacts $M_j$ ($1\leqslant j\leqslant n'$) deux à deux
disjoints et que, dans un voisinage de $\eta$, les $(M_j)_s$ soient non vides.
Comme la réunion des $(M_j)_s$ est $X_s$, les composantes connexes des $(M_j)_s$
sont les composantes connexes de $X_s$, donc leur nombre est $\geqslant n'$.

\oldpage[IV-3]{77}

\begin{lemma}[9.7.2]
\label{IV.9.7.2-fr}
Soient $S$ un préschéma irréductible de point générique $\eta$, $f:X\to S$
un morphisme de présentation finie. Si $X_\eta$ n'est pas réduit, il existe un
voisinage ouvert $U$ de $\eta$ dans $S$ tel que, pour tout $s\in U$, $X_s$ ne
soit pas réduit.
\end{lemma}

En effet, soit $\mathcal N$ le Nilradical de $\mathcal O_X$~; il résulte de
(8.2.13), appliqué suivant la méthode de (8.1.2, \emph{a)}) que $\mathcal N_\eta$
est le nilradical de $\mathcal O_{X_\eta}$, et l'hypothèse est que
$\mathcal N_\eta\neq0$. On conclut de (9.4.5) et (9.2.3) qu'il y a un voisinage
$U$ de $\eta$ tel que, pour tout $s\in U$, $\mathcal N_s$ s'identifie à un Idéal
de $\mathcal O_{X_s}$ et $\mathcal N_s\neq0$~; comme $\mathcal N_s$ est
évidemment contenu dans le Nilradical de $\mathcal O_{X_s}$, on voit que $X_s$
n'est pas réduit pour $s\in U$.

\phantomsection\label{IV.9.7.3-fr}%
\textbf{(9.7.3)} Étant donné un polynôme
$F\in A[\mathrm T_1,\ldots,\mathrm T_n]$, où $A$ est un anneau et les
$\mathrm T_i$ des indéterminées, pour tout homomorphisme d'anneaux
$\rho:A\to B$, nous désignerons par $F_{(\rho)}$ ou $F_{(B)}$ le polynôme de
$B[\mathrm T_1,\ldots,\mathrm T_n]$ obtenu en remplaçant chaque coefficient de
$F$ par son image par $\rho$. Si $k$ est un corps,
$F\in k[\mathrm T_1,\ldots,\mathrm T_n]$ un polynôme non constant et
$X=\operatorname{Spec}(k[\mathrm T_1,\ldots,\mathrm T_n]/(F))$, dire que $X$
est intègre (ou que l'idéal $(F)$ est premier) signifie que $F$ est
\emph{irréductible} (c'est-à-dire que dans toute factorisation $F=F_1F_2$ en
polynômes de $k[\mathrm T_1,\ldots,\mathrm T_n]$, $F_1$ ou $F_2$ est de degré
$0$)~; cela résulte de ce que l'anneau $k[\mathrm T_1,\ldots,\mathrm T_n]$ est
factoriel. On en déduit aussitôt ((4.6.2) et (4.5.2))~:

\begin{lemma}[9.7.4]
\label{IV.9.7.4-fr}
Soient $k$ un corps, $\Omega$ une extension algébriquement close de $k$, $F$ un
polynôme non constant de $k[\mathrm T_1,\ldots,\mathrm T_n]$. Les conditions
suivantes sont équivalentes~:
\begin{enumerate}[label=\emph{\alph*)}]
  \item $X=\operatorname{Spec}(k[\mathrm T_1,\ldots,\mathrm T_n]/(F))$ est
  géométriquement intègre.
  \item $F_{(K)}$ est irréductible pour toute extension $K$ de $k$.
  \item $F_{(\Omega)}$ est irréductible.
\end{enumerate}
Nous dirons dans ce cas que $F$ est \emph{géométriquement irréductible}.
\end{lemma}

\begin{lemma}[9.7.5]
\label{IV.9.7.5-fr}
Soient $A$ un anneau intègre, $K$ son corps des fractions, $F$ un polynôme non
constant de $A[\mathrm T_1,\ldots,\mathrm T_n]$ de degré $d$ tel que $F_{(K)}$
soit géométriquement irréductible. Alors il existe $f\neq0$ dans $A$ tel que
pour tout $x\in D(f)$, $F_{(\boldsymbol{k}(x))}$ soit géométriquement
irréductible.
\end{lemma}

Posons $F=\sum_\alpha c_\alpha\mathrm T^\alpha$ avec comme d'ordinaire
$c_\alpha=c_{\alpha_1\alpha_2\ldots\alpha_n}$,
$\mathrm T^\alpha=\mathrm T_1^{\alpha_1}\mathrm T_2^{\alpha_2}\cdots
\mathrm T_n^{\alpha_n}$,
$|\alpha|=\alpha_1+\alpha_2+\cdots+\alpha_n\leqslant d$ (un au moins des
$c_\alpha$ tels que $|\alpha|=d$ n'étant pas nul). Comme les $c_\alpha$ non
nuls sont inversibles dans $K$, on peut supposer, en remplaçant au besoin $A$
par un anneau $A_g$ ($g\neq0$ dans $A$) que les $c_\alpha$ non nuls sont
inversibles dans $A$. Il en résulte que pour tout $x\in\operatorname{Spec}(A)$,
$F_{(\boldsymbol{k}(x))}$ est \emph{de degré $d$}.

Nous démontrerons d'abord un lemme préliminaire.

\begin{lemma}[9.7.5.1]
\label{IV.9.7.5.1-fr}
Soient $A$ un anneau intègre, $S=\operatorname{Spec}(A)$, $\eta$ le point
générique de $S$, $B$ l'anneau de polynômes
$A[\mathrm T_1,\ldots,\mathrm T_m]$, $(P_i)_{i\in I}$ une famille finie
d'éléments de $B$, $\mathfrak a$ l'idéal de $B$ engendré par les $P_i$. Pour
tout $s\in S$, soit $\mathfrak a_s$ l'idéal de
$\boldsymbol{k}(s)[\mathrm T_1,\ldots,\mathrm T_m]$ engendré par les polynômes
$(P_i)_{(\boldsymbol{k}(s))}$ ($i\in I$). Alors, si
$V(\mathfrak a_\eta)=\varnothing$, il existe un voisinage $U$ de $\eta$ dans
$S$ tel que $V(\mathfrak a_s)=\varnothing$ pour tout $s\in U$.
\end{lemma}

En effet, soit $Y=\operatorname{Spec}(B)$, et soit $Z$ la partie fermée
$V(\mathfrak a)$ de $Y$~; comme $\mathfrak a$ est un idéal de type fini, $Z$
est constructible dans $Y$ $(\mathbf 0_{\mathrm{III}},9.1.5)$ et avec les
notations introduites en (9.4.1), on a $V(\mathfrak a_s)=Z_s$ pour tout
$s\in S$~; comme le morphisme structural $Y\to S$ est de présentation finie,
la conclusion du lemme résulte de (9.5.1) et (9.2.3).

\oldpage[IV-3]{78}

Soit $(p,q)$ un couple d'entiers $>0$ tels que $p+q=d$~; introduisons des
indéterminées $\mathrm T'_\beta$, $\mathrm T''_\gamma$ pour tous les systèmes
d'entiers $\beta$, $\gamma$ tels que $|\beta|\leqslant p$ et
$|\gamma|\leqslant q$~; pour tout système d'entiers $\alpha$ tel que
$|\alpha|\leqslant d$, considérons le polynôme de
\[
B=A[\mathrm T'_\beta,\mathrm T''_\gamma]_{|\beta|\leqslant p,
  |\gamma|\leqslant q}:
\qquad
P_\alpha(\mathrm T'_\beta,\mathrm T''_\gamma)
=\sum_{\beta+\gamma=\alpha}\mathrm T'_\beta\mathrm T''_\gamma-c_\alpha.
\]
Soit $\Omega$ une clôture algébrique de $K$~; dire qu'il existe deux polynômes
$F_1$, $F_2$ de $\Omega[\mathrm T_1,\ldots,\mathrm T_n]$, de degrés respectifs
$p$ et $q$, tels que $F_1F_2=F_{(\Omega)}$, signifie que le système d'équations
$P_\alpha(t'_\beta,t''_\gamma)=0$ ($|\alpha|\leqslant d$) admet une solution
$(t'_\beta,t''_\gamma)$ ($|\beta|\leqslant p$, $|\gamma|\leqslant q$) formée
d'éléments de $\Omega$. Soit $\mathfrak a$ l'idéal de $B$ engendré par les
$P_\alpha$~; l'interprétation précédente, et le théorème des zéros de Hilbert,
montrent que l'hypothèse sur $F_{(K)}$ implique que
$V(\mathfrak a_\eta)=\varnothing$, en désignant par $\eta$ le point générique de
$\operatorname{Spec}(A)$~; le lemme (9.7.5.1) prouve donc que dans un voisinage
de $\eta$, on a $V(\mathfrak a_x)=0$, et par suite pour ces valeurs de $x$,
$F_{(\boldsymbol{k}(x))}$ n'admet pas de factorisation
$F_{(\boldsymbol{k}(x))}=G_1G_2$, où $G_1$, $G_2$ sont des polynômes de degrés
respectifs $p$ et $q$, dont les coefficients sont dans une clôture algébrique
de $\boldsymbol{k}(x)$. Il suffit d'appliquer ce résultat à tous les couples
d'entiers $(p,q)$ tels que $p>0$, $q>0$ et $p+q=d$, pour obtenir la conclusion
du lemme (9.7.5).

\begin{proposition}[9.7.6]
\label{IV.9.7.6-fr}
Soient $S$ un préschéma noethérien intègre de point générique $\eta$, $f:X\to S$
un morphisme de type fini, $\mathcal F$ un $\mathcal O_X$-Module cohérent. Si
$\mathcal F_\eta$ est sans cycle premier associé immergé, il existe un voisinage
$U$ de $\eta$ dans $S$ tel que, pour tout $s\in U$, $\mathcal F_s$ soit sans cycle
premier associé immergé.
\end{proposition}

Comme $X_\eta$ est noethérien, il existe un homomorphisme injectif
$v:\mathcal F_\eta\to\bigoplus_{i=1}^m\mathcal G_i$, où chacun des $\mathcal G_i$ est
irrédondant et $\operatorname{Ass}(\mathcal F_\eta)$ est l'ensemble des $x_i$, où
$\{x_i\}=\operatorname{Ass}(\mathcal G_i)$ (3.2.6)~; si $Z_i$ est l'adhérence de
$\{x_i\}$ dans $X_\eta$, on a $Z_i=\operatorname{Supp}(\mathcal G_i)$ (3.1.4), et
par hypothèse les $Z_i$ sont les composantes irréductibles de
$Z=\operatorname{Supp}(\mathcal F_\eta)$. Il résulte de (8.5.2, (i) et (ii)) (en
restreignant au besoin $S$ à un voisinage de $\eta$) et (8.5.8) qu'il existe des
$\mathcal O_X$-Modules cohérents $\mathcal F_i$ et un homomorphisme injectif
$u:\mathcal F\to\bigoplus_{i=1}^m\mathcal F_i$ tels que $(\mathcal F_i)_\eta=\mathcal G_i$ pour
tout $i$ et $v=u_\eta$. Si $Y_i=\operatorname{Supp}(\mathcal F_i)$, on a
$(Y_i)_s=\operatorname{Supp}((\mathcal F_i)_s)$ pour tout $s\in S$
(\textbf{I}, 9.1.13) et en particulier $(Y_i)_\eta=Z_i$ pour tout $i$.
L'hypothèse entraîne que pour $i\neq j$, $Z_i-(Z_i\cap Z_j)$ est ouvert dense
dans $Z_i$~; on déduit donc de (9.5.3) et (9.5.4) que dans un voisinage de
$\eta$, $(Y_i)_s-((Y_i)_s\cap(Y_j)_s)$ est ouvert et dense dans $(Y_i)_s$.
Supposons que l'on ait prouvé que chacun des $(\mathcal F_i)_s$ est sans cycle
premier associé pour $s\in U$. Alors les éléments de
$\operatorname{Ass}((\mathcal F_i)_s)$ sont les points maximaux de $(Y_i)_s$~;
aucun d'eux ne peut donc appartenir à un $(Y_j)_s$ pour $i\neq j$, et la
proposition sera démontrée. On peut donc supposer que $\mathcal F_\eta$ est
\emph{irrédondant}~; en outre, on peut se borner au cas où
$S=\operatorname{Spec}(A)$ est affine~; $X$ est alors réunion d'un nombre fini
d'ouverts affines $V_j$, et si $V_j\cap\operatorname{Supp}(\mathcal F_\eta)
=\varnothing$, on aura aussi $V_j\cap\operatorname{Supp}(\mathcal F_s)=\varnothing$
pour $s$ voisin de $\eta$ (9.5.1), donc on peut se borner au cas où
$X=\operatorname{Spec}(B)$ est aussi affine et où le morphisme $f:X\to S$ est
dominant~; $A$ est donc un anneau noethérien intègre, \emph{sous-anneau} d'une
$A$-algèbre de type fini $B$, et $\mathcal F=\widetilde M$, où $M$ est un
$B$-module de type fini~; par hypothèse, si $K$ est le corps des fractions de
$A$, le $B_{(K)}$-module $M_{(K)}$ est irrédondant. Soit $\mathfrak q$ l'unique

\oldpage[IV-3]{79}

élément de $\operatorname{Ass}(M_{(K)})$, et soit $\mathfrak p$ l'idéal premier
de $B$ image réciproque de $\mathfrak q$ par l'application canonique
$B\to B_{(K)}$. On sait (5.11.1.1) qu'il existe une filtration finie
$(N_h)_{0\leqslant h\leqslant m}$ de $M_{(K)}$ telle que $N_0=M_{(K)}$,
$N_m=0$, et que $N_h/N_{h+1}$ soit isomorphe à un sous-$B_{(K)}$-module
$\neq0$ de $B_{(K)}/\mathfrak q$. Soit $M_h$ l'image réciproque de $N_h$ par
l'application canonique $M\to M_{(K)}$. Il résulte de (9.4.4) que pour $s$
assez voisin de $\eta$, $(M_h)_{(\boldsymbol{k}(s))}$ s'identifie à un
sous-$B_{(\boldsymbol{k}(s))}$-module de $M_{(\boldsymbol{k}(s))}$, et le
quotient
$(M_h)_{(\boldsymbol{k}(s))}/(M_{h+1})_{(\boldsymbol{k}(s))}$ à un
sous-$B_{(\boldsymbol{k}(s))}$-module $\neq0$ de
$B_{(\boldsymbol{k}(s))}/\mathfrak p_{(\boldsymbol{k}(s))}
=(B/\mathfrak p)_{(\boldsymbol{k}(s))}$. Compte tenu de (3.1.7), on voit qu'on
est ramené à prouver, avec les mêmes notations, que si $B$ est \emph{intègre},
alors $B_{(\boldsymbol{k}(s))}$ n'a pas d'idéaux premiers associés immergés
pour $s$ voisin de $\eta$. Or, en remplaçant au besoin $A$ par $A_g$ et $B$ par
$B_g$ (où $g$ est un élément $\neq0$ de $A$), on peut supposer que $B$ contient
une $A$-algèbre de polynômes $C=A[\mathrm T_1,\ldots,\mathrm T_n]$ telle que
$B$ soit un $C$-module de type fini (Bourbaki, \emph{Alg. comm.}, chap. V,
\S~3, n$^\circ$~1, cor. 1 du th. 1). Comme $B_{(K)}$ est un $C_{(K)}$-module
sans torsion, on peut appliquer de nouveau le raisonnement fait plus haut en
remplaçant $B$, $M$ et $\mathfrak q$ par $C$, $B$ et $(0)$ respectivement, et
il suffit donc de voir que pour $s$ voisin de $\eta$, $C_{(\boldsymbol{k}(s))}$
n'a pas d'idéaux premiers associés immergés. Mais cela est évident puisque
$C_{(\boldsymbol{k}(s))}=\boldsymbol{k}(s)[\mathrm T_1,\ldots,\mathrm T_n]$
est un anneau intègre. C.Q.F.D.

\begin{theorem}[9.7.7]
\label{IV.9.7.7-fr}
Soit $f:X\to S$ un morphisme de présentation finie, et soit $E$ l'ensemble des
$s\in S$ pour lesquels $X_s$ a l'une des propriétés suivantes~: être~:
\begin{enumerate}[label=(\roman*)]
  \item géométriquement irréductible~;
  \item géométriquement connexe~;
  \item géométriquement réduit~;
  \item géométriquement intègre.
\end{enumerate}
Alors $E$ est localement constructible dans $S$.
\end{theorem}

Pour voir que les propriétés envisagées sont constructibles, remarquons d'abord
qu'elles vérifient trivialement la condition (9.2.1, (i)), en vertu de
(4.5.6, (i)) et (4.6.5, (i)). Il reste donc à vérifier (9.2.1, (ii)), donc on
peut supposer $S=\operatorname{Spec}(A)$ affine noethérien et intègre de point
générique $\eta$. En vertu de (4.6.8), il existe une extension finie $K'$ de
$K=\boldsymbol{k}(\eta)$ telle que $(X_\eta)_{(K')}$ soit tel que ses
composantes irréductibles (resp. connexes) soient géométriquement irréductibles
(resp. géométriquement connexes) et que $((X_\eta)_{(K')})_{\mathrm{red}}$ soit
géométriquement réduit. Comme il existe une base de $K'$ sur $K$ formée
d'éléments entiers sur $A$, l'anneau intègre $A'$ engendré par ces éléments est
fini sur $A$ et a $K'$ pour corps des fractions. Posons
$S'=\operatorname{Spec}(A')$~; le morphisme $g:S'\to S$ est fini et dominant,
donc surjectif (\textbf{II}, 6.1.10). Posons $X'=X_{(S')}$, de sorte que si
$\eta'$ est le point générique de $S'$, on a
$X'_{\eta'}=(X_\eta)_{(K')}$~; l'ensemble $E'$ des $s'\in S'$ tels que
$X'_{s'}$ ait l'une des propriétés (i), (ii), (iii) ou (iv) est égal à
$g^{-1}(E)$ (9.2.2, (iv)) ($E$ correspondant bien entendu à la même propriété)~;
$g$ étant surjectif, on a $E=g(E')$ et $S-E=g(S'-E')$~; en outre, $g$ est fermé
et $g^{-1}(\eta)=\{\eta'\}$ puisque $A'$ est intègre et fini sur $A$ (Bourbaki,
\emph{Alg. comm.}, chap. V, \S~2, n$^\circ$~1, cor. 1 de la prop. 1), donc
l'image par $g$ de tout voisinage de $\eta'$ est un voisinage de $\eta$. Le
théorème sera donc démontré si l'on prouve que $E'$ ou $S'-E'$ est un voisinage
de $\eta'$. Autrement dit, on peut désormais supposer que les composantes
irréductibles (resp. connexes) de $X_\eta$ sont géométriquement

\oldpage[IV-3]{80}

irréductibles (resp. géométriquement connexes) et que $(X_\eta)_{\mathrm{red}}$
est géométriquement réduit.

Supposons d'abord que $\eta\in S-E$ pour une des propriétés (i) à (iv). Si
$X_\eta$ n'est pas géométriquement irréductible (resp. géométriquement connexe),
il n'est pas irréductible (resp. connexe) en vertu de l'hypothèse précédente,
donc il en est de même de $X_s$ pour $s$ dans un voisinage de $\eta$ (9.7.1), et
\emph{a fortiori} dans ce voisinage $X_s$ n'est pas géométriquement irréductible
(resp. géométriquement connexe). D'autre part, si $X_\eta$ n'est pas
géométriquement réduit, il n'est pas réduit (sans quoi il serait égal à
$(X_\eta)_{\mathrm{red}}$ qui est géométriquement réduit par hypothèse)~; donc,
dans un voisinage de $\eta$, $X_s$ n'est pas réduit (9.7.2) et \emph{a fortiori}
n'est pas géométriquement réduit. Enfin, si $X_\eta$ n'est pas géométriquement
intègre, ou bien il n'est pas réduit, et alors on vient de voir que $X_s$ n'est
pas réduit (ni \emph{a fortiori} intègre) dans un voisinage de $\eta$~; ou bien
$X_\eta$ est réduit (donc géométriquement réduit par hypothèse), et alors il
n'est pas géométriquement irréductible, et on a vu ci-dessus qu'il en est alors
de même de $X_s$ pour $s$ voisin de $\eta$~; \emph{a fortiori} $X_s$ n'est pas
géométriquement intègre pour ces valeurs de $s$.

Nous allons donc désormais supposer que $\eta\in E$ et examiner séparément
chacune des propriétés considérées.

1$^\circ$ Supposons que $X_\eta$ soit géométriquement intègre. Soit $L$ le corps
des fonctions rationnelles sur $X_\eta$~; l'hypothèse sur $X_\eta$ implique que
$L$ est extension séparable de $K$ (4.6.3), donc extension séparable finie d'une
extension pure $K(\mathrm T_1,\ldots,\mathrm T_n)$ ($\mathrm T_i$
indéterminées)~; il y a par suite un élément $x\in L$, entier sur l'anneau
$K[\mathrm T_1,\ldots,\mathrm T_n]$, et tel que
$L=K(\mathrm T_1,\ldots,\mathrm T_n)(x)$~; soit
$G\in K[\mathrm T_1,\ldots,\mathrm T_n,\mathrm T_{n+1}]$ son polynôme minimal.
Il existe un élément $g\neq0$ de $A$ tel que tous les coefficients $\neq0$ de
$G$ (qui appartiennent à $K$) soient dans l'anneau $A_g$~; remplaçant $A$ par
$A_g$ (ce qui revient à remplacer $S$ par un voisinage de $\eta$), on peut donc
supposer que $G$ a ses coefficients dans $A$~; désignant par $F$ le polynôme
$G$ considéré comme élément de
$A[\mathrm T_1,\ldots,\mathrm T_n,\mathrm T_{n+1}]$, on a donc
$G=F_{(\boldsymbol{k}(\eta))}$. Posons
$Y=\operatorname{Spec}(A[\mathrm T_1,\ldots,\mathrm T_{n+1}]/(F))$, de sorte
que
$Y_\eta=\operatorname{Spec}(K[\mathrm T_1,\ldots,\mathrm T_{n+1}]/(G))$~;
$Y_\eta$ est un schéma intègre ayant $L$ pour corps des fonctions rationnelles.
Comme $X_\eta$ et $Y_\eta$ sont noethériens, il existe un ouvert non vide
$V\subset Y_\eta$ et une immersion ouverte $v:V\to X_\eta$ (nécessairement
dominante) (\textbf{I}, 6.5.1, (ii) et 6.5.4, (ii)). Soit $W$ un ouvert de $Y$
tel que $W\cap Y_\eta=V$~; appliquant (8.8.2, (i)) et (8.10.5, (iii)) suivant
la méthode de (8.1.2, \emph{a)}), on voit qu'en remplaçant au besoin $S$ par un
voisinage de $\eta$, on peut supposer que $v=w_\eta$, où $w:W\to X$ est une
immersion ouverte.

Cela étant, on a vu (4.6.3) que le critère pour qu'un préschéma algébrique
intègre soit géométriquement intègre ne dépend que de son corps de fonctions
rationnelles~; comme $X_\eta$ est géométriquement intègre par hypothèse, il en
est de même de $Y_\eta$, et la définition de $G$ entraîne donc que ce polynôme
est géométriquement irréductible (9.7.4). Appliquant (9.7.5), on voit qu'il y a
un voisinage $U$ de $\eta$ dans $S$ tel que $F_{(\boldsymbol{k}(s))}$ soit
géométriquement irréductible pour tout $s\in U$, donc $Y_s$ est géométriquement
intègre pour $s\in U$ (9.7.4)~; en outre on peut supposer que pour $s\in U$,
$W_s$ n'est pas vide (9.5.1), et par suite est géométriquement intègre (4.6.3)~;
enfin, on peut supposer aussi que

\oldpage[IV-3]{81}

pour $s\in U$, $w_s:W_s\to X_s$ est une immersion ouverte (9.6.1, (x)), et que
son image dans $X_s$ est partout dense (9.5.3). Autrement dit, pour $s\in U$,
il y a dans $X_s$ un ouvert partout dense $W'_s$ qui est géométriquement
intègre~; le critère (4.5.9, \emph{c)}) montre donc déjà que $X_s$ est
géométriquement irréductible pour $s\in U$. Enfin, comme $X_\eta$ est réduit et
par suite n'a pas de cycle premier associé immergé (3.2.1), on peut aussi
supposer que pour $s\in U$, $X_s$ n'a pas de cycle premier associé immergé
(9.7.6)~; soit alors (pour un $s\in U$ fixé) $\Omega$ une extension
algébriquement close de $\boldsymbol{k}(s)$, et soit
$p:(X_s)_{(\Omega)}\to X_s$ la projection canonique~; $p^{-1}(W'_s)$ est un
ouvert dense dans $(X_s)_{(\Omega)}$ (2.3.10) et est intègre par hypothèse~; en
outre $(X_s)_{(\Omega)}$ n'a pas de cycle premier associé immergé (4.2.7), donc
on conclut de (3.2.1) que $(X_s)_{(\Omega)}$ est réduit~; cela achève de prouver
que $X_s$ est géométriquement intègre (4.6.1).

2$^\circ$ Supposons que $X_\eta$ soit géométriquement irréductible~; comme
$X_{\mathrm{red}}$ est aussi de type fini sur $S$ (1.5.4, (vi)) on peut, en
tenant compte de (\textbf{I}, 5.1.8), remplacer $X$ par $X_{\mathrm{red}}$~;
alors $X_\eta$ est aussi intègre, et comme par hypothèse $X_\eta$ est
géométriquement réduit, il est géométriquement intègre. On est alors dans les
conditions de 1$^\circ$ et on en conclut (revenant aux hypothèses initiales) que
$X_s$ est géométriquement irréductible pour $s$ voisin de $\eta$.

3$^\circ$ Supposons $X_\eta$ géométriquement connexe, et soient $Z_i$
($1\leqslant i\leqslant n$) les composantes irréductibles de $X_\eta$~; il
existe (en vertu de (5.10.8.1) appliqué à $\mathfrak T=\{\varnothing\}$) une
application \emph{surjective} $j\mapsto\nu(j)$ d'un intervalle $[1,m]$ de
$\mathbf N$ sur $[1,n]$ telle que
$Z_{\nu(j)}\cap Z_{\nu(j+1)}\neq\varnothing$ pour $1\leqslant j\leqslant m$.
Pour tout $i$, soit $X_i$ l'adhérence de $Z_i$ dans $X$, et soit $Y$ la réunion
des $X_i$~; comme $Y_\eta=X_\eta$ par définition, on peut supposer, en vertu de
(9.5.1), que $Y_s=X_s$ pour tout $s\in S$, donc que $X_s$ est réunion des
$(X_i)_s$. Or, en considérant les sous-préschémas fermés réduits de $X$ ayant
les $X_i$ pour espaces sous-jacents, on voit par 2$^\circ$ qu'il existe un
voisinage $U$ de $\eta$ dans $S$ tel que pour $s\in U$ les $(X_i)_s$ soient
\emph{géométriquement irréductibles} (puisque les $(X_i)_\eta$ peuvent être
supposés géométriquement irréductibles, comme on l'a vu au début). En outre, on
peut aussi supposer que pour $s\in U$, on a
$(X_{\nu(j)})_s\cap(X_{\nu(j+1)})_s\neq\varnothing$ (9.5.1) pour
$1\leqslant j\leqslant m$~; on en conclut aussitôt que $X_s$ est connexe, donc
(4.5.13.1) géométriquement connexe pour $s\in U$.

4$^\circ$ Supposons $X_\eta$ géométriquement réduit~; soient $Z_i$ les
composantes irréductibles de $X_\eta$, $W_i$ l'intérieur de $Z_i$ dans
$X_\eta$~; il y a pour chaque $i$ un ouvert $V_i$ de $X$ tel que
$W_i=V_i\cap X_\eta$ pour tout $i$~; comme les $W_i$ sont ouverts et deux à deux
disjoints et que leur réunion est dense dans $X_\eta$, on peut (9.5.1, 9.5.3 et
9.5.4) supposer que pour $s$ voisin de $\eta$, les $(V_i)_s$ sont des ouverts
deux à deux disjoints dans $X_s$ et que leur réunion est dense dans $X_s$. En
outre, comme les $W_i$ sont géométriquement réduits et ont été supposés au
début géométriquement irréductibles, il résulte du 1$^\circ$ que pour $s$ voisin
de $\eta$, les $(V_i)_s$ sont géométriquement intègres, et \emph{a fortiori}
réduits. D'autre part, on tire de (9.7.6) que pour $s$ voisin de $\eta$, $X_s$
n'a pas de cycle premier associé immergé, puisqu'il en est ainsi de $X_\eta$
qui est réduit (3.2.1)~; on conclut donc de (3.2.1) que $X_s$ est réduit, et de
(4.6.1) qu'il est géométriquement réduit.

\oldpage[IV-3]{82}

Les parties de l'énoncé (9.7.7) concernant les propriétés (i) et (ii) se
généralisent comme suit~:

\begin{proposition}[9.7.8]
\label{IV.9.7.8-fr}
Soient $S$ un préschéma irréductible de point générique $\eta$, et $f:X\to S$
un morphisme de présentation finie. Soit $n$ (resp. $n'$) le nombre géométrique
de composantes irréductibles (resp. connexes) de $X_\eta$ (4.5.2). Alors il
existe un voisinage $U$ de $\eta$ dans $S$ tel que pour tout $s\in U$ le nombre
géométrique de composantes irréductibles (resp. connexes) de $X_s$ soit égal à
$n$ (resp. $n'$).
\end{proposition}

Compte tenu de ce que le nombre géométrique de composantes irréductibles (resp.
connexes) d'un préschéma algébrique est invariant par extension du corps de
base (4.5.6), on voit par la méthode de (9.2.3) que l'on peut se ramener au cas
où $S$ est affine, noethérien et intègre. En outre, en raisonnant comme au début
de la démonstration de (9.7.7), on voit qu'on peut supposer que les composantes
irréductibles (resp. connexes) de $X_\eta$ soient \emph{géométriquement}
irréductibles (resp. \emph{géométriquement connexes}). On sait déjà (9.7.1) que
pour tout $s$ voisin de $\eta$, $X_s$ a au moins $n$ composantes irréductibles
et $n'$ composantes connexes. D'autre part, si $Z_i$ (resp. $Z'_j$) sont les
composantes irréductibles (resp. connexes) de $X_\eta$, $X_i$ (resp. $X'_j$) le
sous-préschéma fermé réduit de $X$ ayant pour espace sous-jacent l'adhérence de
$Z_i$ (resp. $Z'_j$) dans $X$, il résulte de (9.5.1) que dans un voisinage de
$\eta$, les $(X_i)_s$ (resp. $(X'_j)_s$) forment un recouvrement de $X_s$ et que
les $(X'_j)_s$ sont deux à deux disjoints~; comme les $(X_i)_s$ (resp.
$(X'_j)_s$) sont géométriquement irréductibles (resp. géométriquement connexes)
en vertu de (9.7.7) pour $s$ voisin de $\eta$, on voit que $X_s$ a au plus $n$
composantes irréductibles et au plus $n'$ composantes connexes, d'où la
proposition.

\begin{corollary}[9.7.9]
\label{IV.9.7.9-fr}
Soit $f:X\to S$ un morphisme de présentation finie. Pour tout $s\in S$, soit
$n(s)$ (resp. $n'(s)$) le nombre géométrique de composantes irréductibles
(resp. connexes) de $f^{-1}(s)$ (4.5.2). Alors la fonction $s\mapsto n(s)$
(resp. $s\mapsto n'(s)$) est localement constructible dans $S$.
\end{corollary}

Il s'agit de montrer que la propriété «~$X$ est un préschéma algébrique sur un
corps $k$ et le nombre géométrique des composantes irréductibles (resp.
connexes) de $X$ est égal à $n$ (resp. $n'$)~» est constructible. Il résulte de
(4.5.6) que cette propriété vérifie la condition (9.2.1, (i)), et on est donc
ramené au cas où $S$ est affine, noethérien et intègre de point générique
$\eta$~; la conclusion résulte alors de (9.7.8).

\begin{corollary}[9.7.10]
\label{IV.9.7.10-fr}
Soit $X$ un préschéma localement noethérien tel que, si $X'$ est le normalisé de
$X_{\mathrm{red}}$, le morphisme canonique $f:X'\to X$ soit fini. Alors
l'ensemble des points $x\in X$ tels que $X$ soit géométriquement unibranche au
point $x$ est localement constructible dans $X$.
\end{corollary}

En effet, cet ensemble est par définition l'ensemble des points $x\in X$ tels
que le nombre de points géométriques de $f^{-1}(x)$ soit égal à 1. Mais comme
$f$ est fini, ce nombre est aussi le nombre géométrique des composantes
irréductibles de l'espace discret $f^{-1}(x)$ (compte tenu de la définition du
normalisé (\textbf{II}, 6.3.8) et de (4.5.11))~; la conclusion résulte donc de
(9.7.9).

\begin{remark}[9.7.11]
\label{IV.9.7.11-fr}
Soit $Z$ une partie localement constructible de $X$ telle que, pour tout
$s\in S$, $Z\cap X_s$ soit ouvert dans $X_s$, et désignons par $Z_s$ le
sous-préschéma de $X_s$ induit

\oldpage[IV-3]{83}

sur l'ouvert $Z\cap X_s$. Alors, dans le théorème (9.7.7), on peut remplacer
$X_s$ par $Z_s$ sans changer la conclusion~: on le voit en répétant le
raisonnement fait dans (9.6.7, (ii)).
\end{remark}

\begin{proposition}[9.7.12]
\label{IV.9.7.12-fr}
Soient $f:X\to S$ un morphisme de présentation finie, et $g:S\to X$ une
$S$-section de $X$ (\textbf{I}, 2.5.5). Pour tout $s\in S$, soit $X_s^0$ la
composante connexe de $f^{-1}(s)$ qui contient $g(s)$~; alors, la réunion $X^0$
des $X_s^0$ pour $s\in S$ est une partie localement constructible de $X$.
\end{proposition}

Montrons d'abord qu'on peut se ramener au cas où $S$ est affine et
\emph{noethérien}. On peut toujours supposer $S=\operatorname{Spec}(A)$ affine~;
avec les notations de (9.2.3), on a $f=(f_0)_{(S)}$, où $f_0:X_0\to S_0$ est un
morphisme de type fini, et on peut en outre supposer qu'il existe une
$S_0$-section $g_0:S_0\to X_0$ telle que $g=(g_0)_{(S)}$ (8.9.1). Remarquons
maintenant que si $p$ est le morphisme $S\to S_0$, alors, pour tout $s_0\in S_0$,
la composante connexe $(X_0)^0_{s_0}$ de $f_0^{-1}(s_0)$ qui contient $g_0(s_0)$
est \emph{géométriquement connexe} (4.5.13), et par suite, si $s_0=p(s)$, on a
$X_s^0=q^{-1}((X_0)^0_{s_0})$ où $q:X_s\to(X_0)_{s_0}$ est la projection
canonique ((4.5.8) et (4.4.1))~; notre assertion résulte donc de (1.8.2).

Utilisons alors le critère de constructibilité $(\mathbf 0_{\mathrm{III}},9.2.3)$~:
soient $x$ un point de $X$, $Z$ le sous-préschéma réduit de $X$ ayant
$\overline{\{x\}}$ comme espace sous-jacent, $Y$ le sous-préschéma réduit de $S$
ayant $\overline{f(Z)}$ comme espace sous-jacent~; comme la restriction de $f$
à $Z$ se factorise en $Z\to Y\to S$ (\textbf{I}, 5.2.2), on peut remplacer $S$
par $Y$, autrement dit supposer que $f(x)=\eta$, point générique du préschéma
intègre $S$. Par hypothèse, $X_\eta$ est somme de deux sous-préschémas
$X_\eta^0$, $X_\eta^1$, induits sur des ouverts complémentaires de $X_\eta$. En
vertu de (9.5.4) et (9.5.1), on peut donc, en remplaçant au besoin $S$ par un
voisinage de $\eta$, supposer que $X$ est réunion de deux ouverts disjoints
$X^0$, $X^1$ tels que $(X^0)_\eta=X_\eta^0$ et $(X^1)_\eta=X_\eta^1$. Comme
$g:S\to X$ est continu et injectif, $S$ est réunion des deux ouverts disjoints
$g^{-1}(X^0)$ et $g^{-1}(X^1)$~; mais comme $S$ est irréductible, et
\emph{a fortiori} connexe, un de ces deux ouverts est vide, et comme
$g(\eta)\in X_\eta^0$ par définition, on a $g^{-1}(X^1)=\varnothing$. En
d'autres termes, $g$ est une $S$-section de $X^0$~; d'autre part, comme
$(X^0)_\eta$ est géométriquement connexe (4.5.13), il en est de même de
$(X^0)_s$ pour tout $s$ voisin de $\eta$ (9.7.7)~; comme
$g(s)\in(X^0)_s$, on a bien $(X^0)_s=X_s^0$. C.Q.F.D.

\subsection{Décomposition primaire au voisinage d'une fibre générique}
\label{subsection:IV.9.8-fr}

\begin{proposition}[9.8.1]
\label{IV.9.8.1-fr}
Soient $f:X\to S$ un morphisme de présentation finie, $\mathcal F$ un
$\mathcal O_X$-Module quasi-cohérent de présentation finie. Alors l'ensemble
$E$ des $s\in S$ tels que $\mathcal F_s$ soit un $\mathcal O_{X_s}$-Module sans
cycle premier associé immergé est localement constructible.
\end{proposition}

On sait que pour les Modules quasi-cohérents sur les préschémas algébriques, la
propriété d'être sans cycle premier associé immergé est invariante par
changement du corps de base (4.2.7)~; on peut donc se borner au cas où $S$ est
affine, noethérien et intègre de point générique $\eta$, et à prouver que dans
ce cas $E$ ou $S-E$ est un voisinage de $\eta$. On a vu dans (9.7.6) que si
$\eta\in E$, $E$ est un voisinage de $\eta$~; reste donc à considérer le cas où
$\mathcal F_\eta$ a des cycles premiers associés immergés. On peut se borner au
cas où $X$ est affine, et où il y a un sous-$\mathcal O_{X_\eta}$-Module
cohérent $\mathcal H$ de $\mathcal F_\eta$ dont le support $Z'$ est non vide et
\emph{rare} par rapport au support $Z$ de $\mathcal F$ (3.1.3)~; alors, en vertu
de (8.5.2, (i) et (ii))

\oldpage[IV-3]{84}

et (8.5.8), appliqués suivant la méthode de (8.1.2, \emph{a)}), il existe un
sous-$\mathcal O_X$-Module cohérent $\mathcal G$ de $\mathcal F$ tel que
$\mathcal G_\eta=\mathcal H$~; si $Y=\operatorname{Supp}(\mathcal F)$,
$Y'=\operatorname{Supp}(\mathcal G)$, on a par suite
$Y_s=\operatorname{Supp}(\mathcal F_s)$, $Y'_s=\operatorname{Supp}(\mathcal G_s)$
pour tout $s\in S$ (\textbf{I}, 9.1.13) et en particulier $Z=Y_\eta$,
$Z'=Y'_\eta$~; comme $Z-Z'$ est dense dans $Z$ et $Z'\neq\varnothing$, il y a
un voisinage $U$ de $\eta$ tel que pour $s\in U$, $Y_s-Y'_s$ soit dense dans
$Y_s$, et $Y'_s\neq\varnothing$ (9.5.1 et 9.5.3)~; en considérant un point
générique d'une composante irréductible de $Y'_s$ et un voisinage assez petit de
ce point dans $X_s$, on déduit aussitôt de (3.1.3) que $\mathcal F_s$ admet des
cycles premiers associés immergés.

\phantomsection\label{IV.9.8.2-fr}%
\textbf{(9.8.2)} Soient $S$ un préschéma noethérien \emph{intègre}, de point
générique $\eta$, $f:X\to S$ un morphisme de type fini, $\mathcal F$ un
$\mathcal O_X$-Module cohérent. Considérons une décomposition irredondante
réduite $(\mathcal G_i)_{i\in I}$ de $\mathcal F_\eta$ (3.2.6)~; les
$\mathcal G_i$ sont donc des quotients de $\mathcal F_\eta$ et il y a un
homomorphisme injectif
$h:\mathcal F_\eta\to\bigoplus_i\mathcal G_i$~; en outre
$\operatorname{Ass}(\mathcal G_i)$ est réduit à un point $x_i$. Soit $Z_i$
l'adhérence de $\{x_i\}$ dans $X$, de sorte que
$(Z_i)_\eta=\operatorname{Supp}(\mathcal G_i)$. Désignons par $\mathcal J_i$
l'Idéal cohérent de $\mathcal O_X$ définissant le sous-préschéma fermé réduit de
$X$ d'espace sous-jacent $Z_i$ (sous-préschéma encore noté $Z_i$). En vertu de
(8.5.2) et (8.5.8), appliqués suivant la méthode de (8.1.2, \emph{a)}), on peut
(en restreignant au besoin $S$ à un voisinage de $\eta$) supposer qu'il existe
des quotients $\mathcal F_i$ de $\mathcal F$ ($i\in I$) tels que
$\mathcal G_i=(\mathcal F_i)_\eta$ pour tout $i$, et un homomorphisme
$g:\mathcal F\to\bigoplus_i\mathcal F_i$ tel que $h=g_\eta$. En outre
(\textbf{I}, 9.3.5), il existe un entier $m$ tel que
$(\mathcal J_i^m\mathcal F_i)_\eta=(\mathcal J_i)_\eta^m\mathcal G_i=0$ et en
restreignant de nouveau $S$, on peut donc supposer que l'on a aussi
$\mathcal J_i^m\mathcal F_i=0$ (8.5.2.5), donc que le support de $\mathcal F_i$
est contenu dans $Z_i$~; mais comme il est fermé et contient $x_i$, il est
nécessairement égal à $Z_i$.

\begin{proposition}[9.8.3]
\label{IV.9.8.3-fr}
Sous les conditions de (9.8.2), pour tout $s\in S$, et tout $i\in I$, désignons
par $x_{is\alpha}$ ($\alpha\in J_{s,i}$) les points maximaux de $(Z_i)_s$. Il
existe un voisinage $U$ de $\eta$ dans $S$ tel que, pour tout $s\in U$, les
$x_{is\alpha}$ (pour $i\in I$ et, pour chaque $i$, $\alpha\in J_{s,i}$) soient
deux à deux distincts et que $\operatorname{Ass}(\mathcal F_s)$ soit l'ensemble
des $x_{is\alpha}$ (autrement dit, les cycles premiers associés à $\mathcal F_s$
sont les composantes irréductibles des $(Z_i)_s$). En outre, on peut prendre
$U$ tel que, pour que l'adhérence de $\{x_{is\alpha}\}$ dans $X_s$ soit un cycle
premier associé maximal de $\mathcal F_s$, il faut et il suffit que $(Z_i)_\eta$
(adhérence de $\{x_i\}$ dans $X_\eta$) soit un cycle premier associé maximal de
$\mathcal F_\eta$.
\end{proposition}

Il résulte de (3.1.3, \emph{c')}) que pour chaque $i$, il existe un ouvert $W_i$
dans $X_\eta$ tel que $W_i\cap(Z_i)_\eta$ soit non vide, et un
$\mathcal O_{X_\eta}$-Module cohérent $\mathcal H_i$, de support
$W_i\cap(Z_i)_\eta$, tels qu'il y ait un homomorphisme injectif
$v_i:\mathcal H_i\to\mathcal F_\eta|W_i$. Soit $V_i$ un ouvert de $X$ tel que
$V_i\cap X_\eta=W_i$~; appliquant, comme dans (9.8.2), les résultats de (8.5.2)
et (8.5.8), on peut (en restreignant $S$) supposer qu'il existe un Module
cohérent $\mathcal F'_i$, de support $V_i\cap Z_i$ et un homomorphisme
$u_i:\mathcal F'_i\to\mathcal F|V_i$ tels que
$\mathcal H_i=(\mathcal F'_i)_\eta$ et $v_i=(u_i)_\eta$. Nous allons démontrer
qu'il y a un voisinage $U$ de $\eta$ tel que pour $s\in U$, les propriétés
suivantes soient vérifiées~:
\begin{enumerate}[label=\emph{(\roman*)}]
  \item Les $(\mathcal F_i)_s$ sont sans cycle premier associé immergé.
  \item L'homomorphisme
  $g_s:\mathcal F_s\to\bigoplus_i(\mathcal F_i)_s$ est injectif.
  \item $(Z_i)_s\cap(V_i)_s$ est dense dans $(Z_i)_s$,
  $(\mathcal F'_i)_s$ a pour support $(Z_i)_s\cap(V_i)_s$ et
  $(u_i)_s:(\mathcal F'_i)_s\to\mathcal F_s|(V_i)_s$ est injectif.
\end{enumerate}

\oldpage[IV-3]{85}

\begin{enumerate}[label=\emph{(\roman*)}]
  \item[(iv)] Pour $i\neq j$, toute composante irréductible de $(Z_i)_s$ est
  distincte de toute composante irréductible de $(Z_j)_s$.
\end{enumerate}
Or, (i) a déjà été vu (9.7.6)~; (ii) est un cas particulier de (9.4.5)~; (iii)
résulte de même de (9.5.3) et (9.4.5). Enfin, si $i\neq j$,
$(Z_i)_\eta\cap(Z_j)_\eta=(Z_i\cap Z_j)_\eta$ est rare dans $(Z_i)_\eta$ ou
dans $(Z_j)_\eta$~; supposons par exemple que $(Z_i\cap Z_j)_\eta$ soit rare
dans $(Z_j)_\eta$~; alors il résulte de (9.5.3) et (9.5.4) que pour $s$ voisin
de $\eta$, $(Z_i\cap Z_j)_s=(Z_i)_s\cap(Z_j)_s$ est rare dans $(Z_j)_s$, ce qui
montre qu'aucune composante irréductible de $(Z_j)_s$ ne peut être contenue
dans une composante irréductible de $(Z_i)_s$, ni \emph{a fortiori} lui être
égale.

Cela étant, il résulte de (ii) et de (3.1.7) que pour $s\in U$, on a
\[
\operatorname{Ass}(\mathcal F_s)\subset
\bigcup_i\operatorname{Ass}((\mathcal F_i)_s),
\]
de (i) que $\operatorname{Ass}((\mathcal F_i)_s)$ est l'ensemble des
$x_{is\alpha}$ ($\alpha\in J_{s,i}$). D'autre part, en vertu de (iii) et du
critère (3.1.3, \emph{c')}), chacun des $x_{is\alpha}$ ($\alpha\in J_{s,i}$,
$i\in I$) appartient à $\operatorname{Ass}(\mathcal F_s)$. Enfin (iv) signifie
que les $x_{is\alpha}$ sont deux à deux distincts.

Reste à prouver la dernière assertion de l'énoncé. Or, il résulte de (i) que
pour $s$ et pour un $i$ donnés, aucun des ensembles $\{x_{is\alpha}\}$ ne peut
être contenu dans un autre pour $\alpha\in J_{s,i}$. D'autre part, si
$(Z_j)_\eta\subset(Z_i)_\eta$, on peut supposer $U$ pris tel que
$(Z_j)_s\subset(Z_i)_s$ pour $s\in U$ (9.5.1), donc chaque $x_{js\beta}$
appartient à l'adhérence d'un $x_{is\alpha}$~; au contraire, si
$(Z_j)_\eta\cap(Z_i)_\eta$ est rare dans $(Z_j)_\eta$, on a vu en prouvant (iv)
que $(Z_j)_s\cap(Z_i)_s$ est rare dans $(Z_j)_s$, donc aucun des $x_{js\beta}$
n'est adhérent à un $x_{is\alpha}$. En particulier, si $(Z_j)_\eta$ est
maximal, ce qui équivaut à dire que $(Z_j)_\eta\cap(Z_i)_\eta$ est rare dans
$(Z_j)_\eta$ pour \emph{tout} $i\neq j$, on en conclut que
$(Z_j)_s\cap(Z_i)_s$ est rare dans $(Z_j)_s$ pour tout $i\neq j$, donc que tout
$x_{js\alpha}$ ($\alpha\in J_{s,j}$) est \emph{maximal}. C.Q.F.D.

\begin{corollary}[9.8.4]
\label{IV.9.8.4-fr}
Les notations et hypothèses étant celles de (9.8.2), il existe un voisinage $U$
de $\eta$ dans $S$ tel que, pour tout $s\in U$, chacun des $(\mathcal F_i)_s$
soit sans cycle premier associé immergé~; en outre, si
$(\mathcal F_{is\alpha})_{\alpha\in J_{s,i}}$ est l'unique décomposition
irredondante réduite de $(\mathcal F_i)_s$, alors, pour tout $s\in U$, la
famille $(\mathcal F_{is\alpha})_{i\in I,\,\alpha\in J_{s,i}}$ est une
décomposition irredondante réduite de $\mathcal F_s$.
\end{corollary}

La première assertion résulte de (9.8.1) et de la définition des
$\mathcal F_i$. D'autre part, on a vu dans (9.8.3) que l'homomorphisme
$\mathcal F_s\to\bigoplus_i(\mathcal F_i)_s$ est injectif et par définition il
en est de même de chacun des homomorphismes
$(\mathcal F_i)_s\to\bigoplus_\alpha\mathcal F_{is\alpha}$, donc
l'homomorphisme $\mathcal F_s\to\bigoplus_{i,\alpha}\mathcal F_{is\alpha}$ est
injectif. Comme on peut supposer (9.4.5) que pour $s\in U$ chacun des
$(\mathcal F_i)_s$ est un quotient de $\mathcal F_s$, les
$\mathcal F_{is\alpha}$ sont des quotients de $\mathcal F_s$, et il reste à
vérifier (3.2.5) que les $x_{is\alpha}$ sont deux à deux distincts et
appartiennent à $\operatorname{Ass}(\mathcal F_s)$, ce qui a été prouvé dans
(9.8.3).

\begin{proposition}[9.8.5]
\label{IV.9.8.5-fr}
Soient $f:X\to S$ un morphisme de présentation finie, $\mathcal F$ un
$\mathcal O_X$-Module quasi-cohérent de présentation finie. Pour tout $s\in S$,
soit $E(s)$ (resp. $E'(s)$) l'ensemble fini (partie de
$\mathbf Z\cup\{-\infty\}$) des dimensions des cycles premiers associés à
$\mathcal F_s$ (resp. des cycles premiers maximaux associés à $\mathcal F_s$,
c'est-à-dire les composantes irréductibles de $\operatorname{Supp}(\mathcal F_s)$).
Alors les fonctions $s\mapsto E(s)$ et $s\mapsto E'(s)$ sont localement
constructibles dans $S$.
\end{proposition}

\oldpage[IV-3]{86}

Il résulte de (4.2.7) et (4.2.8) que la condition (9.2.1, (i)) est satisfaite
pour les propriétés dont nous voulons démontrer qu'elles sont constructibles.
On est donc ramené au cas où $S$ est affine, noethérien et intègre, et à prouver
que $E$ et $E'$ sont constantes au voisinage du point générique $\eta$ de $S$~;
mais cela résulte de (9.8.3) et (9.5.5).

\begin{proposition}[9.8.6]
\label{IV.9.8.6-fr}
Avec les hypothèses et notations de (9.8.2), soit $i\in I$ tel que $(Z_i)_\eta$
soit maximal (autrement dit, une composante irréductible de
$\operatorname{Supp}(\mathcal F_\eta)$). Alors, il existe un voisinage $U$ de
$\eta$ dans $S$ tel que pour tout $s\in U$ et tout $\alpha\in J_{s,i}$, la
longueur géométrique de $\mathcal F_s$ en $x_{is\alpha}$ (relative à
$\boldsymbol{k}(s)$) (4.7.5) est égale à la longueur géométrique de
$\mathcal F_\eta$ en $x_i$ (relative à $\boldsymbol{k}(\eta)$).
\end{proposition}

On peut évidemment se borner au cas où $S=\operatorname{Spec}(A)$ est affine~;
montrons d'abord qu'on peut se ramener au cas où le sous-préschéma $(Z_i)_\eta$
de $X_\eta$, qui est réduit, est \emph{géométriquement intègre}. Il y a en effet
une extension finie $K'$ de $K=\boldsymbol{k}(\eta)$ telle que
$(((Z_i)_\eta)_{(K')})_{\mathrm{red}}$ soit géométriquement réduit et que les
composantes irréductibles de $((Z_i)_\eta)_{(K')}$ soient géométriquement
irréductibles (4.6.8). Procédons comme dans la démonstration de (9.7.7) en
considérant une sous-$A$-algèbre $A'$ de $K'$ ayant $K'$ pour corps des
fractions et finie sur $A$. On pose $S'=\operatorname{Spec}(A')$ et on considère
le morphisme fini surjectif $g:S'\to S$~; soient ensuite $X'=X_{(S')}$ et
$\mathcal F'=\mathcal F\otimes_{\mathcal O_S}\mathcal O_{S'}$, et soit $\eta'$
le point générique de $S'$. Pour tout $s'\in S'$, soit $s=g(s')$~; si $T$ est
une composante irréductible de $\operatorname{Supp}(\mathcal F_s)$, les
composantes irréductibles $T'_j$ de $T_{(\boldsymbol{k}(s'))}$ sont des
composantes irréductibles de $\operatorname{Supp}(\mathcal F'_{s'})$ et dominent
$T$ (4.2.7) et les multiplicités radicielles de $T$ par rapport à $\mathcal F_s$
et de chaque $T'_j$ par rapport à $\mathcal F'_{s'}$ sont les mêmes (4.7.9). Le
raisonnement de la première partie de (9.7.7) montre donc qu'on peut se borner à
démontrer la proposition pour $X'$ et $\mathcal F'$~; et en vertu du choix de
$K'$, les sous-préschémas réduits ayant pour espaces sous-jacents les composantes
irréductibles de $((Z_i)_\eta)_{(K')}$ sont géométriquement intègres (4.6.1).

Supposons donc désormais $(Z_i)_\eta$ géométriquement intègre~; alors (9.7.7),
il en est de même de $(Z_i)_s$ pour $s$ voisin de $\eta$~; la définition (4.7.5)
montre qu'il suffira donc de prouver que la \emph{longueur} du
$\mathcal O_{X_\eta,x_i}$-module $(\mathcal F_\eta)_{x_i}$ est égale à celle du
$\mathcal O_{X_s,x_{is}}$-module $(\mathcal F_s)_{x_{is}}$ (on a ici supprimé
l'indice $\alpha$, inutile par hypothèse). La question étant évidemment locale
sur $X$, on peut supposer (en se restreignant à un voisinage de $x_i$) que
$\mathcal F=\mathcal F_i$, donc que $\mathcal F_\eta$ est \emph{irredondant} sur
$X=\operatorname{Spec}(B)$ affine, et on écrira $Z$ au lieu de $Z_i$, et $x$
pour le point générique de $Z$ (et de $Z_\eta$). L'anneau noethérien $B$
contient donc $A$ comme sous-anneau, et $\mathcal F=\widetilde M$, où $M$ est un
$B$-module de type fini~; en outre, si $K$ est le corps des fractions de $A$,
le $B_{(K)}$-module $M_{(K)}$ est $\neq0$ et irredondant. Soit $q$ l'unique
élément de $\operatorname{Ass}(M_{(K)})$ et soit $\mathfrak p$ l'idéal premier
de $B$, image réciproque de $q$. En utilisant (5.11.1.1) comme dans la
démonstration de (9.7.6), on se ramène au cas où $B$ est \emph{intègre} et $M$
un sous-module $\neq0$ de $B$~; alors $\mathcal F_\eta$ est un
sous-$\mathcal O_{X_\eta}$-Module non nul de $\mathcal O_{Z_\eta}$, et en vertu
de (9.4.5), pour $s$ voisin de $\eta$, $\mathcal F_s$ est isomorphe à un
sous-$\mathcal O_{X_s}$-Module non nul de $\mathcal O_{Z_s}$~; comme $Z_s$ est
géométriquement intègre, les longueurs de $(\mathcal F_\eta)_x$ et de
$(\mathcal F_s)_{x_s}$ sont toutes deux égales à $1$, ce qui achève la
démonstration.

\oldpage[IV-3]{87}

\begin{corollary}[9.8.7]
\label{IV.9.8.7-fr}
Soient $f:X\to S$ un morphisme de présentation finie, $\mathcal F$ un
$\mathcal O_X$-Module quasi-cohérent de présentation finie. Pour tout $s\in S$,
soit $M(s)$ l'ensemble des couples $(d,m)$ tels qu'il existe une composante
irréductible de $\operatorname{Supp}(\mathcal F_s)$ de dimension $d$ et de
multiplicité radicielle $m$ pour $\mathcal F_s$ (4.7.8). Alors la fonction
$s\mapsto M(s)$ est localement constructible dans $S$.
\end{corollary}

Il résulte de (4.2.7) et (4.7.9) que la condition (9.2.1, (i)) est satisfaite
pour la propriété dont nous voulons démontrer qu'elle est constructible. On est
donc ramené au cas où $S$ est affine, noethérien et intègre et à prouver que $M$
est constante dans un voisinage du point générique de $S$~; mais alors la
proposition résulte de (9.8.3) et (9.8.6).

\begin{proposition}[9.8.8]
\label{IV.9.8.8-fr}
Soient $f:X\to S$ un morphisme de présentation finie, $\mathcal F$ un
$\mathcal O_X$-Module quasi-cohérent de présentation finie. Pour tout $s\in S$,
soit $l(s)$ la somme des multiplicités totales pour $\mathcal F_s$ (relatives à
$\boldsymbol{k}(s)$) des points génériques des composantes irréductibles de
$\operatorname{Supp}(\mathcal F_s)$ (4.7.12). Alors la fonction $s\mapsto l(s)$
est localement constructible dans $S$.
\end{proposition}

Compte tenu de (4.7.12), la condition (9.2.1, (i)) est satisfaite pour la
propriété dont nous voulons démontrer qu'elle est constructible, et on est donc
ramené au cas où $S$ est affine, noethérien et intègre de point générique
$\eta$, et à montrer que $l$ est constante dans un voisinage de $\eta$.
Utilisant les notations de (9.8.2), cela résulte de la définition (4.7.12), du
fait que le nombre géométrique des composantes irréductibles de chaque
$(Z_i)_s$ est constant dans un voisinage de $\eta$ (9.7.8), que la longueur
géométrique de $\mathcal F_s$ en $x_{is\alpha}$ est également à celle de
$\mathcal F_\eta$ en $x_i$ pour chaque $i$ tel que $(Z_i)_\eta$ soit maximal
(9.8.6) et enfin de ce que l'adhérence de $\{x_{is\alpha}\}$ est un cycle
premier associé maximal de $\mathcal F_s$ si et seulement si $(Z_i)_\eta$ est
un cycle premier associé maximal de $\mathcal F_\eta$ (9.8.3).

\begin{remark}[9.8.9]
\label{IV.9.8.9-fr}
On peut préciser de diverses façons les propositions précédentes~;
bornons-nous à un énoncé à titre d'exemple. Disons qu'une partie finie $P$ d'un
$k$-préschéma algébrique $X$ est \emph{saturée} si, pour tout couple de points
$x$, $y$ de $P$, les points génériques des composantes irréductibles de
$\overline{\{x\}}\cap\overline{\{y\}}$ appartiennent aussi à $P$~; pour toute
partie finie $Q$ de $X$, il existe une plus petite partie finie $P$ de $X$
contenant $Q$ et saturée~; nous dirons que $P$ est le \emph{saturé} de $Q$. Pour
tout $\mathcal O_X$-Module cohérent $\mathcal F$, nous appellerons
\emph{squelette primaire} de $\mathcal F$ le système $(P,Q,\omega,d,m)$ où
$Q=\operatorname{Ass}(\mathcal F)$, $P$ est le saturé de $Q$, $\omega$ la
relation d'ordre $\overline{\{x\}}\subset\overline{\{y\}}$ sur $P$, $d$ la
fonction $x\mapsto\dim\overline{\{x\}}$ sur $P$, $m$ la fonction
$x\mapsto\operatorname{long}_{\mathcal O_x}\mathcal F_x$ définie sur l'ensemble
des éléments de $Q$ maximaux pour la relation $\omega$. Appelons d'autre part
\emph{squelette virtuel} tout système $(P,Q,\omega,d,m)$, où $P$ est un
ensemble, $Q$ une partie de $P$, $\omega$ une relation d'ordre sur $P$, $d$ une
application de $P$ dans $\mathbf N$, $m$ une application dans $\mathbf N$ de
l'ensemble des éléments maximaux de $Q$~; on définit de façon évidente la notion
d'\emph{isomorphisme} de deux squelettes virtuels. Enfin, avec les notations
précédentes, appelons \emph{type primaire} de $\mathcal F$ la classe (pour la
relation d'isomorphie des squelettes virtuels) du squelette primaire de
$\mathcal F$. Il résulte de (4.2.6), (4.2.7), (4.2.8), (4.5.1) et (4.7.9) que
si, pour une extension algébriquement close $k'$ de $k$, on pose
$X'=X\otimes_k k'$ et $\mathcal F'=\mathcal F\otimes_k k'$, le
\emph{type primaire} de $\mathcal F'$ est indépendant de

\oldpage[IV-3]{88}

l'extension algébriquement close $k'$ de $k$ considérée~; nous dirons que c'est
le \emph{type primaire géométrique} de $\mathcal F$. Ces définitions posées,
l'énoncé que nous avons en vue est le suivant~:

\phantomsection\label{IV.9.8.9.1-fr}%
\textbf{(9.8.9.1)} Soient $f:X\to S$ un morphisme de présentation finie,
$\mathcal F$ un $\mathcal O_X$-Module quasi-cohérent de présentation finie.
Pour tout $s\in S$, soit $T(s)$ le type primaire géométrique de $\mathcal F_s$.
Alors la fonction $s\mapsto T(s)$ est localement constructible dans $S$.

Compte tenu des remarques qui précèdent, on est ramené comme d'ordinaire à
prouver que si $S$ est affine, noethérien et intègre de point générique $\eta$,
$T(s)$ est constante au voisinage de $\eta$. Raisonnant comme au début de
(9.7.7), on peut supposer que toutes les parties irréductibles de $X_\eta$ qui
interviennent sont géométriquement irréductibles, et alors la proposition
résulte de (9.5.1), (9.5.5), (9.8.3) et (9.8.6)~; nous laissons les détails au
lecteur. On pourrait généraliser en considérant plusieurs Modules cohérents et
en définissant leur «~squelette primaire simultané~», etc. La conclusion
générale de ce qui a été vu depuis le début de ce paragraphe est que pour toutes
les propriétés du type envisagé (et pour un $S$ irréductible) \emph{les
propriétés valables sur la «~fibre générique~» le sont encore sur toutes les
fibres voisines}.
\end{remark}

\subsection{Constructibilité des propriétés locales des fibres}
\label{subsection:IV.9.9-fr}

\begin{proposition}[9.9.1]
\label{IV.9.9.1-fr}
Soient $f:X\to S$ un morphisme localement de présentation finie, $Z$ une
partie localement constructible de $X$ telle que pour tout $s\in S$, $Z_s$ soit
fermé dans $X_s$, $\Phi$ une partie finie de $\mathbf Z\cup\{\pm\infty\}$.
Alors les parties suivantes de $X$ sont localement constructibles~:
\begin{enumerate}
\item[(i)] L'ensemble des $x\in X$ tels que
$\dim_x(Z_{f(x)})$ appartienne à $\Phi$.
\item[(ii)] L'ensemble des $x\in X$ tels que
$\operatorname{codim}_x(Z_{f(x)},X_{f(x)})$ appartienne à $\Phi$.
\item[(iii)] L'ensemble des $x\in X$ tels que l'anneau local
$\mathcal O_{X_{f(x)},x}$ soit équidimensionnel.
\end{enumerate}
\end{proposition}

On notera que les propriétés (i) et (ii) s'expriment encore en disant que les
fonctions $x\mapsto\dim_x(Z_{f(x)})$ et
$x\mapsto\operatorname{codim}_x(Z_{f(x)},X_{f(x)})$ sont
\emph{localement constructibles} dans $X$ ($\mathbf 0_{\mathrm{III}},9.3.1$).

Les questions étant locales sur $X$, on peut se borner au cas où
$S=\operatorname{Spec}(A)$ et $X=\operatorname{Spec}(B)$ sont affines et où
$f$ est un morphisme de présentation finie~; il existe alors un sous-anneau
$A_0$ de $A$ qui est une $\mathbf Z$-algèbre de type fini, un
$A_0$-préschéma de type fini $X_0$ et une partie constructible $Z_0$ de $X_0$
tels que $X=X_0\otimes_{A_0}A$ et $Z=h^{-1}(Z_0)$, où $h:X\to X_0$ est la
projection canonique ((8.9.1) et (8.3.11)). En outre, pour tout $s\in S$, si
$s_0$ est la projection de $s$ dans $S_0=\operatorname{Spec}(A_0)$, on a
\[
X_s=(X_0)_{s_0}\otimes_{k(s_0)}k(s),
\]
et $h_s$ est la projection $X_s\to(X_0)_{s_0}$, on a
$Z_s=h_s^{-1}((Z_0)_{s_0})$. Comme le morphisme $h_s$ est fidèlement plat et
quasi-compact, l'hypothèse que $Z_s$ est fermé dans $X_s$ entraîne que
$(Z_0)_{s_0}$ est fermé dans $(X_0)_{s_0}$ (2.3.12).

Cela étant, la transitivité des fibres (\textbf{I}, 3.6.4) et la prop. (4.2.7)
entraînent que l'ensemble des dimensions des composantes irréductibles de
$Z_s$ contenant $x$ est le même que l'ensemble des dimensions des composantes
irréductibles de $(Z_0)_{s_0}$ contenant $x_0=h_s(x)$. En particulier, on a
$\dim_x(Z_s)=\dim_{x_0}((Z_0)_{s_0})$. D'autre part, si $Z_s^{(\beta)}$ sont les
composantes irréductibles de $Z_s$ contenant $x$ et $X_s^{(\alpha)}$ les
composantes irréductibles de $X_s$

\oldpage[IV-3]{89}

contenant $x$, on a
\[
\operatorname{codim}_x(Z_s,X_s)=
\inf_\beta\left(\sup_\alpha
\bigl(\operatorname{codim}(Z_s^{(\beta)},X_s^{(\alpha)})\bigr)\right),
\]
$(\alpha,\beta)$ variant dans l'ensemble des couples tels que
$x\in Z_s^{(\beta)}\subset X_s^{(\alpha)}$ ($\mathbf 0$, 14.2.6). Comme les
préschémas algébriques irréductibles sont biéquidimensionnels (5.2.1), on peut
écrire, en vertu de ($\mathbf 0$, 14.3.3.1)
\begin{equation}
\tag{9.9.1.1}
\label{IV.9.9.1.1-fr}
\operatorname{codim}_x(Z_s,X_s)=
\inf_\beta\left(\sup_\alpha
\bigl(\dim(X_s^{(\alpha)})-\dim(Z_s^{(\beta)})\bigr)\right)
\end{equation}
avec le même choix des couples $(\alpha,\beta)$. Comme $h_s$ est fidèlement
plat et quasi-compact, pour tout couple formé d'une composante irréductible
$(X_0)_{s_0}^{(\lambda)}$ de $(X_0)_{s_0}$ contenant $x_0$ et d'une composante
irréductible $(Z_0)_{s_0}^{(\mu)}$ de $(Z_0)_{s_0}$ contenant $x_0$ et contenue
dans $(X_0)_{s_0}^{(\lambda)}$, il existe un couple
$(X_s^{(\alpha)},Z_s^{(\beta)})$ du type décrit plus haut et tel que
$Z_s^{(\beta)}$ domine $(Z_0)_{s_0}^{(\mu)}$ et que $X_s^{(\alpha)}$ domine
$(X_0)_{s_0}^{(\lambda)}$ (2.3.5). La formule (9.9.1.1) (et la formule analogue
appliquée à $(X_0)_{s_0}$) montrent alors, en vertu de (4.2.7), que l'on a
\[
\operatorname{codim}_x(Z_s,X_s)=
\operatorname{codim}_{x_0}((Z_0)_{s_0},(X_0)_{s_0}).
\]

On voit ainsi que si $E$ (resp. $E_0$) est l'ensemble des $x\in X$ (resp. des
$x_0\in X_0$) vérifiant l'une des conditions (i), (ii), (iii) de l'énoncé
(resp. la même condition), on a $E=h^{-1}(E_0)$, et en vertu de (1.8.2), on
voit qu'on peut se borner au cas où $A$ est \emph{noethérien}, donc aussi $B$.
Compte tenu de ($\mathbf 0_{\mathrm{III}},9.2.3$), ainsi que de (9.9.1.1), on
est ramené à voir que pour tout $x\in X$, il y a un voisinage $V$ de $x$ dans
$\overline{\{x\}}$ tel que, pour tout $x'\in V$, l'ensemble des dimensions des
composantes irréductibles de $X_{f(x')}$ (resp. $Z_{f(x')}$) contenant $x'$ est
\emph{le même}, et en outre qu'il en est de même de l'ensemble des couples
\[
\bigl(\dim(Z_{f(x')}^{(\beta)}),\dim(X_{f(x')}^{(\alpha)})\bigr)
\]
pour les couples formés d'une composante irréductible $X_{f(x')}^{(\alpha)}$
de $X_{f(x')}$ et d'une composante irréductible $Z_{f(x')}^{(\beta)}$ de
$Z_{f(x')}$ contenue dans $X_{f(x')}^{(\alpha)}$ et contenant $x'$. On peut
évidemment pour cela remplacer $S$ par le sous-préschéma réduit $S'$ de $S$
ayant $\overline{\{f(x)\}}$ pour espace sous-jacent, et $X$ par
$X'=f^{-1}(S')$, les fibres de $X$ et $X'$ aux points de $S'$ étant les mêmes.
Autrement dit, on peut se borner au cas où $S$ est \emph{intègre} et où
$\eta=f(x)$ est son point générique.

Par hypothèse $Z_\eta$ est fermé dans $X_\eta$~; comme $Z$ est constructible,
il résulte de (8.3.11), appliqué suivant la méthode de (8.1.2, a)), que l'on
peut, en remplaçant au besoin $S$ par un voisinage ouvert de $\eta$, supposer
que $Z$ est \emph{fermé} dans $X$. Soient $X_i$ (resp. $Z_j$) les composantes
irréductibles de $X$ (resp. $Z$) contenant $x$~; en vertu de
($\mathbf 0_{\mathrm I},2.1.8$), les $X_i\cap X_\eta$ (resp.
$Z_j\cap X_\eta$) sont les composantes irréductibles de $X_\eta$ (resp.
$Z_\eta$) contenant $x$~; en vertu de (9.5.1), on peut supposer en outre, en
restreignant au besoin $S$ à un voisinage de $\eta$, que les $X_i$ (resp.
$Z_j$) sont exactement les composantes irréductibles de $X$ (resp. $Z$)
rencontrant $\overline{\{x\}}$ et que les
$(X_i)_s\cap\overline{\{x\}}$ et $(Z_j)_s\cap\overline{\{x\}}$ sont non vides
pour tout $s\in S$. Cela étant, il résulte encore de (9.7.1) que l'on peut
supposer, en restreignant $S$, que les couples $(i,j)$ tels que
$(Z_j)_s\subset(X_i)_s$ sont les mêmes pour tout $s\in S$. La conclusion
résulte alors de (9.5.6)~: pour tout $s$ assez voisin de $\eta$, toutes les
composantes irréductibles de $(X_i)_s$ (resp. $(Z_j)_s$) ont une même dimension
égale à celle de $(X_i)_\eta$ (resp. $(Z_j)_\eta$). En outre, si $(i,j)$ est un
couple tel que $Z_j\not\subset X_i$, $(X_i)_\eta$ ne contient pas le

\oldpage[IV-3]{90}

point générique de $(Z_j)_\eta$, donc
$\dim((X_i\cap Z_j)_\eta)<\dim((Z_j)_\eta)$~; par suite, la dimension commune
des composantes irréductibles de $(X_i)_s\cap(Z_j)_s$ est, pour tout $s\in S$,
strictement inférieure à la dimension commune des composantes irréductibles de
$(Z_j)_s$, ce qui prouve qu'aucune des composantes irréductibles de $(Z_j)_s$
n'est contenue dans une composante irréductible de $(X_i)_s$ pour $s\in S$.
C.Q.F.D.

\begin{proposition}[9.9.2]
\label{IV.9.9.2-fr}
Soient $f:X\to S$ un morphisme localement de présentation finie, $\mathcal F$
un $\mathcal O_X$-Module quasi-cohérent de présentation finie, $\Phi$ une
partie finie de $\mathbf Z\cup\{\pm\infty\}$. Les parties suivantes de $X$
sont localement constructibles~:
\begin{enumerate}
\item[(i)] L'ensemble des points $x\in X$ tels que $\mathcal F_{f(x)}$ soit
équidimensionnel au point $x$ (5.1.12).
\item[(ii)] L'ensemble des $x\in X$ tels que les longueurs géométriques de
$\mathcal F_{f(x)}$ relatives à $k(f(x))$, aux points génériques des composantes
irréductibles de $\operatorname{Supp}(\mathcal F_{f(x)})$ contenant $x$
(4.7.5), appartiennent à $\Phi$.
\item[(iii)] L'ensemble des $x\in X$ tels que les dimensions des cycles
premiers associés à $\mathcal F_{f(x)}$ et contenant $x$ appartiennent à
$\Phi$.
\item[(iv)] L'ensemble des $x\in X$ tels que $\mathcal F_{f(x)}$ soit
géométriquement réduit au point $x$ (4.6.17).
\item[(v)] L'ensemble des $x\in X$ tels que $\mathcal F_{f(x)}$ soit
géométriquement intègre au point $x$ (4.6.22).
\item[(vi)] L'ensemble des $x\in X$ tels que
$\operatorname{dim.proj}((\mathcal F_{f(x)})_x)\in\Phi$.
\item[(vii)] L'ensemble des $x\in X$ tels que
$\operatorname{coprof}((\mathcal F_{f(x)})_x)\in\Phi$.
\item[(viii)] L'ensemble des $x\in X$ tels que $\mathcal F_{f(x)}$ possède la
propriété $(S_k)$ au point $x$ (5.7.2).
\item[(ix)] L'ensemble des $x\in X$ tels que $\mathcal F_{f(x)}$ soit un
$\mathcal O_{X_{f(x)}}$-Module de Cohen-Macaulay au point $x$ (5.7.1).
\end{enumerate}
\end{proposition}

(i) Le support $Z$ de $\mathcal F$ est localement constructible et fermé dans
$X$ (8.9.1), et en considérant un sous-préschéma de $X$ ayant $Z$ pour espace
sous-jacent et qui soit de présentation finie sur $S$ (8.9.1), on voit que
l'assertion (i) est un cas particulier de (9.9.1, (iii)).

(ii) Toutes les propriétés envisagées sont locales sur $X$, et on se bornera
donc au cas où $X=\operatorname{Spec}(B)$ et $S=\operatorname{Spec}(A)$ sont
affines et $f$ un morphisme de présentation finie. On garde les notations du
début de la démonstration de (9.9.1), et on suppose en outre $A_0$ choisi de
sorte qu'il existe un $\mathcal O_{X_0}$-Module cohérent $\mathcal F_0$ tel que
$\mathcal F$ soit isomorphe à $\mathcal F_0\otimes_{A_0}A$. Alors ((4.2.7) et
(4.7.9)) l'ensemble des longueurs géométriques de
$(\mathcal F_0)_{f(x_0)}$ aux points génériques des composantes irréductibles de
$\operatorname{Supp}((\mathcal F_0)_{f(x_0)})$ qui contiennent $x_0$ est le
même que l'ensemble analogue pour $x$ et $\mathcal F_{f(x)}$~; autrement dit,
si $E$ (resp. $E_0$) est l'ensemble des $x\in X$ (resp. des $x_0\in X_0$)
vérifiant la condition (ii) de l'énoncé, on a $E=h^{-1}(E_0)$, et en vertu de
(1.8.2), on voit qu'on peut se borner à considérer le cas où $A$ est
\emph{noethérien}. Comme dans la démonstration de (9.9.1), on voit qu'on est
ramené à montrer que, pour tout $x\in X$, il existe un voisinage $V$ de $x$
dans $\overline{\{x\}}$ tel que, pour tout $x'\in V$, l'ensemble des longueurs
géométriques de $\mathcal F_{f(x')}$ aux points génériques des composantes
irréductibles de son support contenant $x'$ soit \emph{le même}. En outre, si
$S'$ est le sous-préschéma réduit de $S$ ayant $\overline{\{f(x)\}}$ pour espace
sous-jacent, et si $X'=f^{-1}(S')$, les fibres de $X$ et de $X'$ aux points de
$S'$ sont les mêmes,

\oldpage[IV-3]{91}

donc on peut remplacer $S$ par $S'$ et $X$ par $X'$, autrement dit supposer que
$S$ est \emph{intègre} et que $\eta=f(x)$ est son point générique.

Or, si l'on pose $Z=\operatorname{Supp}(\mathcal F)$, le même raisonnement que
dans la démonstration de (9.9.1) montre que, si $Z_i$ sont les composantes
irréductibles de $Z$ contenant $x$, on peut supposer que ce sont exactement les
composantes irréductibles de $Z$ rencontrant $\overline{\{x\}}$ et que
$(Z_i)_s\cap\overline{\{x\}}\neq\varnothing$ pour tout $s\in S$. La conclusion
résulte alors de (9.8.3) et (9.8.6).

(iii) On se ramène comme dans (ii) au cas où $S$ est noethérien et intègre et
$\eta=f(x)$ son point générique, en utilisant (4.2.7). Comme dans la
démonstration de (9.9.1), on est ramené à montrer qu'il existe un voisinage $V$
de $x$ dans $\overline{\{x\}}$ tel que, pour tout $x'\in V$, l'ensemble des
dimensions des cycles premiers associés à $\mathcal F_{f(x')}$ qui contiennent
$x'$ soit \emph{le même}. Or, si $Z'_i$ sont les adhérences dans $X$ des cycles
premiers associés à $\mathcal F_\eta$ qui contiennent $x$ (cf. (9.8.2)), il
résulte de (9.8.3) et (9.5.1) que pour tout $s$ assez voisin de $\eta$, les
cycles premiers associés à $\mathcal F_s$ qui rencontrent
$\overline{\{x\}}$ sont des composantes irréductibles des $(Z'_i)_s$ et, en
vertu de (9.5.6), toutes ces composantes irréductibles ont même dimension égale
à $\dim((Z'_i)_\eta)$, d'où la conclusion.

(iv) On raisonne comme dans (iii), en utilisant cette fois (4.7.11). Avec les
mêmes notations que dans (iii), les cycles premiers associés à $\mathcal F_s$
qui sont des composantes irréductibles de $(Z'_i)_s$ sont immergés si et
seulement si $(Z'_i)_\eta$ est un cycle premier associé immergé de
$\mathcal F_\eta$. On conclut déjà que si $x$ appartient à un (resp.
n'appartient à aucun) cycle premier associé immergé de $\mathcal F_\eta$,
l'ensemble des $x'\in\overline{\{x\}}$ tels que $x'$ appartienne à un (resp.
n'appartienne à aucun) cycle premier associé immergé de $\mathcal F_{f(x')}$
est un voisinage de $x$ dans $\overline{\{x\}}$. La conclusion résulte de cette
remarque, de la caractérisation des points où un Module est géométriquement
réduit (4.7.10), et de (ii).

(v) Compte tenu de (4.7.11), on se ramène encore au cas où $S$ est noethérien
et intègre et où $\eta=f(x)$ est son point générique. Soit $Y$ un
sous-préschéma fermé de $X$ ayant $\operatorname{Supp}(\mathcal F)$ pour espace
sous-jacent. Raisonnant comme au début de la démonstration de (9.7.7), on voit
qu'il existe une $A$-algèbre intègre finie $A'$ (de sorte que si
$S'=\operatorname{Spec}(A')$, le morphisme $S'\to S$ est fini et surjectif)
telle que si $Y'=Y_{(S')}$ et si $\eta'$ est le point générique de $S'$, les
composantes irréductibles de $Y'_{\eta'}$ soient géométriquement irréductibles.
D'autre part, si $X'=X_{(S')}$, le morphisme projection $h:X'\to X$ est fini
et surjectif, donc fermé~; par suite, si $x'$ est un point de $X'$ au-dessus de
$x$, on a $h(\overline{\{x'\}})=\overline{\{x\}}$ et l'image par $h$ d'un
voisinage de $x'$ dans $\overline{\{x'\}}$ est un voisinage de $x$ dans
$\overline{\{x\}}$. Compte tenu de (4.7.11) et posant
$\mathcal F'=\mathcal F_{(S')}$, on est donc ramené à prouver que si
$\mathcal F'_{\eta'}$ est (resp. n'est pas) géométriquement intègre au point
$x'$, l'ensemble des $y'\in\overline{\{x'\}}$ tels que
$\mathcal F'_{f'(y')}$ (où $f'=f_{(S')}:X'\to S'$) soit (resp. ne soit pas)
géométriquement intègre au point $y'$ est un voisinage de $x'$ dans
$\overline{\{x'\}}$. On peut donc se borner au cas où $X'=X$ et où les
composantes irréductibles de $Y_\eta$ sont géométriquement irréductibles~; si
l'on désigne par $Z_i$ des parties fermées de $X$ telles que les
$Z_i\cap X_\eta$ soient les

\oldpage[IV-3]{92}

composantes irréductibles de $Y_\eta$, il résulte de (9.7.7), (9.7.8) et
(9.5.3) que pour tout $s$ voisin de $\eta$, les $(Z_i)_s$ sont les composantes
irréductibles de $Y_s$ et qu'elles sont \emph{géométriquement irréductibles}.
Dire qu'en un point $y\in X_s$, $\mathcal F_s$ est géométriquement intègre
signifie alors que $\mathcal F_s$ est géométriquement réduit en ce point et en
outre que $y$ n'appartient qu'à \emph{un seul} des $(Z_i)_s$ (4.6.22). La
conclusion résulte donc d'une part de (iv) et d'autre part de (9.5.1) appliqué à
l'intersection de $\overline{\{x\}}$ et de chaque $Z_i$.

(vi) Gardant les mêmes notations que dans (ii), il résulte de (6.2.1) que l'on
a
\[
\operatorname{dim.proj}((\mathcal F_s)_x)=
\operatorname{dim.proj}(((\mathcal F_0)_{s_0})_{x_0})~;
\]
on peut donc encore se borner au cas où $A$ est \emph{noethérien}. En outre, on
se ramène encore à montrer que, pour tout $x\in X$, il existe un voisinage $V$
de $x$ dans $\overline{\{x\}}$ tel que, pour tout $x'\in V$, on ait
\[
\operatorname{dim.proj}((\mathcal F_{f(x')})_{x'})=
\operatorname{dim.proj}((\mathcal F_{f(x)})_x)~;
\]
et comme ci-dessus, on peut supposer que $S$ est \emph{intègre} et que
$\eta=f(x)$ est son point générique, de sorte que l'on a
$(\mathcal F_\eta)_x=\mathcal F_x$. En vertu du théorème de platitude générique
(6.9.1), on peut, en remplaçant au besoin $S$ par un voisinage ouvert de
$\eta$, supposer que le morphisme $f$ est \emph{plat} et que $\mathcal F$ est
$f$-plat~; on a alors
$\operatorname{dim.proj}((\mathcal F_{f(x')})_{x'})=
\operatorname{dim.proj}(\mathcal F_{x'})$ pour tout $x'\in X$ en vertu de
(6.2.3). Cela étant, en vertu de (6.11.1), on peut (en remplaçant au besoin $X$
par un voisinage ouvert de $x$) supposer que
$\operatorname{dim.proj}(\mathcal F_{x'})\leq
\operatorname{dim.proj}(\mathcal F_x)$ pour tout $x'\in X$. D'autre part, si
$\operatorname{dim.proj}(\mathcal F_x)=n$, il y a par hypothèse un
$\mathcal O_x$-module de type fini $M$ tel que
$\operatorname{Ext}_{\mathcal O_x}^n(\mathcal F_x,M)\neq0$
($\mathbf 0$, 17.2.4). Or, il existe un $\mathcal O_X$-Module cohérent
$\mathcal G$ tel que $M=\mathcal G_x$ (en remplaçant au besoin $X$ par un
voisinage ouvert de $x$ ($\mathbf 0_{\mathrm I}$, 5.3.8))~; en vertu de
(T, 4.2.2), on a donc
$(\mathcal{E}xt_{\mathcal O_X}^n(\mathcal F,\mathcal G))_x\neq0$. Mais
$\mathcal{E}xt_{\mathcal O_X}^n(\mathcal F,\mathcal G)$ est un
$\mathcal O_X$-Module cohérent ($\mathbf 0_{\mathrm{III}}$, 12.3.3), donc son
support $Y$ est fermé ($\mathbf 0_{\mathrm I}$, 5.2.2)~; comme il contient $x$,
il contient aussi $\overline{\{x\}}$, d'où l'on conclut (en appliquant encore
(T, 4.2.2)) que l'on a
$\operatorname{dim.proj}(\mathcal F_{x'})\geq n$ pour tout
$x'\in\overline{\{x\}}$, ce qui achève de prouver l'assertion dans le cas (vi).

(vii) Comme $B$ est une $A$-algèbre de type fini, $X$ est $S$-isomorphe à un
\emph{sous-schéma fermé} d'un $S$-schéma de la forme
$Y=\operatorname{Spec}(A[T_1,\ldots,T_r])$~; si $j:X\to Y$ est l'injection
canonique, et $\mathcal G=j_*(\mathcal F)$, on a
$\mathcal G_s=(j_s)_*(\mathcal F_s)$ pour tout $s\in S$, et, en vertu de
($\mathbf 0$, 16.4.11), on peut se borner à démontrer l'assertion pour $Y$ et
$\mathcal G$. Autrement dit, on peut supposer que $B=A[T_1,\ldots,T_r]$, de
sorte que chacun des schémas $X_s=\operatorname{Spec}(k(s)[T_1,\ldots,T_r])$
est \emph{régulier} ($\mathbf 0$, 17.3.7). Soit alors
$W=\operatorname{Supp}(\mathcal F)$, de sorte que
$W_s=\operatorname{Supp}(\mathcal F_s)$ (\textbf{I}, 9.1.13)~; on a, par
(6.11.2.1)
\begin{equation}
\tag{9.9.2.1}
\label{IV.9.9.2.1-fr}
\operatorname{coprof}((\mathcal F_{f(x)})_x)=
\operatorname{dim.proj}((\mathcal F_{f(x)})_x)-
\operatorname{codim}_x(W_{f(x)},X_{f(x)}).
\end{equation}
Mais comme $W$ est constructible (8.9.1) et chaque $W_s$ fermé, il résulte de
(vi) et de (9.9.1, (ii)) que les deux fonctions du second membre de (9.9.2.1)
sont constructibles~; il en est donc de même de leur différence ce qui achève
de prouver la proposition dans le cas (vii).

(viii) Soit $U_n$ l'ensemble des $x\in X$ tels que
$\operatorname{coprof}((\mathcal F_{f(x)})_x)\leq n$, et posons
$Z_n=X-U_n$~; il résulte de (vii) que les $Z_n$ sont constructibles~; en outre,
comme

\oldpage[IV-3]{93}

la fonction $x\mapsto\dim_x(W_{f(x)})$ est constructible en vertu de
(9.9.1, (i)), elle ne prend qu'un nombre fini de valeurs, donc les nombres
$\dim(W_{f(x)})$ ont une borne supérieure finie $m$ lorsque $x$ parcourt $X$~;
comme
$\operatorname{coprof}((\mathcal F_{f(x)})_x)\leq
\dim((\mathcal F_{f(x)})_x)\leq\dim(W_{f(x)})$, on voit que l'on a
$Z_n=\varnothing$ pour $n\geq m$. Enfin, il résulte de (6.11.2, (i)) que pour
tout $n$ et tout $s\in S$, $(Z_n)_s$ est \emph{fermé} dans $X_s$. D'après
(5.7.4), l'ensemble des $x\in X$ où $(\mathcal F_{f(x)})_x$ possède la
propriété $(S_k)$ est l'ensemble des $x\in X$ vérifiant toutes les relations
\begin{equation}
\tag{9.9.2.2}
\label{IV.9.9.2.2-fr}
\operatorname{codim}_x((Z_n)_{f(x)},W_{f(x)})>n+k
\end{equation}
pour tout $n\geq0$~; comme cette relation est automatiquement vérifiée pour
$n\geq m$, on n'a à considérer les relations (9.9.2.2) que pour $0\leq n<m$.
Mais en vertu de (9.9.1, (ii)), l'ensemble $V_{n,k}$ des $x$ vérifiant (9.9.2.2)
est constructible, et il en est de même de l'intersection de ces ensembles pour
$0\leq n<m$.

(ix) L'assertion résulte ici aussitôt de (vii), l'ensemble des $x\in X$ où
$(\mathcal F_{f(x)})_x$ est un module de Cohen-Macaulay étant défini par la
relation $\operatorname{coprof}((\mathcal F_{f(x)})_x)=0$.

\begin{corollary}[9.9.3]
\label{IV.9.9.3-fr}
Soient $f:X\to S$ un morphisme de présentation finie, $\mathcal F$ un
$\mathcal O_X$-Module de présentation finie, $\mathbf P$ l'une des propriétés
(i) à (ix) de (9.9.2). Alors l'ensemble des $s\in S$ tels que la propriété
$\mathbf P$ soit vraie en tous les points $x\in X_s$, est localement
constructible dans $S$.
\end{corollary}

En effet, son complémentaire est l'image par $f$ du complémentaire de
l'ensemble $E$ des points où $\mathbf P$ est vraie. Comme $E$ est localement
constructible dans $X$, il en est de même de $X-E$, donc $f(X-E)$ est
localement constructible dans $S$ en vertu du théorème de Chevalley (1.8.4).

\begin{proposition}[9.9.4]
\label{IV.9.9.4-fr}
Soit $f:X\to S$ un morphisme localement de présentation finie. L'ensemble des
$x\in X$ tels que $X_{f(x)}$ ait au point $x$ l'une (déterminée) des propriétés
suivantes~:
\begin{enumerate}
\item[(i)] être géométriquement régulier~;
\item[(ii)] posséder la propriété $(R_k)$ géométrique~;
\item[(iii)] être géométriquement normal~;
\item[(iv)] être géométriquement réduit (i.e. séparable)~;
\item[(v)] être géométriquement ponctuellement intègre~;
\end{enumerate}
est une partie localement constructible de $X$.
\end{proposition}

Pour les propriétés (iv) et (v), cela résulte de (9.9.2, (iv) et (v)) appliqué à
$\mathcal F=\mathcal O_X$. Pour les autres propriétés, compte tenu de (6.7.8),
on se ramène, comme au début de (9.9.2) au cas où $S$ est noethérien et intègre
et de point générique $\eta=f(x)$. En outre, en vertu du théorème de platitude
générique (6.9.1), on peut, en remplaçant $S$ par un voisinage ouvert de
$\eta$, supposer que $f$ est un morphisme \emph{plat}. Dire que $X_{f(y)}$ est
géométriquement régulier au point $y$ signifie alors que $f$ est
\emph{régulier} au point $y$ (6.8.1), et on sait que l'ensemble $L$ de ces
points est \emph{ouvert} dans $X$ (6.8.7), ce qui prouve la proposition dans le
cas (i).

Pour prouver le cas (ii), posons $F=X-L$, qui est \emph{fermé} dans $X$. Dire
que $X_s$

\oldpage[IV-3]{94}

vérifie la propriété géométrique $(R_k)$ au point $y\in f^{-1}(s)$ signifie, ou
bien que l'on a $y\in L$, ou bien que les points génériques $z_i$ des
composantes irréductibles de l'ensemble fermé $F_s$ qui contiennent $y$
vérifient la relation $\dim(\mathcal O_{X_s,z_i})\geq k+1$ (compte tenu de
(4.2.7) et (5.2.3))~; autrement dit, les points $y\in F_s$ où $X_s$ vérifie la
propriété $(R_k)$ géométrique sont ceux tels que
$\operatorname{codim}_y(F_s,X_s)\geq k+1$ (5.1.2). La conclusion résulte donc
de (i) et de (9.9.1, (ii)).

Enfin, dans le cas (iii), la conclusion résulte de (ii), de (9.9.2, (viii)), du
fait que la propriété $(S_k)$ est stable par extension du corps de base (6.7.1)
et enfin du critère de Serre (5.8.6).

\begin{corollary}[9.9.5]
\label{IV.9.9.5-fr}
Soit $f:X\to S$ un morphisme de présentation finie, et notons $\mathbf P$
l'une des propriétés (i) à (v) de (9.9.4). Alors l'ensemble des $s\in S$ tels
que la propriété $\mathbf P$ soit vraie en tous les points $x\in X_s$, est
localement constructible dans $S$.
\end{corollary}

La démonstration à partir de (9.9.4) est la même que celle de (9.9.3) à partir
de (9.9.2).

\begin{proposition}[9.9.6]
\label{IV.9.9.6-fr}
Soient $f:X\to S$ un morphisme localement de présentation finie,
$\mathcal L_\bullet$ un complexe formé de $\mathcal O_X$-Modules
quasi-cohérents de présentation finie~; pour tout $s\in S$, soit
$(\mathcal L_\bullet)_s$ le complexe
$\mathcal L_\bullet\otimes_{\mathcal O_S}k(s)$ de
$\mathcal O_{X_s}$-Modules cohérents de type fini. Alors, pour un entier $n$
donné, l'ensemble des $x\in X$ tels que
$(\mathcal H_n((\mathcal L_\bullet)_{f(x)}))_x=0$ est localement constructible
dans $X$.
\end{proposition}

On peut se borner au cas où $\mathcal L_i=0$ sauf pour $i=0$, 1 ou 2, et où
$n=1$. En outre, la question étant locale sur $X$, on peut se borner au cas où
$S=\operatorname{Spec}(A)$ et $X=\operatorname{Spec}(B)$ sont affines, $B$
étant une $A$-algèbre de présentation finie. Il existe alors un sous-anneau
noethérien $A_0$ de $A$, un $A_0$-préschéma de type fini $X_0$ et un complexe
$\mathcal L_\bullet^{(0)}$ de $\mathcal O_{X_0}$-Modules cohérents, nuls sauf en
dimensions 0, 1 et 2 et tels que
$X=X_0\otimes_{A_0}A$ et
$\mathcal L_\bullet=\mathcal L_\bullet^{(0)}\otimes_{A_0}A$. Pour tout
$s\in S$, si $s_0$ est la projection de $s$ dans
$S_0=\operatorname{Spec}(A_0)$, on a
$X_s=(X_0)_{s_0}\otimes_{k(s_0)}k(s)$, et le morphisme projection
$X_s\to(X_0)_{s_0}$ est fidèlement plat~; on en conclut que l'on a
\[
\mathcal H_n((\mathcal L_\bullet)_s)=
\mathcal H_n((\mathcal L_\bullet^{(0)})_{s_0})
\otimes_{k(s_0)}k(s),
\]
et par suite, si $E$ (resp. $E_0$) est l'ensemble des $x\in X$ (resp.
$x_0\in X_0$) tels que
\[
(\mathcal H_n((\mathcal L_\bullet)_{f(x)}))_x=0
\quad\text{(resp.}\quad
(\mathcal H_n((\mathcal L_\bullet^{(0)})_{f_0(x_0)}))_{x_0}=0\text{)},
\]
on a $E=h^{-1}(E_0)$, où $h:X\to X_0$ est la projection canonique. En vertu
de (1.8.2), on peut donc se borner au cas où $A$ est \emph{noethérien}~; il
s'agit de voir ($\mathbf 0_{\mathrm{III}},9.2.3$) que si $x\in X$ est tel que
$(\mathcal H_n((\mathcal L_\bullet)_{f(x)}))_x=0$ (resp. $\neq0$), il existe un
voisinage ouvert $V$ de $x$ dans $\overline{\{x\}}$ tel que, pour tout
$x'\in V$, on ait
$(\mathcal H_n((\mathcal L_\bullet)_{f(x')}))_{x'}=0$ (resp. $\neq0$).
Remplaçant $S$ par le sous-préschéma réduit de $S$ ayant
$\overline{\{f(x)\}}$ pour espace sous-jacent, on peut supposer que $S$ est
intègre et que $f(x)=\eta$ est son point générique. Alors, en restreignant $S$
à un voisinage ouvert de $\eta$, on peut supposer que pour tout $s\in S$, on a
$(\mathcal H_n(\mathcal L_\bullet))_s=
\mathcal H_n((\mathcal L_\bullet)_s)$ (9.4.3), et par suite, si $Z$ est le
support de $\mathcal H_n(\mathcal L_\bullet)$, le support de
$\mathcal H_n((\mathcal L_\bullet)_s)$ est $Z_s=Z\cap X_s$
(\textbf{I}, 9.1.13.1). L'hypothèse est que
$Z_\eta\cap\overline{\{x\}}=\varnothing$ (resp. $\neq\varnothing$). Comme
$Z\cap\overline{\{x\}}$ est fermé dans l'espace noethérien $X$, on conclut de
(9.5.1) qu'il y a un voisinage $U$ de $\eta$ dans $S$ tel que, pour tout
$s\in U$, on ait $Z_s\cap\overline{\{x\}}=\varnothing$ (resp.
$\neq\varnothing$). L'ensemble $V=f^{-1}(U)\cap\overline{\{x\}}$ répond donc à
la question.
